\documentclass[11pt]{amsart}
\usepackage{amsmath,amssymb,amsthm,mathtools}
\usepackage[margin=1.15in]{geometry}
\setlength{\parskip}{4pt plus 1pt}
\usepackage{hyperref}
\usepackage{array}
\usepackage{booktabs}
\usepackage{longtable}
\usepackage{xcolor}
\usepackage{enumitem}
% The workspace TeX installation has bitmap AMS fonts; character expansion
% requires scalable outlines and aborts pdfTeX. Protrusion remains enabled.
\usepackage[final,expansion=false]{microtype}
\emergencystretch=2.5em
\tolerance=1200

%% Theorem environments
\newtheorem{theorem}{Theorem}[section]
\newtheorem{proposition}[theorem]{Proposition}
\newtheorem{lemma}[theorem]{Lemma}
\newtheorem{corollary}[theorem]{Corollary}
\newtheorem{conjecture}[theorem]{Conjecture}
\theoremstyle{definition}
\newtheorem{definition}[theorem]{Definition}
\newtheorem{example}[theorem]{Example}
\theoremstyle{remark}
\newtheorem{remark}[theorem]{Remark}
\newtheorem{notation}[theorem]{Notation}
\newtheorem{convention}[theorem]{Convention}
\newtheorem{problem}[theorem]{Problem}

%% Negative-result environments
\theoremstyle{definition}
\newtheorem{nogo}{No-go}[section]
\newtheorem{obstruction}{Wall}[section]
\newtheorem{witness}{Witness}[section]

%% Math macros
\newcommand{\RH}{\mathrm{RH}}
\renewcommand{\Re}{\operatorname{Re}}
\renewcommand{\Im}{\operatorname{Im}}
\newcommand{\half}{\tfrac{1}{2}}
\newcommand{\gam}{\gamma}
\newcommand{\Lam}{\Lambda}
\newcommand{\eps}{\varepsilon}
\newcommand{\supp}{\operatorname{supp}}
\newcommand{\artanh}{\operatorname{artanh}}
\newcommand{\cW}{\mathcal{W}}
\newcommand{\cS}{\mathcal{S}}
\newcommand{\cD}{\mathcal{D}}
\newcommand{\cH}{\mathcal{H}}
\newcommand{\cO}{\mathcal{O}}
\newcommand{\cL}{\mathcal{L}}
\newcommand{\R}{\mathbb{R}}
\newcommand{\C}{\mathbb{C}}
\newcommand{\Z}{\mathbb{Z}}
\newcommand{\Q}{\mathbb{Q}}
\newcommand{\N}{\mathbb{N}}
\newcommand{\Dd}{\mathbb{D}}
\newcommand{\Tt}{\mathbb{T}}
\newcommand{\Pp}{\mathsf{P}}
\newcommand{\Bb}{\mathsf{B}}
\newcommand{\Cc}{\mathsf{C}}
\newcommand{\Uu}{\mathsf{U}}
\newcommand{\Aa}{\mathsf{A}}
\DeclareMathOperator{\Tr}{Tr}
\DeclareMathOperator{\Spec}{Spec}
\DeclareMathOperator{\Res}{Res}
\DeclareMathOperator{\diag}{diag}
\DeclareMathOperator*{\esssup}{ess\,sup}

\title[Limits of the positivity family]{The Limits of the Positivity Family in the
Li--Keiper Route to the Riemann Hypothesis\\[0.35em]
\normalsize\textmd{Sixteen closed transfer mechanisms, adaptive cutoffs, and
density-relaxed criteria}}

\author{David Alejandro Trejo Pizzo}
\thanks{Independent Researcher. \texttt{dtrejopizzo@gmail.com}}

\begin{document}
\maketitle

\begin{abstract}
The direct route from the Li--Keiper criterion to the Riemann Hypothesis developed here
terminates in a compact--tail inequality $(\star)$, bounding the prime counting error weighted by
a Laguerre kernel.  The cutoff audit in this paper corrects an important overstatement of its
required margin.  The constants $2002/501$ and $A_n/1001$ describe only a convenient sufficient
transfer obtained by replacing the actual cutoff by the raw validity floor $T\ge1000$.  Since
admissibility is upward closed, the intrinsic cofinal statement is instead
$4\lambda_n-A_n>0$: for each strict positive margin one may move the cutoff far enough out to
prove $(\star)$.  Moreover $4\lambda_n-A_n\ge0$ for all large $n$ already proves the Riemann
Hypothesis directly by Li's criterion.  Thus no \emph{quantitative} claim that
``semidefiniteness falls short by a divergent amount'' survives the cutoff correction.

What survives, and is the thesis of the title, is sharper and independent of the cutoff
arithmetic. Every instrument in the catalogue certifies that some \emph{auxiliary} object is
nonnegative: a measure, a square, a Gram form, a totally positive kernel, a completely monotone
function. The object whose nonnegativity is actually at stake is the completed cocycle
$H_u^{r}K_u^{r-1}$, and we prove that it is not of that kind. Its Euler and Gamma factors are
positive for every real exponent by infinite divisibility; its polar factor is signed for every
exponent, with an exactly computable negative density at $x_0=\tfrac12\log2$ equal to
$-r+(r-1)/\sqrt2$, which is bounded away from zero and worsens linearly in $r$. The target is
therefore a signed cancellation between two divergent blocks of the same order, $A_n$ and
$B_n=A_n-\lambda_n$, with a margin proportional to either. That is a coercivity estimate, not a
semidefiniteness one.

This paper is the record of a systematic failure. Sixteen documented attacks on that inequality
were closed, each with an exact witness rather than a report of difficulty: local
positive-semidefiniteness of an order kernel; the global Gram/energy factorization of the same
kernel; the non-local zero--lobe Gram; M\"obius inversion as a positive adjoint; the positive
Euler--Gamma Koszul complex; the nonnegative Jordan hierarchy $\mu*\log^{k}\ge0$; Selberg's
$\mu*\log^{2}\ge0$ with its $O(x)$ summation and the associated Riccati induction; radial Abel
and Ces\`aro--Fej\'er positivity of the incomplete Li block; real order and monotonicity of a
correlated-shift cocycle; the three-channel Jordan--Beta expansion with a positive polar square;
tower-by-tower algebraic squares; global compound-Poisson structure together with
$\mathrm{PF}_2/\mathrm{TP}_2$, log-concavity, decreasing variation and stochastic order;
complete monotonicity in the sense of Bernstein; Stein--Mecke generator positivity;
Fej\'er--Carath\'eodory's $\Re\Phi\le1$; the fourth power as a sum of squares; and a favourable
tail sign extracted from an absolute Vinogradov--Korobov envelope.

The sixteen failures are nevertheless strongly related.  Every attack tries to transfer the
nonnegativity of an auxiliary kernel, measure, moment sequence, or square to the compact
certificate.  At a \emph{prefixed} finite cutoff, and using only the two-sided tail envelope, the
bare conclusion $D_n^{[4]}=4\lambda_n-A_n\ge0$ leaves an uncontrolled signed tail; at the
coarse floor $T=1000$ the stronger sufficient condition is exactly
$D_n^{[4]}\ge A_n/1001$.  The catalogue proves that the proposed transfers do not supply that
extra information.  It does not prove that positivity of $D_n^{[4]}$ is globally inadequate,
nor that proportional coercivity is necessary for RH.  Three universal falsifiers implement
the transfer failures: an exact negative Laplace density at $x_0=\tfrac12\log2$, strictly below the first
prime atom, where the Euler channel has no support; off-line symmetric quartets, one of them
placed uniformly inside the Laguerre bulk and one giving
$16\bigl(1-\cosh(na)\cos(n\vartheta)\bigr)<0$ already at the first target index $n=150$; and
sign changes of the genuinely arithmetic blocks, certified in rational arithmetic
($H(2969)<-21<0<110<H(3167)$; a $\mathrm{PF}_2$ minor equal to $-1/2916$ at the primes $2$ and
$3$).

We isolate the $\tfrac12\log2$ phenomenon as a structural wall for the positive-measure
transfer in its own right --- the
arithmetic-propagation wall localized in the Laplace variable --- and prove that it is not a
near miss: the defect is $-r+(r-1)/\sqrt2<0$ for every $r>1$, failing by the constant factor
$r(1-1/\sqrt2)+1/\sqrt2$. We then describe two genuinely weaker successor criteria.  The
dominant-mode argument behind Bombieri--Lagarias implies that, if RH is false, exponentially
negative Li coefficients occur on a syndetic set.  Consequently it is enough to prove any
subexponential lower bound outside an exceptional set of logarithmic density zero, or on
consecutive blocks of unbounded length.  A smooth criterion that retains exactly this weakening is
the bounded Fermi average
$H_X^{-1}\sum_{n\le X}[n(1+e^{t(\lambda_n+b_n)})]^{-1}\to0$, where
$b_n\to\infty$ and $\log(1+b_n)=o(n)$.  We give its exact regularized prime--pole form, while
explicitly leaving its estimate open.  The exponential partition
$\sum_{n\le X}e^{-t\lambda_n}/n=O_t(\log X)$ is retained only as a stronger algebraic auxiliary:
it already forces a lower bound at every degree.  Finally, we prove
that operator-norm coercivity for the whitened arithmetic jet is \emph{false} rather than hard,
and distinguish the extremal (gapless, marginal) problem from the fixed-family problem that
$(\star)$ poses.  Fixed-vector non-alignment remains one possible stronger route, not a logically
necessary form of the missing theorem.

Nothing in this paper proves $(\star)$, the Riemann Hypothesis, or its negation. Every
no-go is stated relative to an explicitly declared list of admissible inputs, and none of them
refutes the existence of a signed non-local inequality special to the real von Mangoldt weights
$\Lambda(m)$.  The new density and partition criteria are equivalences and not estimates of the
real arithmetic side.
\end{abstract}

\tableofcontents
\newpage

%% ======================================================================
\section{Prologue: what this paper is, and why a negative result is worth writing}
\label{sec:prologue}
%% ======================================================================

There is a particular kind of experience in long research programmes which is almost never
written down, because the conventions of publication reward the announcement of a theorem and
are silent about the announcement of a category error. One works for a long time on a target
inequality. One tries a mechanism; it fails, and one records the failure. One tries a second,
structurally different mechanism; it fails too, for what appears to be an unrelated reason. One
tries a third. After the tenth or twelfth such episode the suspicion becomes unavoidable that
the failures are not independent events --- that they are one failure, seen from twelve angles,
and that the reason they keep recurring is not that the mechanisms are weak but that, for the
specific transfer they were asked to perform, they all use the same missing input.

This paper is the written form of that suspicion, made precise and proved. The investigation it
reports on had reduced the Riemann Hypothesis, along a route we shall recall in
\S\ref{sec:target}, to a single unconditional inequality about the Li--Keiper coefficients
$\lambda_n$ of the Riemann zeta function. The reduction is exact, the inequality is explicit,
and its truth would close the Hypothesis. Over the course of the work reported here, sixteen
distinct attacks on that inequality were opened and closed. Each was closed with what we shall call a
\emph{witness}: a rational number, an exact identity, a specific pair of primes, a specific
index $n$ --- something one can check, rather than a report that a route ``did not seem to
work''. The discipline of witnesses is what makes the present synthesis possible: because every
death is a computation, one can ask what the computations have in common.

They have this in common. Every one of the sixteen mechanisms is an instrument for proving that
some auxiliary object is nonnegative. Their outputs are: a Gram matrix is positive semidefinite; a
sequence is Polya frequency of order two; a measure is positive; a function is completely
monotone; a form is a sum of squares; a random variable dominates another stochastically; the
real part of an analytic function is bounded by one. These are all excellent instruments, and
the identities they produced along the way are, in several cases, genuinely new and worth
keeping. But their conclusion is always of one shape, namely
\[
   (\text{something})\ \ge\ 0 .
\]
At the time these mechanisms were built, the compact inequality was transferred through the
uniform floor $T\ge1000$ to the stronger target
\[
   \lambda_n\ \ge\ c\,A_n,\qquad c>\tfrac14\ \text{fixed},
\]
a bound in which the right-hand side is of size $\asymp n\log n$.  At that fixed floor the
algebraic gap between $D_n^{[4]}\ge0$ and the declared sufficient transfer is exactly
$A_n/1001$.  This is a valid diagnosis of those transfer attempts, but it is not the intrinsic
logical frontier: the actual cutoff may be enlarged separately for each $n$.  The corrected
frontier for cofinal cutoffs is the strict inequality $D_n^{[4]}>0$, while the non-strict
inequality $D_n^{[4]}\ge0$ for all large $n$ already implies RH directly.  Accordingly this
paper no longer promotes coercivity as a necessary category for a proof of RH.

The sharpest and most useful form of the lesson is chronological. The \emph{declared fixed-cutoff
transfer} changed type before we distinguished it from the master criterion. The reductions that
improved that sufficient margin --- from
$2\lambda_n\ge A_n$ down through $3\lambda_n\ge A_n$ to
$\tfrac{2002}{501}\lambda_n\ge A_n$ --- were real progress, and were correctly recorded as
such for a cutoff controlled only by its floor. But in pushing that sufficient exponent
\emph{below} the integer four they moved this particular transfer across the boundary of what a
bare square, positive measure, or semidefinite form certifies.
The exponent four is precisely where factorization lives: $H_u^4K_u^3=(H_u^2K_u^{3/2})^2$, and
$\mathrm{PF}_2$, sums of squares and Gram representations are all built for even exponents.
At the fixed floor, the inequality one gets from those structures, namely
$4\lambda_n-A_n\ge0$, does not by itself imply the compact certificate through the unsigned tail
estimate; the convenient stronger condition is
$4\lambda_n-A_n\ge A_n/1001$. Six of the sixteen mechanisms were opened without making this
distinction. Their individual witnesses remain valid, but the former conclusion that
semidefiniteness itself could never close RH is withdrawn.

We think this is worth writing down for three reasons.

First, because it is the truth about how this piece of the work went, and an account that reports
only its live fronts is not reporting on itself. We therefore collect in \S\ref{sec:scope} the
corrections and retractions this paper carries, several of them to our own earlier claims.

Second, because the negative result still shrinks the search space for the \emph{specific
transfers} catalogued here.  Their witnesses rule out positive-measure, local-square and
operator-norm shortcuts.  They do not force the successor to be coercive: after the quantifier
audit in \S\ref{sec:successor}, a nonlinear density estimate in the degree is another, strictly
weaker possibility.  We also prove there that operator-norm coercivity for the whitened
arithmetic jet is false, while fixed-vector non-alignment remains an optional stronger route.

Third, because several of the identities established on the way to these deaths are exact,
reusable and, as far as our bibliographic audit could determine, not in the literature in the
combination in which they appear here: the global energy identity for the order kernel
$w(\max(x,y))$ against $d\psi-dx$ (Theorem~\ref{thm:max-energy}); the exact Hessian cancellation
that collapses that mechanism (Theorem~\ref{thm:hessian-zero}); the local tower identities
$b_u(p^k)=\nabla^2\bigl((k+1)Q^k\bigr)$ and
$(J_u^{*r})(p^k)=\nabla^r\bigl[\binom{k+r-1}{r-1}Q^k\bigr]$
(Theorem~\ref{thm:tower-diff}); the exact hypergeometric density
$-rc\,e^{-q}\,{}_1F_1(1-r;2;cq)$ of the real polar channel (Theorem~\ref{thm:polar-1F1}); the
signed L\'evy triple $\nu_A^{[3]}+2\nu_K-3\nu_P$ with its exact negative value at
$\tfrac12\log2$ (Theorem~\ref{thm:stein-witness}); and the Fej\'er coefficient identity
$q_m=\tfrac{m+1}{u}\bigl(1-\sigma_m(\Phi;1)\bigr)$ together with the proof that the holomorphy
it requires is itself equivalent to the Hypothesis (Theorem~\ref{thm:fejer-rh}). A no-go proof
that produces exact identities is not a wasted computation; it is a computation whose conclusion
happens to be negative.

\subsection*{What we claim and what we do not}

We claim: the reductions of \S\ref{sec:target}, including upward closure of admissible cutoffs,
the cofinal exponent-four criterion (Theorem~\ref{thm:four}), the fixed-ratio sufficient form
(Theorem~\ref{thm:scalefree}) and the algebraic fixed-floor identity
(Theorem~\ref{thm:threshold}); the sixteen no-gos of
\S\S\ref{sec:mech-gram}--\ref{sec:mech-envelope}, each within the list of admissible inputs
declared at its head; the classification theorem of \S\ref{sec:synthesis}; the structural wall
of \S\ref{sec:halflog2}; and the negative result about operator-norm coercivity in
\S\ref{sec:successor}; and the density, block and nonlinear partition criteria proved there.

We do not claim: $(\star)$; the Riemann Hypothesis; the negation of any of the sixteen target
inequalities. A no-go in this paper says that a \emph{specified} mechanism cannot produce the
target from a \emph{specified} list of inputs. In particular, not one of them excludes the
existence of a signed, non-local inequality valid for the actual von Mangoldt sequence
$\Lambda(m)$ --- and every one of the sixteen sections ends by saying so explicitly, because
that inequality is exactly the thing that is missing.

\begin{convention}[the standard of proof in this paper]\label{conv:standard}
Every sign decision quoted as a certificate is exact. The rational witnesses of
\S\ref{sec:falsifiers} are computed in exact rational arithmetic with outward-rounded intervals
for the logarithms of primes; no floating-point comparison enters any statement labelled a
theorem, proposition or witness. Two numerical diagnostics appear in the text
(Remarks~\ref{rmk:diag-cubic} and \ref{rmk:diag-global}); both are explicitly labelled as
directional evidence, and neither is used in the proof of anything.
\end{convention}

\begin{convention}[unconditionality]\label{conv:cond}
Every statement in this paper is unconditional unless it is explicitly said to hold under the
Riemann Hypothesis, and every imported statement is marked accordingly at the point of import.
We are pedantic about this because one of the errors we had to correct was exactly the importation
of a conditional asymptotic as though it were unconditional; see Remark~\ref{rmk:erratum}.
\end{convention}

%% ======================================================================
\section{The target, its exact reductions, and the margin ladder}\label{sec:target}
%% ======================================================================

This section fixes notation and recalls, with proofs where they are short, the chain that
produces the inequality whose attempted proof is the subject of the paper. Nothing in this
section is negative; everything in it is an exact identity or an unconditional estimate. The
reader who knows the chain may skip to \S\ref{sec:ladder}, and the reader who wants only the
punchline may skip to Theorem~\ref{thm:threshold}.

\subsection{The Li--Keiper coefficients and their archimedean block}\label{sec:li}

Let
\[
   \xi(s)=\half s(s-1)\pi^{-s/2}\Gamma(s/2)\zeta(s)
\]
be the completed zeta function, entire of order one, satisfying $\xi(s)=\xi(1-s)$. Li's
coefficients \cite{Li1997} are
\[
   \lambda_n=\frac{1}{(n-1)!}\frac{d^n}{ds^n}
   \Bigl[s^{n-1}\log\xi(s)\Bigr]_{s=1},\qquad n\ge1,
\]
and admit the Hadamard expansion
\begin{equation}\label{eq:li-zeros}
   \lambda_n=\sum_\rho\Bigl[1-\Bigl(1-\frac1\rho\Bigr)^{\!n}\Bigr],
\end{equation}
the sum being over the non-trivial zeros, taken in conjugate pairs.

\begin{theorem}[Li's criterion]\label{thm:li}
$\lambda_n\ge0$ for all $n\ge1$ if and only if the Riemann Hypothesis holds.
\end{theorem}

This is \cite{Li1997}; see also the complements of Bombieri and Lagarias \cite{BombieriLagarias1999},
whose Theorem~1 gives the general form of the criterion for a multiset of complex numbers closed
under $z\mapsto1-\bar z$, and whose part (c) is the exponential-blowup statement we shall use in
\S\ref{sec:falsifiers}. The Cayley substitution
\begin{equation}\label{eq:cayley}
   w_\rho=1-\frac1\rho,\qquad\text{equivalently}\qquad \rho=\frac{1}{1-w_\rho},
\end{equation}
maps the critical line $\Re s=\half$ to the unit circle $|w|=1$, the half-plane
$\Re s>\half$ to the open unit disc, and turns \eqref{eq:li-zeros} into
$\lambda_n=\sum_\rho(1-w_\rho^n)$. Under the Hypothesis each conjugate pair contributes
$2-2\cos(n\vartheta_\rho)=|1-w_\rho^n|^2\ge0$; off the line a quartet contributes an expression
containing $\cosh(na)$ with $a\ne0$, which is the mechanism of every falsifier in this paper.

Separating the explicit-formula contributions of the completed factor from those of the Euler
product gives the decomposition on which the whole route rests:
\begin{equation}\label{eq:split}
   \lambda_n=A_n+\lambda_n^{\mathrm{prime}} ,
\end{equation}
where the \emph{archimedean block} is the exactly summable expression
\begin{equation}\label{eq:An}
   A_n=1-\frac n2(\gam+\log4\pi)
   +\sum_{r\ \mathrm{odd}}\Bigl[\Bigl(1-\frac1r\Bigr)^{\!n}-1+\frac nr\Bigr],
   \qquad A_n\sim\tfrac12 n\log n,
\end{equation}
and the prime block is
\begin{equation}\label{eq:lamprime}
   \lambda_n^{\mathrm{prime}}
   =-n+\int_1^\infty\bigl(\psi(y)-y\bigr)f_n'(y)\,dy
   =\int_1^\infty\bigl(\psi(y)-y+1\bigr)f_n'(y)\,dy ,
\end{equation}
with $f_n$ the Laguerre test function recalled below. We record once and for all that the
boundary term $-n$ in \eqref{eq:lamprime} arises from $\psi(1)=0$ and must never be discarded;
an earlier version of this work discarded it and had to retract the consequence.

\begin{notation}[the Laguerre kernel]\label{not:kernel}
Write $L_n^{(\alpha)}$ for the generalized Laguerre polynomials. Put
\[
   E(u)=\psi(e^u)-e^u,\qquad K_n(u)=e^{-u}L^{(2)}_{n-1}(u),\qquad a=\log2 .
\]
The exponent $\alpha=2$ is not a free parameter: it is forced by the exact integration-by-parts
identity $\bigl(L_n^{(1)}\bigr)'=-L^{(2)}_{n-1}$, which is what converts \eqref{eq:lamprime}
into the compact form used below. Replacing it requires rebuilding every boundary term in the
chain.
\end{notation}

\subsection{The compact--tail identity and the inequality $(\star)$}\label{sec:a1}

\begin{theorem}[compact--tail decomposition]\label{thm:compact-tail}
For every $n\ge1$ and every $T>\log 2$,
\begin{equation}\label{eq:CnT}
   C_n(T)=\lambda_n-\tfrac14A_n-R_n(T),
   \qquad\text{where}\qquad
   R_n(T)=-\int_T^\infty E(u)K_n(u)\,du ,
\end{equation}
and $C_n(T)\ge0$ if and only if
\begin{equation}\label{eq:a1}
   \tag{$\star$}
   \int_{\log2}^{T}\bigl(\psi(e^u)-e^u\bigr)e^{-u}L^{(2)}_{n-1}(u)\,du
   \ \le\ q_n:=\tfrac34A_n+1-L_n^{(1)}(\log2).
\end{equation}
\end{theorem}

We shall refer to \eqref{eq:a1} throughout as \emph{the compact inequality}, and to the
statement that it holds for every $n\ge150$ with an admissible cutoff $T=T_n$ as $(\star)$
without further qualification. It is the sole unconditional assertion that the route recalled here
leaves to be proved: granted \eqref{eq:a1} for all $n\ge150$, together with the tail estimate of
\S\ref{sec:a0} and the finitely many verified indices $1\le n\le149$, Li's criterion
(Theorem~\ref{thm:li}) gives the Riemann Hypothesis. Everything negative in this paper is a
statement about attempts to prove $(\star)$.

The reserve $q_n$ must be used in this exact form. Its asymptotic $\tfrac38n\log n$ is correct
but has already caused one retraction in this work, because the difference between $q_n$
and its leading term is of the same order as the quantities being compared at moderate $n$.

There is a one-parameter family worth recording, because it clarifies that the coefficient
$\tfrac14$ is not a lucky accident and cannot be improved for free.

\begin{theorem}[the $\theta$-family]\label{thm:theta}
For $\theta\in(0,1)$ set
$C_n^\theta(T)=(1-\theta)A_n-n-\int_0^TE K_n$. Then
\begin{equation}\label{eq:theta}
   C_n^\theta(T)=\lambda_n-\theta A_n-R_n(T),
   \qquad
   q_{n,\theta}=(1-\theta)A_n+1-L_n^{(1)}(\log2),
\end{equation}
and $C_n^\theta(T)\ge0$ is equivalent to the inequality \eqref{eq:a1} with $q_n$ replaced by
$q_{n,\theta}$; we write $(\star_\theta)$ for that variant, so that
$(\star)=(\star_{1/4})$. Moreover, if $T_n(\theta)$ denotes the minimal admissible cutoff for which the
tail estimate $|R_n(T_n(\theta))|\le\theta A_n$ is available from a fixed effective prime
number theorem, then $\theta\mapsto T_n(\theta)$ is non-increasing, and the two values
$\theta,\theta'$ are connected by the \emph{signed transport}
\begin{equation}\label{eq:transport}
   \Delta_{n,\theta}=-\int_{T_n(1/4)}^{T_n(\theta)}E(u)K_n(u)\,du
   =R_n\bigl(T_n(\tfrac14)\bigr)-R_n\bigl(T_n(\theta)\bigr),
\end{equation}
whose net effect on the requirement satisfies the two-sided bracket
$-2\theta A_n\le\mathcal N_n(\theta)\le\tfrac12A_n$ with indeterminate sign.
\end{theorem}

The content of the bracket is that lowering $\theta$ raises the reserve and simultaneously
pushes $T_n(\theta)$ outwards, transferring tail mass into the compact block that
\eqref{eq:a1} must control. In the limit $\theta\to0$ the requirement tends to the full Li
inequality. This is a redistribution of the content of the Hypothesis, not a discount, and we
therefore fix $\theta=\tfrac14$ throughout, which is the value at which \eqref{eq:a1} is stated.
For definiteness the effective triple used to realize $T_n(\theta)$ is the
Johnston--Yang bound \cite{JohnstonYang}, Theorem~1.4, with
$A=1$, $U_0=1000$ and
$\eta(u)=0.1853\,u^{3/5}(\log u)^{-1/5}-1.801\log u$; the only feature of it that matters below
is the floor $U_0=1000$.

\subsection{The tail estimate, and the factor its published form discards}\label{sec:a0}

The tail $R_n(T_n)$ is controlled unconditionally by a Vinogradov--Korobov type zero-free region
\cite{Korobov,MTY}, in the following shape.

\begin{theorem}[unconditional tail estimate]\label{thm:a0}
For every admissible cutoff $T_n$ there is $0<B_n\le A_n$ with
\begin{equation}\label{eq:a0-raw}
   |R_n(T_n)|\ \le\ \frac{B_n}{4}\int_{T_n}^\infty(1+u)^{-2}\,du .
\end{equation}
\end{theorem}

The integral in \eqref{eq:a0-raw} evaluates exactly, and the published form of the estimate
discards the result. Retaining it costs nothing.

\begin{corollary}[the retained tail estimate]\label{cor:a0plus}
\begin{equation}\label{eq:a0plus}
   |R_n(T_n)|\ \le\ \frac{B_n}{4(1+T_n)}\ \le\ \frac{A_n}{4(1+T_n)} .
\end{equation}
\end{corollary}

\begin{proof}
$\int_{T_n}^\infty(1+u)^{-2}du=(1+T_n)^{-1}$, and $B_n\le A_n$.
\end{proof}

No hypothesis has been added and the Riemann Hypothesis has not been used. The commonly quoted
form $|R_n(T_n)|\le\tfrac14A_n$ is \eqref{eq:a0plus} with the factor $(1+T_n)^{-1}$ thrown away.
As we now show, that thrown-away factor is worth a factor of two in the required margin --- and,
more importantly for this paper, it is what moves the target across the boundary of the
semidefinite family.

\subsection{The margin ladder}\label{sec:ladder}

\begin{definition}\label{def:margin}
For real $r>1$ put $D_n^{[r]}:=r\lambda_n-A_n$.
\end{definition}

\begin{theorem}[the admissible exponent range]\label{thm:range}
Let $n\ge150$, let $T_n$ be an admissible cutoff, and suppose $D_n^{[r]}\ge0$. Then
$C_n(T_n)\ge0$, and hence $(\star)$ holds at $n$, provided
\begin{equation}\label{eq:rmax}
   r\ \le\ r_{\max}(T_n):=\frac{4(1+T_n)}{T_n+2} .
\end{equation}
\end{theorem}

\begin{proof}
From $D_n^{[r]}\ge0$ we get $\lambda_n\ge A_n/r$. Insert this and \eqref{eq:a0plus} into
\eqref{eq:CnT}:
\[
   C_n(T_n)\ \ge\ A_n\Bigl(\frac1r-\frac14-\frac{1}{4(1+T_n)}\Bigr).
\]
The bracket is nonnegative exactly when $\tfrac1r\ge\tfrac14+\tfrac1{4(1+T_n)}
=\tfrac{T_n+2}{4(1+T_n)}$, which is \eqref{eq:rmax}.
\end{proof}

The exponent range \eqref{eq:rmax} is governed by the cutoff, and the cutoff is not a fixed
quantity. This is the point on which the sharp form of the target turns, so we isolate it.

\begin{lemma}[admissibility is upward closed]\label{lem:upward}
Fix $n$ and $\theta$. The defining requirement for a cutoff --- that the declared effective prime
number theorem give $|E(u)|\le C\,e^{u-\eta(u)}$ with
$\eta(u)\ge(n+1)\log(1+u)+\log\bigl(3An^2/(\theta A_n)\bigr)$ for every $u\ge T$ --- is a
pointwise condition imposed on a ray. Hence if $T$ is admissible and $T'\ge T$, then $T'$ is
admissible.
\end{lemma}

\begin{proof}
The condition at $T'$ is the condition at $T$ restricted to a smaller ray.
\end{proof}

Since $r_{\max}(T)=4-\tfrac{4}{T+2}$ is strictly increasing in $T$, Lemma~\ref{lem:upward} gives
\begin{equation}\label{eq:sup4}
   \sup\bigl\{r_{\max}(T):T\ \text{admissible}\bigr\}=4,
   \qquad\text{and}\qquad r_{\max}(T)<4\ \text{for every admissible }T .
\end{equation}
The supremum is approached and never attained. Consequently the sharp form of the target is the
\emph{strict} inequality at the integer exponent four.

\begin{theorem}[the frontier is the exponent four]\label{thm:four}
\begin{enumerate}[label=\textup{(\roman*)},leftmargin=*]
\item If $\lambda_n>\tfrac14A_n$ for every $n\ge150$, then for each such $n$ there is an
admissible cutoff $T_n$ at which $(\star)$ holds; together with Theorem~\ref{thm:li} and the
verified indices $1\le n\le149$ this gives the Riemann Hypothesis.
\item Conversely, if for some $n$ the inequality $(\star)$ holds at arbitrarily large admissible
cutoffs, then $\lambda_n\ge\tfrac14A_n$.
\item For every admissible cutoff $T$, the semidefinite conclusion $D_n^{[4]}\ge0$ is
insufficient \emph{by itself to force $C_n(T)\ge0$ through the two-sided estimate
\eqref{eq:a0plus} at that prescribed $T$}, since $r_{\max}(T)<4$.
\end{enumerate}
\end{theorem}

\begin{proof}
(i) Put $\delta_n=\lambda_n-\tfrac14A_n>0$ and choose
$T_n=\max\bigl\{T_n(\tfrac14),\,A_n/(4\delta_n)\bigr\}$, admissible by Lemma~\ref{lem:upward}.
By Corollary~\ref{cor:a0plus},
$|R_n(T_n)|\le A_n/\bigl(4(1+T_n)\bigr)<A_n/(4T_n)\le\delta_n$, so
$C_n(T_n)=\delta_n-R_n(T_n)>0$, which is $(\star)$ by Theorem~\ref{thm:compact-tail}.
(ii) By \eqref{eq:a0plus}, $R_n(T)\to0$ as $T\to\infty$, so
$C_n(T)\to\lambda_n-\tfrac14A_n$, and a limit of nonnegative numbers is nonnegative.
(iii) Immediate from \eqref{eq:sup4} and \eqref{eq:rmax}.
\end{proof}

We may therefore state the sharp \emph{cofinal-cutoff sufficient target} as
\begin{equation}\label{eq:target}
   \boxed{\ D_n^{[4]}=4\lambda_n-A_n\ >\ 0
   \qquad (n\ge150),\ \ \text{unconditionally.}\ }
\end{equation}
The distinction matters.  If the non-strict inequality $D_n^{[4]}\ge0$ holds for every
$n\ge150$, then $\lambda_n\ge A_n/4>0$ there and Li's criterion, together with the verified
finite range, proves RH directly; no transfer to $(\star)$ is needed.  Strictness is required
only when one insists on realizing the compact certificate at a finite adapted cutoff using
\eqref{eq:a0plus}.

\begin{remark}[the declared instantiation, and why it is not the frontier]\label{rmk:1000}
It is instructive to see what a fixed choice of cutoff costs. The effective triple quoted in
\S\ref{sec:a1} is valid from $U_0=1000$ onwards, so $T_n\ge1000$ for every $n$, and taking that
floor gives the single exponent
\begin{equation}\label{eq:rstar}
   r_*:=r_{\max}(1000)=\frac{4\cdot1001}{1002}=\frac{2002}{501}=3.996007984\ldots ,
\end{equation}
valid simultaneously for all $n$; the corresponding sufficient condition is
$D_n^{[4]}\ge A_n/1001$. But $1000$ is only the threshold of validity of the imported estimate,
not the cutoff the estimate produces. The condition of Lemma~\ref{lem:upward} forces
$\eta(T)\gtrsim(n+1)\log T$, hence $T_n\gg n^{5/3}$ up to logarithmic factors; numerically the
minimal admissible cutoffs at $\theta=\tfrac14$ run from about $7.9\times10^{4}$ at $n=8$ to
about $5.5\times10^{8}$ at $n=1000$. Using the cutoff actually available rather than its floor
replaces the sufficient slack $A_n/1001$ by $A_n/(1+T_n)$, a reduction by the factor
$(1+T_n)/1001\approx5.5\times10^{5}$ at $n=1000$ and growing like $n^{5/3}$ thereafter. In
absolute terms the slack $A_n/(1+T_n)=O\bigl(n^{-2/3}\log n\bigr)$ tends to zero. What it never
does is vanish, by \eqref{eq:sup4}.
\end{remark}

Two remarks on the ladder. First, the descent $2\to3\to4^-$ is a genuine weakening at each
step: $2\lambda_n-A_n\ge0$ implies $3\lambda_n-A_n\ge\tfrac12A_n>0$, while the algebraic value
$\lambda_n=A_n/3$ satisfies the cubic margin and violates the quadratic one; and
$3\lambda_n-A_n\ge0$ implies $4\lambda_n-A_n\ge\tfrac13A_n>0$ but not conversely. Second, the
whole ladder is a statement about one number.

\begin{theorem}[fixed-ratio sufficient transfer]\label{thm:scalefree}
Let $c>\tfrac14$ be fixed and suppose $\lambda_n\ge c\,A_n$ for all $n\ge150$. If the admissible
cutoffs satisfy
\begin{equation}\label{eq:Tc}
   T_n\ \ge\ \frac{2-4c}{4c-1},
\end{equation}
then $(\star)$ holds for all $n\ge150$. Conversely, no value $c\le\tfrac14$ can force
$C_n(T_n)\ge0$ at a prescribed finite cutoff from the two-sided estimate
\eqref{eq:a0plus} alone, since the resulting lower bound has no sign.
\end{theorem}

\begin{proof}
Apply Theorem~\ref{thm:range} with $r=1/c$. The condition \eqref{eq:rmax} reads
$(1+T_n)(4c-1)\ge1$, i.e.\ $T_n\ge(4c-1)^{-1}-1=(2-4c)/(4c-1)$.
\end{proof}

\begin{corollary}\label{cor:1000}
The pair $\bigl(c,T\bigr)=\bigl(501/2002,\,1000\bigr)$ satisfies \eqref{eq:Tc} with equality:
$4c-1=\tfrac{2004}{2002}-1=\tfrac{1}{1001}$, so $(4c-1)^{-1}-1=1000$.
\end{corollary}

Corollary~\ref{cor:1000} says that the constant $2002/501$ and the cutoff floor $1000$ are the
same datum written twice.  Thus the following is a useful \emph{uniform fixed-ratio sufficient
condition}, but it is not the intrinsic frontier:
\begin{equation}\label{eq:frontier}
   \boxed{\ \exists\,c>\tfrac14\ \text{fixed}:\qquad
   \lambda_n\ \ge\ c\,A_n\quad\text{for all }n\ge150 .\ }
\end{equation}
The value $c=501/2002$ (equivalently the exponent $2002/501$) is one instantiation of
\eqref{eq:frontier}, tied to the coarse floor of the declared effective prime number theorem.
For the actual $n$-dependent cutoffs the necessary slack tends to zero, and allowing adapted
cofinal cutoffs gives Theorem~\ref{thm:four}.  No finite cutoff reaches $c=\tfrac14$ through the
unsigned transfer, but this does not make a fixed $c>\tfrac14$ logically necessary for RH.

\begin{remark}[what \eqref{eq:frontier} says about the primes]\label{rmk:75}
By \eqref{eq:split}, the inequality \eqref{eq:frontier} is
\[
   \lambda_n^{\mathrm{prime}}\ \ge\ -(1-c)A_n,
   \qquad 1-c\approx\tfrac34 ,
\]
that is: \emph{the prime side may not destroy more than about three quarters of the archimedean
block}. Under the Hypothesis the prime side destroys $O(\sqrt n\log n)$ of a block of size
$\sim\tfrac12n\log n$, i.e.\ asymptotically none of it. The requirement therefore carries a
cushion of a factor four against the truth. It carries no cushion at all against the available
tools, and that gap between the two --- generous against reality, hopeless against method --- is
the subject of this paper.
\end{remark}

\begin{remark}[the exact condition versus the sufficient one]\label{rmk:overshoot}
The exact form of $(\star)$ is $\lambda_n\ge\tfrac14A_n+R_n(T_n)$, with the true signed tail.
At a fixed $T_n$, replacing $R_n(T_n)$ by its ceiling from \eqref{eq:a0plus} gives a sufficient
condition. The overshoot is bounded: it lies
in $\bigl[0,\,A_n/(2(1+T_n))\bigr]\subseteq[0,A_n/2002]$, which is half of the slack
$A_n/1001$ appearing in Theorem~\ref{thm:threshold} only when one replaces $T_n$ by its floor.
For the actual cutoffs it is $o(1)$ in absolute size.  It would matter further only if a mechanism
could exploit the sign of $R_n(T_n)$, and
\S\ref{sec:mech-envelope} proves that no such mechanism is available from the tail estimate,
the prime number theorem and eventual monotonicity alone.
\end{remark}

\begin{remark}[a binding erratum, and why the labelling convention of
Convention~\ref{conv:cond} exists]\label{rmk:erratum}
An earlier stage of this work imported Lagarias's asymptotic for the Li coefficients of
automorphic $L$-functions \cite{Lagarias2007} with the sign of the incomplete block inverted.
With $k_n(s)=(1+1/s)^n-1$, the residue of $k_n(s)\,(-L'/L)(s+1)$ at $s=\rho-1$ is
$-m_\rho k_n(\rho-1)$, and retaining the orientation of the contour throughout gives, in the
case of $\zeta$,
\[
   \lambda_n=A_n+\lambda_n(\sqrt n)+O(\sqrt n\log n)
   \qquad\text{unconditionally},
\]
where $\lambda_n(T)=\sum_{|\Im\rho|<T}\bigl[1-(1-1/\rho)^n\bigr]$ is the incomplete Li
coefficient, \emph{not} $A_n-\lambda_n(\sqrt n)+O(\sqrt n\log n)$. Consequently the small
remainder is $\widetilde\eps_n:=A_n-\lambda_n+\lambda_n(\sqrt n)=O(\sqrt n\log n)$, and the
quantity $\eps_n^-:=A_n-\lambda_n-\lambda_n(\sqrt n)$, which the earlier treatment used as
small, satisfies $\eps_n^-=\widetilde\eps_n-2\lambda_n(\sqrt n)$ and has \emph{no}
unconditional $O(\sqrt n\log n)$ bound. Every corollary resting on such a bound was withdrawn.
The clean asymptotic $\lambda_n=\tfrac12n\log n+C_1n+O(\sqrt n\log n)$ without the incomplete
block holds only under the Riemann Hypothesis, and must not be quoted as unconditional. We mention this here
because it is the reason the present paper labels every import, and because it illustrates the
failure mode that the discipline of witnesses is designed to catch: an $O(\cdot)$ can hide
precisely the sign one is trying to establish.
\end{remark}

%% ======================================================================
\section{Two certificate types and the fixed-cutoff transfer}\label{sec:dichotomy}
%% ======================================================================

We now record a distinction that remains useful for auditing the sixteen mechanisms, but whose
scope must be stated precisely.  The quantitative identity in \S\ref{sec:threshold} concerns the
single sufficient transfer obtained from the floor $T\ge1000$.  It does not show that the master
criterion sits beyond all semidefinite arguments: $D_n^{[4]}\ge0$ for all large $n$ already
implies RH directly, and $D_n^{[4]}>0$ realizes $(\star)$ at adapted cofinal cutoffs.

\subsection{The two classes}

\begin{definition}[semidefiniteness certificate]\label{def:semidef}
A \emph{semidefiniteness certificate} for a quantity $X$ is a proof of $X\ge0$ obtained by
exhibiting $X$ as a value of an object belonging to one of the following classes, or as a
positive combination of such values:
\begin{enumerate}[label=\textup{(S\arabic*)},leftmargin=*]
\item a positive semidefinite quadratic form, Gram matrix, or reproducing kernel
      \cite{Aronszajn,HornJohnson};
\item a sum of squares in an algebra, or an algebraic square $C^2$ of an operator or symbol
      \cite{Reznick,Marshall,Scheiderer};
\item a Polya frequency or totally positive sequence or kernel, or a log-concave sequence
      \cite{Karlin,Schoenberg,BrandenSurvey};
\item an integral against a positive measure --- a Laplace, Stieltjes or Herglotz
      representation, a L\'evy measure, a compound-Poisson law, a completely monotone function
      \cite{Widder,SSV,Sato,SteutelVanHarn};
\item a stochastic or convolution order relation between two laws \cite{ShakedShanthikumar};
\item a nonnegative trigonometric or Fej\'er-type kernel, or a Carath\'eodory/Schur bound
      \cite{Fejer,Caratheodory,Toeplitz,Schur};
\item a graded complex with a positive metric, so that the object of interest is a norm or a
      Hodge Laplacian eigenvalue.
\end{enumerate}
\end{definition}

\begin{definition}[coercivity certificate]\label{def:coerc}
A \emph{coercivity certificate} for a quantity $X$ relative to a scale $S>0$ is a proof of
$X\ge c\,S$ with $c>0$ fixed and independent of the parameters. Concretely it is one of:
\begin{enumerate}[label=\textup{(C\arabic*)},leftmargin=*]
\item a spectral gap: $\inf\operatorname{spec}(\mathcal A)\ge c$ for a self-adjoint operator
      $\mathcal A$;
\item a Poincar\'e or log-Sobolev inequality, i.e.\ a lower bound on a Dirichlet form by a
      variance \cite{BGL,Gross};
\item a lower bound on a Rayleigh quotient for a \emph{fixed} vector, obtained from a
      non-alignment or quantitative-unique-continuation estimate;
\item a resolvent estimate $\|(\mathcal A-z)^{-1}\|\le C$ on a region, yielding quantitative
      invertibility rather than mere injectivity.
\end{enumerate}
\end{definition}

The two definitions are not exclusive as a matter of logic --- a spectral gap is, after all, a
semidefiniteness statement for $\mathcal A-c$ --- and we are not proposing a theorem of
mathematical logic. We are proposing an accurate description of what the available instruments
actually deliver when applied to the object at hand. The point is that in each of the sixteen
mechanisms below, the positivity being exploited is a positivity of an \emph{auxiliary}
object --- the L\'evy measure, the local coefficient sequence, the kernel, the metric --- and the
inequality that the mechanism produces about $\lambda_n$ has zero on its right-hand side. To get
a proportional margin at a prescribed cutoff one must either shift the auxiliary object by a
fixed amount (which is what the witnesses of \S\ref{sec:falsifiers} forbid for the listed
transfers) or estimate the target quantity directly.  Such a proportional margin is optional,
not logically necessary for the master criterion.

\subsection{The algebraic gap at the declared cutoff floor}\label{sec:threshold}

Here is the quantitative statement for the historical fixed-floor transfer. The exponent four is
where every instrument in class (S1)--(S3) naturally lives, because $4=2\cdot2$ and because the
completed cocycle admits the formal factorization
$H_u^4K_u^3=\bigl(H_u^2K_u^{3/2}\bigr)^2$ recorded in \S\ref{sec:sos}. Replacing the actual
cutoff by $1000$ pushes the convenient sufficient exponent strictly below four. The resulting
algebraic identity is exact.

\begin{theorem}[the fixed-floor threshold]\label{thm:threshold}
Put $\delta=4-r_*=2/501$. Then for every $n$,
\begin{equation}\label{eq:threshold-id}
   D_n^{[r_*]}=D_n^{[4]}-\delta\lambda_n
   =\frac{1001\,D_n^{[4]}-A_n}{1002},
\end{equation}
and consequently
\begin{equation}\label{eq:threshold}
   \boxed{\ D_n^{[r_*]}\ge0
   \qquad\Longleftrightarrow\qquad
   D_n^{[4]}=4\lambda_n-A_n\ \ge\ \frac{A_n}{1001}. \ }
\end{equation}
\end{theorem}

\begin{proof}
$D_n^{[r_*]}=r_*\lambda_n-A_n=(4-\delta)\lambda_n-A_n=D_n^{[4]}-\delta\lambda_n$. Since
$\lambda_n=(D_n^{[4]}+A_n)/4$ and $\delta/4=1/1002$,
\[
   D_n^{[r_*]}=D_n^{[4]}-\frac{D_n^{[4]}+A_n}{1002}
   =\frac{1002D_n^{[4]}-D_n^{[4]}-A_n}{1002}
   =\frac{1001D_n^{[4]}-A_n}{1002}.
\]
The equivalence \eqref{eq:threshold} is immediate.
\end{proof}

\begin{corollary}[the shortfall of the floor-$1000$ transfer]\label{cor:shortfall}
Any argument that tries to reach $(\star)$ at the declared floor solely by inserting
$D_n^{[4]}\ge0$ into the two-sided tail estimate falls short of that particular sufficient
condition by $A_n/1001$, and
\[
   \frac{A_n}{1001}=\frac{n\log n}{2002}\bigl(1+o(1)\bigr)\longrightarrow\infty .
\]
\end{corollary}

The corollary explains why the floor-$1000$ versions of the quartic transfer cannot be repaired
by a bounded correction: their declared deficit grows with the archimedean block.  It must not be
read as a deficit at the actual $n$-dependent cutoff, where the required slack is
$A_n/(1+T_n)=o(1)$, or as a deficit in Li's criterion, where $D_n^{[4]}\ge0$ already suffices.

\begin{remark}[the chronology of the mismatch]\label{rmk:chronology}
Before the retention of the factor $(1+T_n)^{-1}$ in Corollary~\ref{cor:a0plus}, the sufficient
target was the \emph{strong margin} $D_n^{[2]}=2\lambda_n-A_n\ge0$, and the exponent two is where
a Gram or sum-of-squares certificate would have delivered exactly the right conclusion. The
successive weakenings to $r=3$ and then to $r=r_*<4$ were real improvements in a uniform
fixed-cutoff sufficient condition.  They also obscured the stronger cutoff fact: $r_*$ was never
intrinsic.  Six mechanisms were evaluated against that unnecessarily strong benchmark.  Their
explicit internal falsifiers remain valid; the global diagnosis ``positivity cannot suffice''
does not.
\end{remark}

\subsection{What the title asserts, precisely}\label{sec:thesis}

Theorem~\ref{thm:threshold} and Corollary~\ref{cor:shortfall} are statements about one
historical transfer, and Remark~\ref{rmk:1000} has just deprived them of the global reading they
were originally given. The thesis of this paper does not rest on them. It rests on the following,
which is proved in \S\ref{sec:halflog2} and is independent of every choice of cutoff, exponent
and normalization.

\begin{proposition}[the object at stake is not a positive object]\label{prop:thesis}
Let $r>1$ and $\eps>0$ be fixed. The normalized completed cocycle
$\mathcal S_u^{[r]}=H_u^{r}K_u^{r-1}$ is not the Laplace transform of a positive measure for all
sufficiently small $u>0$; nor is it, along the correlated path $u=c\eps$, in the limit that
extracts the coefficient. Its arithmetic and Gamma factors \emph{are} positive for every real
exponent, by infinite divisibility \eqref{eq:nuA}, \eqref{eq:nuK}. The obstruction is carried
entirely by the polar factor and is quantitative: at $x_0=\tfrac12\log2$, strictly below the
first prime atom, the first variation of the completed inverse equals
\[
   e^{-\eps x_0}\Bigl(-r+\frac{r-1}{\sqrt2}\Bigr)
   =-e^{-\eps x_0}\Bigl[\Bigl(1-\tfrac1{\sqrt2}\Bigr)r+\tfrac1{\sqrt2}\Bigr]\ <\ 0 ,
\]
a defect bounded away from zero which worsens linearly in $r$.
\end{proposition}

\begin{proof}
Theorem~\ref{thm:halflog2}, equation \eqref{eq:diag-expansion} for the correlated path, and
Proposition~\ref{prop:not-near-miss} for the quantitative form.
\end{proof}

This is the exact sense in which semidefiniteness is not coercivity here. A semidefiniteness
certificate establishes $X\ge0$ by exhibiting $X$ as a member of a class every member of which is
nonnegative by construction; Definition~\ref{def:semidef} lists the classes that were tried.
Proposition~\ref{prop:thesis} says that the quantity at stake belongs to none of them, provably,
with a rational witness, and that the failure is not a small perturbation to be absorbed by a
sharper constant but a fixed proportion of the object. What is then left to prove is that a
difference of two divergent blocks of the same order has a definite sign with a proportional
margin --- a coercivity estimate in the sense of Definition~\ref{def:coerc}.

Two disclaimers, both of which we mean literally. This does not assert that no positivity
argument whatever can succeed: it asserts that the sixteen classes actually attempted certify
objects other than the one in question. And it does not assert that a proportional bound is
logically necessary for the Riemann Hypothesis; \S\ref{sec:cushion} and
Theorem~\ref{thm:density-block} show that it is not.

\subsection{The affine obstruction: why no Gram representation exists at all}

Before the catalogue, we isolate the most elementary reason a semidefinite certificate cannot
represent the target. It applies to any attempt to write the target inequality as the
nonnegativity of a quadratic form in the arithmetic data, and it is the abstract shell of what
happens concretely in \S\ref{sec:m1-global} and \S\ref{sec:m3}.

\begin{lemma}[affine functionals admit no Gram representation]\label{lem:affine-gram}
Let $\mathcal A(c)=q-\sum_{j=1}^mc_j$ be a non-constant affine functional of
$c=(c_1,\dots,c_m)\in\R^m$. Then there is no positive semidefinite
$G\in\R^{(m+1)\times(m+1)}$ with
\begin{equation}\label{eq:gram}
   \mathcal A(c)=\begin{pmatrix}1&c^{\!\top}\end{pmatrix}
   G\begin{pmatrix}1\\ c\end{pmatrix}
   \qquad\text{for all }c\text{ in a nonempty open set.}
\end{equation}
\end{lemma}

\begin{proof}
Write $G=\begin{pmatrix}G_{11}&G_{1c}\\ G_{c1}&G_{cc}\end{pmatrix}$. The right-hand side of
\eqref{eq:gram} is a polynomial of degree two in $c$ whose Hessian is $2G_{cc}$; the left-hand
side has vanishing Hessian, so $G_{cc}=0$. If $G\succeq0$ and $G_{cc}=0$, then every principal
$2\times2$ minor containing the first coordinate, namely
$G_{11}G_{jj}-G_{1j}^2=-G_{1j}^2$, must be nonnegative, forcing $G_{1c}=0$. Hence the
right-hand side is the constant $G_{11}$, contradicting non-constancy.
\end{proof}

\begin{remark}\label{rmk:affine-scope}
Lemma~\ref{lem:affine-gram} does not say that a positivity proof of $\mathcal A(c)\ge0$ is
impossible; it says that the proof cannot proceed by \emph{identically} rewriting $\mathcal A$ as
a positive form in the same variables. One may of course add signed terms to cancel the Hessian,
but then one is left with a difference of Gram forms, and controlling the negative part returns
one to $\mathcal A$ itself. One may also pass to nonlinear variables such as
$\sqrt{\Lambda(m)}$, but then cross terms appear which must vanish for independent weights,
leaving the oscillatory diagonal already refuted. Both escapes are executed and closed in
\S\ref{sec:m3}.
\end{remark}

%% ======================================================================
\section{Bibliographic context: where each instrument comes from}\label{sec:biblio}
%% ======================================================================

A catalogue of failures is only as trustworthy as its bibliographic hygiene. It is embarrassingly
easy to spend months rediscovering an obstruction that a published paper recorded decades ago, and
it is equally easy to claim novelty for an identity that is a special case of something standard. This
section therefore does two things. It situates each of the seven positivity classes of
Definition~\ref{def:semidef} in the literature it comes from, so that the reader can see which
classical theorem each mechanism was trying to invoke. And it states, class by class, what the
present account duplicates and what it does not.

The discipline behind this section was imposed on us by an earlier episode in which a mechanism was
developed at length before an external search revealed that its analytic content was already
published. Two routes were eliminated on those grounds in the course of the present work: the
scalar least-common-multiple form of the prime correlator, whose content is contained in Coffey's
effective work on the Li coefficients and the Stieltjes constants \cite{Coffey2005,Coffey2008}; and
the model-space/de~Branges route, whose relation between Weil positivity, model spaces and
de~Branges spaces \cite{deBranges} is covered by Suzuki \cite{Suzuki2025}, so that the novelty
claimed for it did not exist.

\subsection{The route itself}

The criterion used here is Li's \cite{Li1997}, with the complements of Bombieri and
Lagarias \cite{BombieriLagarias1999}; the coefficients had been introduced in the power-series form
by Keiper \cite{Keiper1992}. The literature on them divides sharply along the conditional line, and
Convention~\ref{conv:cond} exists because of it. Unconditional statements: Lagarias's asymptotic
with the incomplete block retained \cite{Lagarias2007}, whose sign we correct in
Remark~\ref{rmk:erratum}; the effective computations of Coffey \cite{Coffey2005,Coffey2008}; the
numerical surveys of Ma\'slanka \cite{Maslanka}. Conditional statements: the clean asymptotic
$\lambda_n=\tfrac12n\log n+C_1n+O(\sqrt n\log n)$, which requires the Hypothesis; Voros's
sharpenings of the criterion \cite{Voros2006,VorosBook}, whose conclusions concern variants that are
sensitive to the Hypothesis and do not transfer automatically to $(\star)$; the extension to the
Selberg class of Omar and Mazhouda \cite{OmarMazhouda}, which also indicates that a generalized base
point is admissible only if held fixed. Positivity conditions adjacent to Li's are surveyed by
Conrey and Li \cite{ConreyLi}. For general background we use \cite{Titchmarsh,Edwards,Ivic,
Davenport,MontgomeryVaughan,IwaniecKowalski,ConreyNotices,Bombieri,Sarnak}, and for the original
memoir \cite{Riemann}.

The identification of $(\star)$ with a Weil positivity statement
(Theorem~\ref{thm:weil-a1}) places the target inside Weil's framework \cite{Weil1952}, whose
function-field counterpart \cite{Weil1948,Deligne} is the source of the analogy that motivates the
whole positivity family; Yoshida's Hermitian forms \cite{Yoshida} and Burnol's invariant analysis
\cite{Burnol} are the closest classical treatments of the quadratic form itself, and Connes's trace
formula \cite{connes1999} together with the Connes--Consani archimedean positivity
\cite{ConnesConsani2021} and Riemann--Roch programme \cite{ConnesConsaniRR} are where the
positivity is expected to come from geometrically. Deninger's conjectural cohomology
\cite{Deninger} and ordinary Hodge theory \cite{Voisin} are the models for the graded-complex
mechanism of \S\ref{sec:koszul}, whose algebra is the standard Koszul construction
\cite{Eisenbud}.

\subsection{Class (S1)--(S2): Gram forms and sums of squares}

The Gram mechanisms of \S\ref{sec:mech-gram} invoke nothing more than the elementary theory of
reproducing kernels \cite{Aronszajn} and Sylvester's law of inertia \cite{HornJohnson}; the second
is what makes Remark~\ref{rmk:inertia} conclusive rather than merely suggestive. The indefinite
setting one is pushed into when the Hessian cancels is that of Krein spaces \cite{Bognar}, and the
relevant operator-theoretic background is \cite{ReedSimon,Simon,Davies}; the commutant lifting
apparatus \cite{FoiasFrazho} is the natural home of the interpolation form of the same question.

The sums-of-squares mechanism of \S\ref{sec:sos} is the arithmetic shadow of a very well understood
gap in real algebraic geometry: nonnegative is not the same as a sum of squares. Hilbert
\cite{Hilbert1888} settled the classification of the cases where they agree; Motzkin
\cite{Motzkin} and Choi and Lam \cite{ChoiLam} supplied the standard counterexamples;
Reznick \cite{Reznick} and Scheiderer \cite{Scheiderer} describe the general obstruction, and the
modern accounts are \cite{Marshall,PrestelDelzell,Blekherman}, with the Positivstellensatz results
of Schm\"udgen \cite{Schmudgen} and Putinar \cite{Putinar} and the computational side in
\cite{Lasserre,Parrilo}. Our No-go~\ref{ng:sos-limit} is not an instance of that classical gap ---
it is sharper and more specific, since the square genuinely exists and is destroyed by the
extraction limit --- but the classical gap is the reason one should have expected trouble.

\subsection{Class (S3): total positivity}

Total positivity is the tool that converts coefficient positivity into control of oscillation, and
it is exactly the tool Theorem~\ref{thm:pf2} removes. The theory is Karlin's \cite{Karlin}; the
characterization of Polya frequency sequences by their generating functions is due to Aissen,
Schoenberg and Whitney \cite{ASW} and Edrei \cite{Edrei}, with Schoenberg's smoothing operators
\cite{Schoenberg} as the variation-diminishing mechanism; matrix-theoretic treatments are in
\cite{FallatJohnson}. The log-concavity literature \cite{Stanley,BrandenSurvey} and the
P\'olya--Schur programme of Borcea and Br\"and\'en \cite{BorceaBranden} describe how far such
structures normally propagate through natural operations --- which is why their failure on the
exponent lattice, at the second minor and with the primes $2$ and $3$, is worth recording as a
witness rather than as an inconvenience.

\subsection{Class (S4)--(S5): positive measures, complete monotonicity, stochastic order}

The apparatus of \S\ref{sec:mech-measure} is entirely classical. Bernstein's theorem
\cite{Bernstein} and its modern treatment \cite{SSV}, Widder's Laplace transform \cite{Widder},
Hausdorff's moment sequences \cite{Hausdorff} and the moment-problem literature
\cite{Akhiezer,ShohatTamarkin} are what Corollary~\ref{cor:bernstein} contradicts. The infinite
divisibility of the Euler and Gamma channels is a L\'evy--Khinchine statement
\cite{Sato,SteutelVanHarn,Feller}, and the fact that the Gamma ratio is a generalized Gamma
convolution places it in Bondesson's class \cite{Bondesson}. The coupling
$X(u_2)=X(u_1)+Z$ is stochastic ordering in the sense of \cite{ShakedShanthikumar,MullerStoyan},
and Remark~\ref{rmk:order} is the observation that those orders act on monotone and convex test
classes, which is precisely the class the Laguerre kernel is not in.

The Stein--Mecke successor of \S\ref{sec:stein} rests on Mecke's equation \cite{Mecke}, in the form
given in \cite{LastPenrose}, and belongs to the Stein's-method circle of ideas
\cite{ChenGoldsteinShao,BarbourChen,NourdinPeccati}. We emphasize that
Theorem~\ref{thm:stein} is an exact identity in that tradition and is not itself the object of a
no-go; what dies is the inference from generator positivity, whose premise
Theorem~\ref{thm:stein-witness} falsifies.

\subsection{Class (S6): nonnegative kernels and Tauberian transfer}

The Fej\'er--Carath\'eodory mechanism of \S\ref{sec:fejer} is built from Fej\'er's kernel
\cite{Fejer}, Carath\'eodory's coefficient bound \cite{Caratheodory}, the Toeplitz--Herglotz
representation \cite{Toeplitz,Herglotz} and Schur's contractivity \cite{Schur}; the harmonic
analysis background is \cite{Zygmund,Katznelson}. The Abel and Ces\`aro mechanism of
\S\ref{sec:abel} lives in Tauberian theory \cite{HardyDivergent,Korevaar}, and the obstruction that
closes both --- a power series whose coefficients are eventually of one sign must have a real
singularity on its circle of convergence --- is Pringsheim's theorem, for which we follow Landau's
exposition \cite{Landau} and Remmert \cite{Remmert}. The recurring point, which we state once here
because it applies to both, is that the classical Tauberian conditions are conditions on
\emph{decay}, whereas an off-line quartet contributes growth of order $\cosh(na)$; every kernel
with decaying weights is therefore blind to the only event that matters.

\subsection{Coercivity: one optional stronger successor class}

The instruments of Definition~\ref{def:coerc} come from a different literature altogether:
Poincar\'e and logarithmic Sobolev inequalities \cite{Gross,BGL}, and the spectral theory of
semibounded operators \cite{ReedSimon,Davies}. It is worth noting that this literature is organized
around lower bounds by a fixed multiple of a scale and around the phenomenon that
Proposition~\ref{prop:no-gap} identifies: in the absence of a spectral gap, one may instead bound
a fixed vector's quotient. That is the technical content of Problem~\ref{prob:nonalign}.  The
density-relaxed criteria of \S\ref{sec:density-criterion} show that this is not the only possible
successor class.

\subsection{Analytic and arithmetic inputs used as inputs, not as mechanisms}

The effective prime number theory is imported and never re-derived: the Vinogradov--Korobov
zero-free region \cite{Vinogradov,Korobov} with Ford's explicit form \cite{Ford} and the current
explicit constants \cite{MTY}, the effective error term of Johnston and Yang
\cite{JohnstonYang} which supplies the triple realizing $T_n(\theta)$ in
Theorem~\ref{thm:theta}, and the verified height of Platt and Trudgian \cite{PlattTrudgian} with
B\"uthe's partial-hypothesis estimates \cite{Buthe} --- the last used only as verification, never as
an input to the uniform statement, since by design the proof of $(\star)$ for all $n\ge150$ may
not depend on finite zero localization. Large-value and density technology
\cite{GuthMaynard,Montgomery} does not enter, for a reason that deserves stating plainly: counting
zeros and locating them are different problems, and no density estimate excludes a single off-line
zero.

The arithmetic identities used are Chebyshev's $\Lambda=\mu*\log$ and Selberg's
$\mu*\log^2=\Lambda\log+\Lambda*\Lambda$ \cite{Selberg1949,Erdos1949,Apostol}; the Jordan totient
appearing in \eqref{eq:jordan} is classical, and its analytic quotients are studied in
\cite{MSSS}, so no novelty is claimed for \eqref{eq:jordan} itself --- only for the combination of
the continuous cocycle, the polar pairing and the Laguerre pullback. The Laguerre asymptotics that
underlie the stationary map \eqref{eq:stationary} are Szeg\H{o}'s \cite{Szego}, with the uniform
Plancherel--Rotach regime treated in \cite{Erdelyi,Olver} and its modern Riemann--Hilbert form in
\cite{DKMVZ}; we stress, as \S\ref{sec:m3} does, that the phase $2\sqrt{nu}$ is the small-ratio
approximation and not the uniform bulk phase, and that the falsifier of \S\ref{sec:m3} is
constructed to survive replacement by the correct uniform phase.

\subsection{Adjacent criteria we did not use, and why}

Three well-known reformulations are deliberately absent from the mechanisms above, and their absence
is itself a bibliographic finding.

The Nyman--Beurling criterion \cite{Beurling}, in B\'aez-Duarte's strengthened form
\cite{BaezDuarte}, is an exact equivalence; it therefore relocates the difficulty into an
approximation problem without weakening it, which is the pattern of circularity this paper is at
pains to avoid.

The de~Bruijn--Newman flow, with the non-negativity of the constant \cite{RodgersTao} and the
effective upper bounds \cite{PolymathDBN}, gives an inequality in the direction opposite to the one
needed; the same is true of the hyperbolicity/Jensen-polynomial reformulations
\cite{GORZ}, which are equivalences.

The Davenport--Heilbronn function \cite{DavenportHeilbronn} remains the standard external falsifier
for any mechanism that uses only functional symmetry, and is the model for the artificial divisors
of \S\ref{sec:falsifiers}. We use the algebraically simpler symmetric quartets instead, because
they can be placed at prescribed indices and uniformly inside the bulk, which the
Davenport--Heilbronn zeros cannot.

\subsection{Summary of the audit}

Duplicated, and therefore not claimed: the Jordan totient and its positivity; the Selberg identity;
the compound-Poisson representation of a multiplicative law; Bernstein's theorem and its use; the
Fej\'er--Carath\'eodory mechanism itself; the classical failure of nonnegativity to imply sums of
squares; the relation between Weil positivity, model spaces and de~Branges spaces; the effective
prime number theory. Not found in the literature in the combination used here, and therefore stated
as new in this paper: the global energy identity for the order kernel against $d\psi-dx$ and its
exact Hessian cancellation (Theorems~\ref{thm:max-energy}, \ref{thm:hessian-zero}); the local tower
differences \eqref{eq:tower-2}, \eqref{eq:tower-r} and the coupled operator square
\eqref{eq:tower-square}; the confluent hypergeometric form \eqref{eq:1F1} of the real polar
density; the signed L\'evy triple \eqref{eq:nu3} with its exact value at $\tfrac12\log2$; the
equivalence between the holomorphy of the regulated Cayley discs and the Hypothesis
(Theorem~\ref{thm:fejer-rh}); the fixed-floor identity \eqref{eq:threshold}, the cutoff-upward
closure and the cofinal exponent-four criterion; the density/block extraction of
Theorem~\ref{thm:syndetic-li} from the Bombieri--Lagarias dominant modes; and the sign correction
of Remark~\ref{rmk:erratum}, which is an audit of \cite{Lagarias2007} performed here and not an
official corrigendum.

%% ======================================================================
\section{Closed mechanisms I: order kernels, Gram forms and positive complexes}
\label{sec:mech-gram}
%% ======================================================================

The five mechanisms of this section all attempt to exhibit the prime-side functional as a value
of a positive semidefinite form. They are the most natural first attempts, because the target is
a quadratic-looking object and the arithmetic weights $\Lambda(m)$ are nonnegative. All five
fail, and --- this is the informative part --- they fail in three genuinely different ways: the
kernel is indefinite (\S\ref{sec:m1-local}); the kernel is definite but its Hessian is exactly
cancelled by the compensator it comes with (\S\ref{sec:m1-global}); or the representation is
forbidden for the affine reason of Lemma~\ref{lem:affine-gram} (\S\ref{sec:m3}). The last two
mechanisms fail because the positive structure one builds annihilates precisely the arithmetic
current it was built to carry (\S\S\ref{sec:m5-gram}--\ref{sec:koszul}).

\subsection{The order kernel $w(\max(i,j))$ and local positive semidefiniteness}
\label{sec:m1-local}

\subsubsection*{Admissible inputs}
Two levels are declared, and each use must say which it employs.
\emph{Finite algebraic level:} real numbers $q_1<\dots<q_M$ with $q_j>1$; weights $\ell_j>0$
with $\Psi_j=\sum_{i\le j}\ell_i$ and $d_j=q_j-1-\Psi_{j-1}-\ell_j/2$; arbitrary reals
$w_1,\dots,w_M$. \emph{Regulated arithmetic level:} $q_j$ runs over the prime powers with
$\ell_j=\Lambda(q_j)$; for $n\ge0$ and $\eps>0$,
\begin{equation}\label{eq:tau}
   \tau(x)=x^{-1-\eps}L_n(\log x),\qquad w(x)=\tau'(x),\qquad E(x)=\psi(x)-x+1 ,
\end{equation}
and the prime number theorem is used only to kill the boundary as $X\to\infty$ at fixed
$\eps>0$.

The mechanism starts from an exact identity which is worth stating on its own, because it is the
cleanest available algebraic form of the prime-side functional and because the finite boundary
term it contains is mandatory and easy to lose.

\begin{theorem}[regulated finite identity with its boundary]\label{thm:m1-finite}
Let $X>1$ not be a prime power, put $C_{n,\eps}(X)=-\int_1^XE(x)\tau'(x)\,dx$, and let
\[
   \mathcal H(X)=\sum_{q_j\le X}\ell_j(q_j-1)-\frac{\psi(X)^2}{2}.
\]
Then $\mathcal H(X)=-\int_1^XE-\tfrac12E(X)^2$, $d\mathcal H=\sum_j\ell_jd_j\delta_{q_j}$, and
\begin{equation}\label{eq:m1-finite}
   C_{n,\eps}(X)=\sum_{q_j<X}\ell_jd_jw(q_j)
   -\frac12\int_1^XE(x)^2\tau''(x)\,dx
   +\frac12E(X)^2w(X).
\end{equation}
\end{theorem}

\begin{proof}
Write $G(x)=\int_1^xE$. Since $G'=E$, one integration by parts gives
$C(X)=-G(X)w(X)+\int_1^XG\tau''$. Substituting $G=-\mathcal H-E^2/2$ and integrating
$-\int\mathcal H\tau''$ in the Stieltjes sense produces
$-\mathcal H(X)w(X)+\int w\,d\mathcal H$; the two boundary contributions combine to
$-[G(X)+\mathcal H(X)]w(X)=\tfrac12E(X)^2w(X)$, which is \eqref{eq:m1-finite}.
\end{proof}

At the finite algebraic level, grouping the double sum by its larger index gives the identity
that names the mechanism,
\begin{equation}\label{eq:m1-max}
   \sum_{j=1}^M\ell_jd_jw_j
   =\sum_{j=1}^M\ell_j(q_j-1)w_j
   -\frac12\sum_{i,j\le M}\ell_i\ell_j\,w_{\max(i,j)} ,
\end{equation}
and the hope is that $K(i,j)=w_{\max(i,j)}$ is positive semidefinite. For a globally
non-increasing sequence with limit $w_\infty\ge0$ it is, by the telescoping representation
\begin{equation}\label{eq:m1-telescope}
   w_{\max(i,j)}=w_\infty+\sum_{k\ge\max(i,j)}(w_k-w_{k+1})
   =w_\infty+\sum_k(w_k-w_{k+1})\mathbf1_{i\le k}\mathbf1_{j\le k},
\end{equation}
which exhibits $K$ as a nonnegative combination of rank-one projections onto prefix vectors. The
proposal was to split $\tau'$ into blocks of monotone variation and use \eqref{eq:m1-telescope}
on each decreasing block. The exact block decomposition is
\begin{equation}\label{eq:m1-blocks}
\begin{split}
   \frac12\sum_{i,j\le M}\ell_i\ell_jw_{\max(i,j)}
   =\sum_{b=1}^B\Bigl\{&P_{b-1}\sum_{j\in I_b}\ell_jw_j
   +\frac{w_{m_b}}{2}L_b^2
   +\frac12\sum_{k=m_{b-1}+1}^{m_b-1}(w_k-w_{k+1})L_{b,k}^2\Bigr\},
\end{split}
\end{equation}
with $I_b=(m_{b-1},m_b]$, $P_{b-1}=\sum_{i\le m_{b-1}}\ell_i$,
$L_{b,k}=\sum_{m_{b-1}<i\le k}\ell_i$ and $L_b=L_{b,m_b}$; inside a block one applies
$w_{\max(i,j)}=w_{m_b}+\sum_{k=\max(i,j)}^{m_b-1}(w_k-w_{k+1})$, and for $i$ in an earlier block
and $j\in I_b$ one has $w_{\max(i,j)}=w_j$, the two orientations of the pair absorbing the factor
$\tfrac12$.

\begin{nogo}[a decreasing block is not a positive form]\label{ng:m1-local}
The sentence ``each decreasing block gives a positive semidefinite form'' omits three
sign-deciding terms, and the remaining signed correlator has no fixed sign for the real von
Mangoldt weights.
\end{nogo}

\noindent Equation \eqref{eq:m1-blocks} exhibits exactly what is omitted: the rank-one term
$w_{m_b}L_b^2/2$, which is negative whenever $w_{m_b}<0$; the cross terms
$P_{b-1}\sum_{j\in I_b}\ell_jw_j$, which are unsigned; and the increasing blocks, where
$w_k-w_{k+1}<0$. Keeping the first moment and the max-form together does not repair this: by
Abel summation on the finite range,
\begin{equation}\label{eq:m1-abel}
   \sum_{j=1}^M\ell_jd_jw_j=\sum_{j=1}^{M-1}H_j(w_j-w_{j+1})+H_Mw_M,
   \qquad H_j=\sum_{i\le j}\ell_i(q_i-1)-\frac{\Psi_j^2}{2},
\end{equation}
which collapses the two pieces to a single signed correlator against the increments of $w$. The
discrete boundary term $H_Mw_M$ is mandatory; writing only the sum over a finite prefix is
false. And $H$ itself oscillates. Writing
$H(x)=\sum_{q\le x}\Lambda(q)(q-1)-\psi(x)^2/2$ for the arithmetic version, an exact rational
computation with outward-rounded prime logarithms certifies
\begin{equation}\label{eq:H-witness}
   H(2969)<-21,\qquad H(3167)>110 .
\end{equation}
Hence neither a global orientation of $H$, nor a global orientation of the increments
$w_j-w_{j+1}$, nor blockwise semidefiniteness can prove the target. The route closes because the
positivity was located in a factor that the arithmetic does not respect.

\begin{remark}[the kernel is not hiding a positive part]\label{rmk:inertia}
One might suspect that \eqref{eq:m1-blocks} is a lossy decomposition and that some finer
factorization of $K$ is positive. It is not. At ordered points $x_1<\dots<x_M$, with
$\delta_j=w_j-w_{j+1}$ for $j<M$ and $\delta_M=w_M$, and $v_j$ the prefix indicator vector,
\begin{equation}\label{eq:inertia}
   K=\sum_{j=1}^M\delta_j\,v_jv_j^{\!\top} ,
\end{equation}
and the matrix with columns $v_j$ is triangular and invertible. By Sylvester's law of inertia
\cite{HornJohnson}, the inertia of $K$ is \emph{exactly} the tally of signs among
$w_1-w_2,\dots,w_{M-1}-w_M,w_M$. So \eqref{eq:inertia} is the maximal factorization, and for the
Laguerre weight it is indefinite: with $a=1+\eps$ and $u=\log x$,
\[
   \tau''(x)=x^{-a-2}\bigl\{L_n''(u)-(2a+1)L_n'(u)+a(a+1)L_n(u)\bigr\},
\]
whose brace equals $\tfrac{n(n-1)}{2}+(2a+1)n+a(a+1)>0$ at $x=1$ and has leading term
$a(a+1)(-1)^nu^n/n!$ as $u\to\infty$. For every odd $n$ the weight changes sign.
\end{remark}

\subsection{The global energy of the same kernel, and an exactly vanishing Hessian}
\label{sec:m1-global}

The local no-go leaves open the possibility that the mechanism works globally: perhaps the
compensator $-\tfrac12\int E^2\tau''$ in \eqref{eq:m1-finite} is itself a positive energy, and
the sum of the two terms is coercive. It is a positive energy. It is the energy of the very same
kernel, and this is fatal.

Fix $X>1$, not a prime power, and $w\in C^1([1,X])$. Put
\begin{equation}\label{eq:measures}
   d\mu=\sum_{q_j<X}\ell_j\delta_{q_j},\qquad
   d\nu=\mathbf1_{[1,X]}(x)\,dx,\qquad
   d\sigma=d\mu-d\nu ,
\end{equation}
so that $E(x)=\sigma([1,x])=\psi(x)-x+1$, and define
\begin{align}
   Q_X(w;\sigma)&=\frac12\iint_{[1,X]^2}w(\max\{x,y\})\,d\sigma(x)\,d\sigma(y),
   \label{eq:QX}\\
   S_X(w;\mu)&=\int_{[1,X]}(x-1)w(x)\,d\mu(x)
   -\frac12\iint_{[1,X]^2}w(\max\{x,y\})\,d\mu(x)\,d\mu(y).
   \label{eq:SX}
\end{align}
For the von Mangoldt measure, \eqref{eq:SX} is exactly $\sum_{q_j<X}\ell_jd_jw(q_j)$, by
\eqref{eq:m1-max}.

\begin{theorem}[the compensator is the max-energy]\label{thm:max-energy}
With the conventions above,
\begin{equation}\label{eq:max-energy}
   Q_X(w;\sigma)=\frac12E(X)^2w(X)-\frac12\int_1^XE(x)^2w'(x)\,dx
   =\frac{w(X)}{2}E(X)^2+\frac12\int_1^X\bigl[-w'(t)\bigr]E(t)^2\,dt .
\end{equation}
In particular, for $w=\tau'$ the second term is precisely the compensator appearing in
\eqref{eq:m1-finite}, boundary included.
\end{theorem}

\begin{proof}
The pushforward of $\sigma\otimes\sigma$ under $(x,y)\mapsto\max\{x,y\}$ has cumulative function
$(\sigma\otimes\sigma)\{\max(x,y)\le t\}=E(t)^2$, an identity which absorbs the atomic diagonal
automatically. Hence $2Q_X=\int_{[1,X]}w(t)\,d\bigl(E(t)^2\bigr)$, and Stieltjes integration by
parts with $E(1)=0$ gives \eqref{eq:max-energy}.
\end{proof}

The second form in \eqref{eq:max-energy} is a genuine positive mixture of squares of prefix
masses whenever $w(X)\ge0$ and $-w'\ge0$. For the real Laguerre weight those hypotheses fail
uniformly, by Remark~\ref{rmk:inertia}. But the decisive fact is the following, which holds for
every finite measure and requires no arithmetic at all.

\begin{theorem}[global collapse of the mechanism]\label{thm:hessian-zero}
For every finite measure $\mu$ on $[1,X]$ --- not only the von Mangoldt measure ---
\begin{equation}\label{eq:collapse}
   S_X(w;\mu)+Q_X(w;\mu-\nu)=-\int_1^XE(x)w(x)\,dx ,
   \qquad E(x)=\mu([1,x])-(x-1).
\end{equation}
Consequently, for any finite signed measure $\eta$ with $F(x)=\eta([1,x])$ and every $t\in\R$,
\begin{equation}\label{eq:affine-exact}
   S_X(w;\mu+t\eta)+Q_X(w;\mu+t\eta-\nu)
   =S_X(w;\mu)+Q_X(w;\mu-\nu)-t\int_1^XF(x)w(x)\,dx ,
\end{equation}
so the second variation vanishes identically:
$\tfrac{d^2}{dt^2}\bigl\{S_X+Q_X\bigr\}=0$.
\end{theorem}

\begin{proof}
Expand \eqref{eq:QX} into its three components. The elementary integrals are
$\int_1^Xw(\max\{q,y\})dy=(q-1)w(q)+\int_q^Xw$ and
$\tfrac12\iint_{[1,X]^2}w(\max\{x,y\})dxdy=\int_1^X(y-1)w(y)dy$. Adding \eqref{eq:SX}, the
$\mu\times\mu$ terms cancel exactly, as do the terms $(q-1)w(q)$. By Fubini,
$\int\int_q^Xw(y)dy\,d\mu(q)=\int_1^X\mu([1,y])w(y)dy$, and the remaining sum is
$-\int_1^X\{\mu([1,y])-(y-1)\}w(y)dy$, which is \eqref{eq:collapse}. Since the right-hand side of
\eqref{eq:collapse} is affine in $\mu$, \eqref{eq:affine-exact} follows.
\end{proof}

\begin{nogo}[the two Hessians cancel, boundary included]\label{ng:m1-global}
The only arithmetic quadratic component of the mechanism cancels before any estimate is made. No
Gram form survives to provide coercivity; what remains is the original linear functional.
\end{nogo}

\noindent Concretely: in \eqref{eq:SX} the coefficient of $t^2$ is
$-\tfrac12\iint w(\max)\,d\eta\,d\eta$, and in \eqref{eq:QX} it is the exact opposite. The
boundary bookkeeping matters and is easy to get wrong: applying
Theorem~\ref{thm:max-energy} to $\eta$ gives
\begin{equation}\label{eq:hess-boundary}
   \iint_{[1,X]^2}w(\max\{x,y\})\,d\eta(x)\,d\eta(y)
   =w(X)F(X)^2-\int_1^XF(x)^2w'(x)\,dx ,
\end{equation}
and omitting the term $w(X)F(X)^2$ would manufacture a spurious residual Hessian. Thus
factorizing the compensator factorizes only \emph{one} of the two cancelling terms; keeping both
coupled, as the discipline adopted here requires, leaves neither a positive nor a negative
quadratic form.

There is a second, quantitative reason the route is closed. Once \eqref{eq:affine-exact} holds,
taking absolute values no longer estimates a collective energy; it estimates the original linear
functional. With $w=\tau'$ and the regulator removed, the pointwise Vinogradov--Korobov envelope
produces an absolute load $\mathcal B_n$ for which
\begin{equation}\label{eq:load}
   \mathcal B_n\ \ge\ \exp\Bigl(\tfrac23n\log n+\tfrac13n\log\log n-O(n)\Bigr),
\end{equation}
against a first-difference budget of size $O(\log n)$. So the absolute-value exit is not merely
inconclusive: it loses a superpolynomial factor. This is the factorial loss that any
term-by-term use of the nonnegativity of $\Lambda$ incurs, in its sharpest local form.

\subsection{The non-local zero--lobe Gram}\label{sec:m3}

\subsubsection*{Admissible inputs}
The Weil autocorrelation of the Laguerre test; the stationary map between zero heights and
Laguerre lobes; conjugation and the functional equation; the geometry of Laguerre oscillation.
Nothing about the real values $\Lambda(m)$.

Put $f_n(x)=\mathbf1_{[0,\infty)}(x)L^{(1)}_{n-1}(x)$, $f^{\#}(x)=e^x\overline{f(-x)}$, and
$h_n=f_n*f_n^{\#}$. With $w=(s-1)/s$ its Laplace transform is
\begin{equation}\label{eq:hn-symbol}
   \mathcal Lh_n(s)=\bigl(1-w^n\bigr)\bigl(1-w^{-n}\bigr)=2-w^n-w^{-n},
\end{equation}
and the completed explicit formula, with pole, Gamma, primes, conjugation and boundary retained
before evaluation, gives $\mathcal W_\zeta(h_n)=2\lambda_n$. Combining with \eqref{eq:theta}:

\begin{theorem}[$(\star)$ as a Weil autocorrelation]\label{thm:weil-a1}
For every $\theta\in(0,1)$ and every admissible $T$,
\begin{equation}\label{eq:weil-a1}
   C_n^\theta(T)=\tfrac12\mathcal W_\zeta\bigl(f_n*f_n^{\#}\bigr)-\theta A_n-R_n(T),
\end{equation}
so $(\star_\theta)$ holds if and only if the right-hand side of \eqref{eq:weil-a1} is
nonnegative.
\end{theorem}

This is already a non-local factorization: $h_n$ is an autocorrelation on the whole additive
line and the prime side has not been separated from the archimedean side. The trouble is that
the word ``autocorrelation'' makes the symbol positive only after one identifies the functional
involution $s\mapsto1-s$ with complex conjugation. On $s=\half+it$ one has $1-s=\bar s$ and
\eqref{eq:hn-symbol} becomes $|1-w^n|^2\ge0$; off that line $w^{-1}\ne\bar w$ and the reciprocal
product has no sign. The identification is equivalent to being on the critical line, which is
what one is trying to prove.

\begin{nogo}[the stationary map forgets the coordinate that decides the sign]\label{ng:m3}
The zero--lobe pairing retains the height of a zero and discards its abscissa. The discarded
coordinate is exactly the one that can make a quartet's contribution exponentially negative.
\end{nogo}

\noindent In the regime where the Laguerre phase is $2\sqrt{nu}$, a zero of positive height
$\gam$ contributes phase $\gam u-2\sqrt{nu}$ with derivative $\gam-\sqrt{n/u}$, giving the
stationary point
\begin{equation}\label{eq:stationary}
   \gam\longmapsto u_\gam=\frac{n}{\gam^2}.
\end{equation}
The abscissa $\beta$ does not enter \eqref{eq:stationary} at all; it enters only the amplitude
$e^{(\beta-1/2)u}$, i.e.\ in Cayley coordinates the modulus of
$w_\rho=1-1/\rho=e^{-a+i\vartheta}$. A non-critical quartet contributes to
$\mathcal W_\zeta(h_n)$ exactly
\begin{equation}\label{eq:quartet-8}
   8-8\cosh(na)\cos(n\vartheta),
\end{equation}
and \eqref{eq:stationary} discards $a$. Moreover the map is many-to-one: the count of relevant
zeros is $N(\sqrt n)\sim(\sqrt n/4\pi)\log n$ while there are $\asymp n$ Laguerre lobes.

The falsifier can be placed \emph{inside the bulk}, not at an edge, which is what makes it
decisive. Let $n\ge300$ with $100\mid n$, and set
\begin{equation}\label{eq:m3-quartet}
   \vartheta=\frac{\pi}{50},\qquad
   r=\frac{1+\cos\vartheta}{2},\qquad
   w=re^{i\vartheta},\qquad \rho=\frac{1}{1-w}.
\end{equation}
With $D=|1-w|^2=1-2r\cos\vartheta+r^2$ one computes
$\beta=(1-r\cos\vartheta)/D$, $\gam=r\sin\vartheta/D$, and
\[
   \beta-\tfrac12=\frac{1-r^2}{2D}>0,\qquad
   1-\beta=\frac{r(r-\cos\vartheta)}{D}>0 ,
\]
so $\half<\beta<1$ strictly. Completing to the quartet
$\{\rho,\bar\rho,1-\rho,1-\bar\rho\}$ preserves conjugation and the functional equation, and its
two members of positive height share the same $\gam$, so \eqref{eq:stationary} assigns them the
same lobe. Since $d\gam/dr=\sin\vartheta(1-r^2)/D^2>0$ on $(0,1)$, evaluating at the endpoints
$r=\cos\vartheta$ and $r\uparrow1$ gives $\cot\vartheta<\gam<\tfrac12\cot(\vartheta/2)$, and the
elementary inequalities $\sin x<x$, $\cos x>1-x^2/2$, $\tan x>x$, $3<\pi<22/7$ give
\[
   \cot\vartheta>\frac1\vartheta-\frac\vartheta2>\frac{175}{11}-\frac{11}{350}>15,
   \qquad
   \tfrac12\cot\tfrac\vartheta2<\frac1\vartheta=\frac{50}{\pi}<\frac{50}{3}<\sqrt{300}\le\sqrt n .
\]
Hence $\tfrac{9}{2500}n<u_\gam<\tfrac{1}{225}n$: the stationary point lies in a fixed compact
subinterval of $(0,4n)$, using neither the hard nor the soft Plancherel--Rotach edge. Finally,
with $a=-\log r>0$ and $100\mid n$ one has $n\vartheta=2\pi(n/100)$, so $\cos(n\vartheta)=1$ and
\eqref{eq:quartet-8} equals $8-8\cosh(na)<0$; repeating the quartet with multiplicity $M$
multiplies this by $M$ without changing any height. By contrast, the critical pair
$\half\pm i\gam$ with multiplicity two has the \emph{same} height multiset, hence the same
stationary data, and contributes $8\{1-\cos(n\phi)\}\ge0$. The inputs the map sees agree; the
signs do not. Therefore no sign rule depending only on heights, lobe occupancy, conjugation and
the functional equation can represent the completed form.

There is an independent algebraic closure on the prime side. At a finite cutoff, writing the
signed lobe charges as $c=(c_1,\dots,c_m)$, the inequality $(\star)$ is affine,
$\mathcal A(c)=q_{n,\theta}-\sum_jc_j$, and Lemma~\ref{lem:affine-gram} forbids a positive
semidefinite Gram representation of it outright. This is the abstract reason the mechanism has no
chance, and it is worth noticing that it is the same reason as Theorem~\ref{thm:hessian-zero}:
the target is affine in the arithmetic data, and semidefinite certificates are quadratic objects.

\subsection{M\"obius inversion as a positive adjoint}\label{sec:m5-gram}

\subsubsection*{Admissible inputs}
$\Lambda=\mu*\log$ in the Dirichlet convolution algebra; absolute convergence at
$a=1+\eps>1$; the Laguerre addition and generating formulas; pairing of the pole with the prime
term before $\eps\downarrow0$. Not admissible: a sign from the M\"obius coefficients separately,
a norm of each factor separately, the functional equation without the real weights, or a metric
constructed after assuming the desired positivity.

The exact structure this mechanism was built to exploit is genuine and worth recording. With
$a=1+\eps$, put
\begin{equation}\label{eq:MB}
   \mathfrak M_j(a)=\sum_{d\ge1}\frac{\mu(d)}{d^a}L_j(\log d),
   \qquad
   \mathfrak B_k(a)=\sum_{m\ge1}\frac{\log m}{m^a}L^{(-1)}_k(\log m),
\end{equation}
both absolutely convergent. The Laguerre addition formula
$L_n(x+y)=\sum_{j=0}^nL_j(x)L^{(-1)}_{n-j}(y)$ together with $\Lambda=\mu*\log$ gives the exact
degree convolution
\begin{equation}\label{eq:degree-conv}
   \sum_{r\ge2}\frac{\Lambda(r)}{r^a}L_n(\log r)
   =\sum_{j=0}^n\mathfrak M_j(a)\,\mathfrak B_{n-j}(a),
\end{equation}
with no absolute value and no Cauchy--Schwarz in the degree. In generating-function form, with
$t=z/(1-z)$,
\begin{equation}\label{eq:MB-gen}
   \sum_{j\ge0}\mathfrak M_j(a)z^j=\frac{1}{1-z}\cdot\frac{1}{\zeta(a+t)},
   \qquad
   \sum_{k\ge0}\mathfrak B_k(a)z^k=-\zeta'(a+t),
\end{equation}
and, pairing with the continuous term
$J_{n,\eps}=\int_1^\infty x^{-1-\eps}L_n(\log x)dx=(\eps-1)^n\eps^{-n-1}$,
\begin{equation}\label{eq:m5-paired}
   C_{n,\eps}=\sum_{j=0}^n\mathfrak M_j(1+\eps)\mathfrak B_{n-j}(1+\eps)-J_{n,\eps},
   \qquad
   \sum_{n\ge0}C_{n,\eps}z^n=-(1-z)\frac{F_\eps'(z)}{F_\eps(z)},
\end{equation}
where $F_\eps(z)=(\eps+t)\zeta(1+\eps+t)$. Both terms of \eqref{eq:m5-paired} diverge separately
as $\eps\downarrow0$: it is a paired identity, not a licence to estimate the pieces.

\begin{nogo}[the Dirichlet inverse cannot be a Hilbert adjoint]\label{ng:m5-gram}
Realizing M\"obius inversion as the adjoint of the divisor operator in a positive metric is
impossible already on the minimal divisor poset $\{1,p\}$; the true adjoint reintroduces the comb
of quotients $n/m$ that the holomorphic inversion avoids.
\end{nogo}

\noindent The structural reason is short. Dirichlet inversion is a holomorphic, unipotent
operation on the divisor poset: it is lower-triangular with ones on the diagonal in the natural
basis. A Hilbert square, by contrast, introduces the adjoint, and the adjoint of the divisor
operator involves the quotients $n/m$, i.e.\ the very off-diagonal structure whose cancellation
was the goal. On $\{1,p\}$, requiring $Z^{-1}=Z^{\dagger}$ in a positive metric forces the
metric to be degenerate. Filtering by Euler factors does not help, because the filtration is not
monotone for the Laguerre response; and separating the norms of $\mathfrak M$ and $\mathfrak B$
in \eqref{eq:degree-conv} recreates exactly the exponential loss \eqref{eq:load}. What survives
is a coordinate, not a mechanism: the degree convolution \eqref{eq:degree-conv} is an exact and
useful rewriting whose estimation is again the original problem.

\subsection{The positive Euler--Gamma Koszul complex}\label{sec:koszul}

\subsubsection*{Admissible inputs}
A finite set of primes $P$; the boolean divisor cube
$\mathcal D_P=\{\prod_{p\in S}p:S\subseteq P\}$; the Hilbert space
$V_P=\ell^2(\mathcal D_P)$ with its canonical orthonormal basis and the ordinary positive metric;
the quotient generators $N_r$ and exterior multiplication $\eps_r$.

The last of the Gram-type mechanisms is the most structurally ambitious: build a graded complex
whose Hodge Laplacian cancels the mixed prime channels, so that positivity of the metric
delivers the missing sign. The complex exists and does cancel them.

\begin{theorem}[the minimal positive complex]\label{thm:koszul}
On the boolean divisor cube the Koszul differential
\begin{equation}\label{eq:koszul}
   d_E=\sum_rN_r\otimes\eps_r
\end{equation}
satisfies $d_E^2=0$, with a positive ordinary metric --- no supertrace and no Krein form. On
$\{1,p,q,pq\}$ its Hodge Laplacian cancels the mixed channels $p/q$ and $q/p$ exactly.
\end{theorem}

\begin{nogo}[the complex annihilates the current it was built to carry]\label{ng:koszul}
The same differential makes the M\"obius--divisor connection null-homotopic, hence zero in
cohomology.
\end{nogo}

\noindent Because $N_r=d_E\iota_r+\iota_rd_E$ for the contraction $\iota_r$, and
\begin{equation}\label{eq:MdZ}
   M\,\delta Z=(\log p)N_p+(\log q)N_q ,
\end{equation}
the connection $M\delta Z$ is a sum of null-homotopic operators, hence induces zero on
cohomology. Already on the poset $\{1,p\}$, the hermitized channel one wishes to declare a
boundary is exactly the channel that must survive as the prime current. The complex satisfies the
quotient-cancellation requirement and fails the current-survival requirement, and it fails the
second before the Gamma budget can even appear.

\begin{remark}[scope]\label{rmk:koszul-scope}
This is a no-go for the divisor Koszul complex and, more generally, for every complex in which
each quotient generator $N_p$ is null-homotopic. It does not exclude a non-local differential
which kills only quadratic combinations of quotients without making the $N_p$ exact --- but such
an object would still have to \emph{construct}, not postulate, the exact Gamma term. Nor does it
speak to graded resolutions incorporating the prime powers $p^k$ as additional data; the model
here is the squarefree boolean cube.
\end{remark}

%% ======================================================================
\section{Closed mechanisms II: nonnegative arithmetic coefficients}\label{sec:mech-coeff}
%% ======================================================================

The mechanisms of \S\ref{sec:mech-gram} sought positivity in a kernel. Those of the present
section seek it where it undeniably is: in the coefficients of the arithmetic Dirichlet series
themselves. The Euler product is a genuine and inexhaustible source of nonnegativity, and it is
natural to hope that a hierarchy of nonnegative arithmetic functions, pushed through the Laguerre
transform, delivers a sign. It does not, and the reason is uniform across all three mechanisms:
the pole must be paired with the prime term before the regulator is removed, and pairing a
nonnegative hierarchy with the pole turns it into a \emph{difference} of two nonnegative
hierarchies.

\subsection{The Jordan cocycle and the hierarchy $\mu*\log^k\ge0$}\label{sec:jordan}

\subsubsection*{Admissible inputs}
The absolutely convergent series of $A_u(s)=\zeta(s-u)/\zeta(s)$ in $\Re s>1+u$; M\"obius
inversion; the Laguerre generating identity; Abel summation; the real order of
$F(s)=(s-1)\zeta(s)$ on $(1,\infty)$.

For $u\ge0$ and $\Re s>1+u$, M\"obius inversion gives
\begin{equation}\label{eq:jordan}
   A_u(s)=\frac{\zeta(s-u)}{\zeta(s)}=\sum_{m\ge1}\frac{J_u(m)}{m^s},
   \qquad
   J_u(m)=m^u\prod_{p\mid m}\bigl(1-p^{-u}\bigr)\ \ge\ 0 ,
\end{equation}
the Jordan totient in its analytic normalization; on a prime power, $J_u(p^0)=1$ and
$J_u(p^j)=(q-1)q^{j-1}$ with $q=p^u$. Differentiating \eqref{eq:jordan} in $u$ at $u=0$
generates the full hierarchy $\mu*\log^k\ge0$, of which the first two members are Chebyshev's
$\Lambda=\mu*\log$ and Selberg's $\mu*\log^2=\Lambda\log+\Lambda*\Lambda$. All of these are
nonnegative arithmetic functions, and their positivity is exact and unconditional.

\begin{nogo}[the polar pairing converts the hierarchy into a difference]\label{ng:jordan}
Keeping the pole coupled, as any correct treatment must, replaces the positive hierarchy by
$D_k-\tfrac{k}{s-1}D_{k-1}$, whose Cayley--Laguerre coefficients already change sign at low
degree; and completing the cocycle so as to restore positivity produces a Schur family whose
contractivity is equivalent to the Hypothesis.
\end{nogo}

\noindent The first jet of the coupled cocycle is exactly the regulated cost functional
$C_{n,\eps}$ of \eqref{eq:m5-paired}, and an exact rational computation with outward-rounded
intervals gives
\begin{equation}\label{eq:jordan-witness}
   C_1<0<C_6 ,
\end{equation}
so no sign or convexity statement can be read off coefficient by coefficient. Second, the
falsifier $X_a(s)=\xi(s+a)\xi(s-a)$ with $0<a<\half$ satisfies $X_a(1-s)=X_a(s)$, has all the
Euler-side positivity of the hierarchy in its half-plane of absolute convergence, and has zeros
at $\half\pm a+i\gam$, off the line of symmetry. Hence ``nonnegative Jordan coefficients plus the
functional equation'' cannot force the Cayley sign. Third, and this is the sharpest form, the
completed family $\xi(s-u)/\xi(s)$ is Schur for all $u\downarrow0$ if and only if the Hypothesis
holds, and its half-completion has coefficient exactly $-\tfrac12\Delta D_n$: the mechanism
returns the very gate it was supposed to open. There is no intermediate inequality anywhere in
the family.

\subsection{Selberg's positivity, the $O(x)$ summation, and the Riccati induction}
\label{sec:selberg}

\subsubsection*{Admissible inputs}
$b:=\Lambda\log+\Lambda*\Lambda=\mu*\log^2\ge0$; Selberg's effective summatory formula
$\sum_{m\le x}b(m)=x\log x-\ldots+O(x)$ in the classical form; partial summation; the Laguerre
degree recurrence; the Riccati identity for $Z^{-1}\delta Z$.

Write $a=1+\eps$, $t=z/(1-z)$, $s=a+t$, and
\[
   p_n=\sum_{m\ge2}\frac{\Lambda(m)}{m^a}L_n(\log m),\qquad
   j_n=\int_1^\infty x^{-a}L_n(\log x)\,dx=\frac{(\eps-1)^n}{\eps^{n+1}},\qquad
   c_n=p_n-j_n .
\]
Selberg's identity induces an exact recurrence in the Laguerre degree for $c_n$, and the Riccati
identity for the logarithmic derivative of the zeta germ is its differentiated form. Both are
exact and both are, by themselves, sign-free.

\begin{nogo}[a positive convolution with a matching cancellation is not a one-sided
inequality]\label{ng:selberg}
Splitting the double pole makes the centred Selberg coefficient cancel another term of size
$\eps^{-n-1}$ in the recurrence. The classical $O(x)$ summation controls both pieces only at that
same scale, so at $\eps=1/n$ its dual norm is at least $n(n-1)^n-1$ against a budget of size
$O(\log n)$.
\end{nogo}

\noindent The failure here is quantitative rather than structural, and the scale is worth naming
because it recurs: any exit that estimates the two colliding pieces separately pays
$\eps^{-n-1}$, which at the natural choice $\eps=1/n$ is $n^{n+1}$-sized. Compare
\eqref{eq:load}: the two losses are the same phenomenon in different coordinates. A residue
estimate strong enough to be sufficient would again contain the correlator $C_{n,\eps}$
explicitly, i.e.\ would already be the theorem. What this no-go closes is the \emph{autonomous}
use of Selberg--Riccati; it does not close the existence of a new signed theorem about the real
weights that estimates the coupled residue directly.

\subsection{$\mathrm{PF}_2$, $\mathrm{TP}_2$ and log-concavity on the exponent lattice}
\label{sec:pf2}

\subsubsection*{Admissible inputs}
Multiplicativity of the shifted divisor coefficients; the exact local formula on a prime tower;
the theory of Polya frequency and totally positive sequences \cite{Karlin,Schoenberg,BrandenSurvey};
the variation-diminishing property of $\mathrm{TP}$ kernels.

Put $b_u=J_u*J_u$, so that $A_u(s)^2=\sum_mb_u(m)m^{-s}$ with $b_u\ge0$, multiplicative, and
\begin{equation}\label{eq:bu-tower}
   b_u(p^k)=2(q-1)q^{k-1}+(k-1)(q-1)^2q^{k-2},\qquad q=p^u,\ k\ge1 .
\end{equation}
The hope is that the lattice sequence $k\mapsto b_u(p^k)r^k$, $r=p^{-1-\eps}$, is a Polya
frequency sequence, so that the Euler product acts as a variation-diminishing operator and the
oscillations of the Laguerre test cannot increase in number. This would be a very strong
structural statement, and it is false at the second minor.

\begin{theorem}[the exact $\mathrm{PF}_2$ minor]\label{thm:pf2}
With $a_k=b_u(p^k)r^k$,
\begin{equation}\label{eq:pf2}
   a_1^2-a_0a_2=r^2(Q-1)(Q-3),\qquad Q=p^u ,
\end{equation}
so the sequence is neither log-concave nor $\mathrm{PF}_2$ --- hence not
$\mathrm{PF}_\infty$ --- for every $1<Q<3$. For the general convolution power $J_u^{*r}$ the
corresponding minor is
\begin{equation}\label{eq:pf2-general}
   w_1^2-w_0w_2=\frac{\rho^2r}{2}(Q-1)\bigl((r-1)Q-(r+1)\bigr),
   \qquad w_k=\bigl(J_u^{*r}\bigr)(p^k)\rho^k .
\end{equation}
\end{theorem}

\begin{witness}[two real primes, exact rational minor]\label{wit:pf2}
Take $u=1$, $\eps=2$, and write $w(m)=b_1(m)m^{-3}$. Then
\[
   w(3)=\frac{4}{27},\qquad w(6)=\frac{1}{27},\qquad w(12)=\frac{5}{432},
\]
and the $\mathrm{PF}_2$ minor on the fibre $3\cdot2^k$ is
\begin{equation}\label{eq:pf2-witness}
   \det\begin{pmatrix}w(6)&w(12)\\ w(3)&w(6)\end{pmatrix}
   =w(6)^2-w(3)w(12)=-\frac{1}{2916}\ <\ 0 .
\end{equation}
Normalizing multiplies this by a positive constant and does not change its sign. The witness uses
the actual primes $2$ and $3$; it does not replace the arithmetic by an abstract divisor.
\end{witness}

\begin{witness}[the quartic minor]\label{wit:pf2-quartic}
With $r=4$, $p=2$, $u=\half$, $\eps=1$, so $\rho=\tfrac14$ and $Q=\sqrt2$, equation
\eqref{eq:pf2-general} gives
\begin{equation}\label{eq:pf2-quartic}
   w_1^2-w_0w_2=\frac{11-8\sqrt2}{8}\ <\ 0 ,
\end{equation}
since $121<128$. The fourth convolution power is not $\mathrm{PF}_2$ either.
\end{witness}

\begin{nogo}[total positivity is unavailable on the exponent lattice]\label{ng:pf2}
The local factors are not $\mathrm{TP}_2$, so the Euler product cannot be presented as a
variation-diminishing operator built from totally positive local blocks, and no
$\mathrm{PF}_\infty$, $\mathrm{TP}$ or decreasing-variation argument is available for the gate.
\end{nogo}

\noindent It is worth being clear about why this closure is more than a technical annoyance.
Variation diminution is the one classical tool that converts positivity of coefficients into
control of \emph{oscillatory} test functions, which is precisely what the Laguerre kernel is.
Theorem~\ref{thm:pf2} removes it. After that removal, the positivity of $b_u$ is a statement
about a measure with no bearing on the sign of its pairing with a sign-changing test, which is
the general shape of the mismatch this paper is about.

%% ======================================================================
\section{Closed mechanisms III: positive measures, complete monotonicity and stochastic order}
\label{sec:mech-measure}
%% ======================================================================

We come now to the family of mechanisms that seemed, for a while, the most promising, and whose
death is the most instructive. The completed cocycle is an infinitely divisible object: its
arithmetic factor is a compound-Poisson law over the additive semigroup generated by the
$\log p$, its Gamma factor has an explicitly positive L\'evy density, and the two are exactly
ordered in the shift parameter. If the \emph{whole} completed cocycle were the Laplace transform
of a positive measure, several classical routes --- Bernstein's theorem, stochastic ordering,
Stein--Mecke generator positivity --- would apply at once. It is not, and the obstruction sits at
one exactly computable point.

\subsection{The exact compound-Poisson structure}\label{sec:cp}

\subsubsection*{Admissible inputs}
Independence of the prime exponents; the exact local generating function; the L\'evy--Khinchine
representation for compound-Poisson laws \cite{Sato,SteutelVanHarn}; Mecke's equation
\cite{LastPenrose}; Bernstein's theorem on completely monotone functions \cite{Widder,SSV};
stochastic and convolution order \cite{ShakedShanthikumar}.

Fix $0<u<\eps$, $s_0=1+\eps$, and normalize the arithmetic coefficients to a probability law
\begin{equation}\label{eq:cp-law}
   \mathbb P(N=m)=\frac{b_u(m)\,m^{-1-\eps}}{A_u(1+\eps)^2} .
\end{equation}
Since $b_u$ is multiplicative, writing $N=\prod_pp^{K_p}$ makes the exponents $K_p$ independent,
and with $r=p^{-1-\eps}$, $Q=p^u$, $R=Qr$ the normalized local generating function is
\begin{equation}\label{eq:cp-local}
   G_p(z)=\mathbb Ez^{K_p}
   =\Bigl(\frac{1-R}{1-r}\Bigr)^{\!2}\Bigl(\frac{1-rz}{1-Rz}\Bigr)^{\!2},
\end{equation}
each unsquared factor being the generating function of a zero-inflated geometric law:
$\mathbb P(X_p=0)=\tfrac{1-R}{1-r}$ and
$\mathbb P(X_p=k)=\tfrac{1-R}{1-r}(R-r)R^{k-1}$ for $k\ge1$, with total mass exactly one. Taking
logarithms in \eqref{eq:cp-local},
\begin{equation}\label{eq:cp-levy}
   \log G_p(z)=2\sum_{j\ge1}\frac{R^j-r^j}{j}\bigl(z^j-1\bigr),
\end{equation}
so $X=\log N$ is compound-Poisson with L\'evy measure
\begin{equation}\label{eq:nuA}
   \nu_A=2\sum_p\sum_{j\ge1}
   \frac{p^{-j(1+\eps-u)}-p^{-j(1+\eps)}}{j}\,\delta_{j\log p}\ \ge\ 0 .
\end{equation}
This is a product over \emph{all} primes: the interactions between distinct primes are produced
by exponentiating the sum, not discarded tower by tower. Moreover if $0<u_1<u_2<\eps$ then
$\nu_A(u_2)-\nu_A(u_1)\ge0$, so one may couple
$X(u_2)=X(u_1)+Z_{u_1,u_2}$ with an independent nonnegative compound-Poisson increment. The
arithmetic law is thus increasing in convolution order, hence in stochastic order.

\begin{remark}[why the ordering does not transfer]\label{rmk:order}
Stochastic order orders expectations of \emph{monotone} functions, and its convex refinements
order convex ones. The Laguerre tests of the gate are neither: they change sign $n$ times. The
ordering \eqref{eq:nuA} is a true and exact structural fact about the arithmetic factor which is
simply orthogonal to the question being asked. This is the mismatch in its mildest form: the
positive structure exists, is exactly identified, and does not couple to the target.
\end{remark}

The corresponding Gamma channel is also positive. The normalized Gamma ratio
$K_u(s)=\pi^{u/2}\Gamma(1+(s-u)/2)/\Gamma(1+s/2)$ is the Laplace transform of
$\eta=-\tfrac12\log V$ with $V\sim\mathrm{Beta}\bigl(\tfrac{3+\eps-u}{2},\tfrac u2\bigr)$, and
\begin{equation}\label{eq:nuK}
   \nu_K(x)\,dx=\frac{e^{-(3+\eps-u)x}-e^{-(3+\eps)x}}{x\bigl(1-e^{-2x}\bigr)}\,dx\ \ge\ 0 .
\end{equation}
Both channels are infinitely divisible for \emph{every} real exponent, not only integers, which
is what allows the near-quartic exponent $r_*$ of \eqref{eq:rstar} to be used at all:
$A_u^r$ has L\'evy measure $r\,\nu_A/2$ and $K_u^{r-1}$ has $(r-1)\nu_K$.

\subsection{The polar channel is signed, and no value of the correlation parameter repairs it}
\label{sec:polar}

The third channel is the pole, and it is where all the sign lives. Write
$P_u(s)=1-\tfrac{u}{s-1}$.

\begin{theorem}[exact polar densities]\label{thm:polar}
In the coordinate $t=s-1-\eps$, after the substitutions $u=c\eps$ and $r=\eps y$:
\begin{enumerate}[label=\textup{(\roman*)},leftmargin=*]
\item the squared pole $P_u^2$ has inverse Laplace transform
      $\delta_0(dr)+c(cr-2)e^{-r}dr$, whose continuous density is strictly negative on
      $0<r<2/c$ for \emph{every} $c>0$;
\item the cubed pole $P_u^3$ has continuous density
      $c\bigl(-3+3cr-\tfrac12c^2r^2\bigr)e^{-r}dr$, negative near $r=0$ and for large $r$, for
      every $c>0$.
\end{enumerate}
\end{theorem}

\begin{proof}
Expand $(1-u/(s-1))^k$ in powers of $(s-1)^{-1}$ and use
$(\eps+t)^{-(j+1)}=\int_0^\infty\tfrac{y^j}{j!}e^{-\eps y}e^{-ty}dy$ termwise; substitute
$u=c\eps$, $r=\eps y$. For (ii) the quadratic $-\tfrac12c^2r^2+3cr-3$ has roots
$cr=3\pm\sqrt3$.
\end{proof}

For real, non-integer exponents the polar density does not truncate, but it sums in closed form,
which is the technical fact that made the near-quartic reduction usable.

\begin{theorem}[the real polar density is a confluent hypergeometric]
\label{thm:polar-1F1}
For real $r>0$ and $s-1=\eps+t$,
\[
   \Bigl(1-\frac{u}{s-1}\Bigr)^{\!r}
   =\sum_{j\ge0}(-1)^j\binom rj\frac{u^j}{(\eps+t)^j},
\]
whose inverse Laplace transform is
$\delta_0(dy)+e^{-\eps y}\sum_{j\ge1}(-1)^j\binom rj\frac{u^jy^{j-1}}{(j-1)!}\,dy$. After
$u=c\eps$ and $q=\eps y$ the continuous part sums exactly to
\begin{equation}\label{eq:1F1}
   -\,r\,c\,e^{-q}\;{}_1F_1(1-r;2;cq)\;dq ,
\end{equation}
a convergent identity for all $q$ --- not an asymptotic expansion. For $3<r<4$ the channels
$j=1,3$ of the expansion are negative and $j=2$, $j\ge4$ are positive; \eqref{eq:1F1} is
negative at $q=0$.
\end{theorem}

\begin{remark}[why separating channels is forbidden]\label{rmk:no-split}
The divergence that the full combination must cancel is visible already in the $m=1$ Euler
subchannel, with the Gamma factor omitted:
\begin{equation}\label{eq:m1-divergence}
   [z^1]\frac{1}{1-z}\Bigl(1-\frac{u}{\eps+z/(1-z)}\Bigr)^{\!2}
   =(1-c)^2+\frac{2c(1-c)}{\eps},
\end{equation}
which diverges like $\eps^{-1}$ for $0<c<1$; restoring the Gamma factor changes this only by
$O(1)$, since $\partial_s\log K_u(s)=O(u)$. Hence separating $m=1$ from $m>1$ destroys exactly
the cancellation that makes the limit finite. Every no-go in this section respects the rule that
the three channels stay coupled until after the estimate.
\end{remark}

\subsection{Complete monotonicity fails at $x_0=\tfrac12\log2$}\label{sec:cm}

We now come to the witness that recurs, in five different disguises, throughout the rest of the
paper. It is worth stating in the most concrete available form first.

\begin{theorem}[exact inverse density of the completed cocycle]\label{thm:cm}
Take the rational parameters $u=1$, $\eps=2$, $c=\tfrac12$. Then
$V\sim\mathrm{Beta}(2,\tfrac12)$, the normalized density of $\eta=-\tfrac12\log V$ is
\begin{equation}\label{eq:beta-density}
   f(x)=\frac32\,\frac{e^{-4x}}{\sqrt{1-e^{-2x}}},\qquad x>0 ,
\end{equation}
the inverse transform of the polar factor is $\delta_0-2e^{-2x}dx+xe^{-2x}dx$, and the
absolutely continuous part of the convolution is, with $Q=\sqrt{1-e^{-2x}}$,
\begin{equation}\label{eq:h-exact}
   h(x)=\frac32e^{-2x}\Bigl(\frac1Q-4Q+\artanh Q\Bigr),
\end{equation}
using the exact evaluations
$\int_0^xe^{-2y}f(x-y)dy=\tfrac32e^{-2x}Q$ and
$\int_0^xye^{-2y}f(x-y)dy=\tfrac32e^{-2x}\bigl(\artanh Q-Q\bigr)$.
At the exact point $x_0=\tfrac12\log2$ one has $Q=1/\sqrt2$ and
\begin{equation}\label{eq:cm-witness}
   \frac1Q-4Q+\artanh Q=-\sqrt2+\artanh\tfrac1{\sqrt2}\ <\ 0 ,
\end{equation}
because $\artanh q=q\int_0^1\frac{ds}{1-q^2s^2}<2q=\sqrt2$ at $q=1/\sqrt2$.
\end{theorem}

\begin{corollary}[Bernstein fails for the real completed cocycle]\label{cor:bernstein}
The arithmetic law \eqref{eq:cp-law} has positive mass at $X=0$ and no support in
$0<X<\log2$. Since $0<x_0<\log2$, the inverse Laplace density of the \emph{complete} cocycle at
$x_0$ equals $h(x_0)$ times the mass at zero times positive normalizing constants, and is
therefore strictly negative. By uniqueness of the Laplace transform for finite signed measures,
no other positive measure represents it. Hence, by Bernstein's theorem,
\begin{equation}\label{eq:not-cm}
   t\longmapsto\frac{\mathcal S_1(3+t)}{\mathcal S_1(3)}
   \quad\text{is not completely monotone on }t\ge0 .
\end{equation}
\end{corollary}

\begin{nogo}[no positive global representation, hence no Bernstein, no stochastic order, no
$\mathrm{TP}$]\label{ng:cm}
Three routes close simultaneously: (i) $\mathrm{PF}_\infty$, $\mathrm{TP}$ and decreasing
variation obtained by factorizing the product over prime coordinates, killed by
Theorem~\ref{thm:pf2} and Witness~\ref{wit:pf2}; (ii) representing the completed cocycle as the
transform of a positive probability law and applying stochastic order, killed by
Corollary~\ref{cor:bernstein}; (iii) transferring the increasing coupling
$X(u_2)=X(u_1)+Z$ directly to the Laguerre coefficients, killed because the polar factor is a
deconvolution and the oscillatory test does not preserve the order.
\end{nogo}

\noindent We emphasize what is \emph{not} concluded. Corollary~\ref{cor:bernstein} does not imply
that any coefficient of the gate has the wrong sign; the gate may well hold. What it excludes is a
proof of the gate through positivity of the completed generator.

\subsection{Stein--Mecke: the identity that keeps every other prime}\label{sec:stein}

The successor left alive by \S\ref{sec:cm} was the one global identity that does not require
positivity anywhere: Mecke's equation for a compound-Poisson law, which retains inside the
expectation all primes other than the selected jump. Executed on the cubic cocycle it produces a
genuinely global recurrence, and then dies for the same reason at the same point.

Let $d_u=J_u*J_u*J_u\ge0$, $s_0=1+\eps$, $u=c\eps$, and let $\mu_3$ be defined by
$\mathcal L\mu_3(t)=\mathcal S_u^{[3]}(s_0+t)/\mathcal S_u^{[3]}(s_0)$ where
$\mathcal S_u^{[3]}=H_u^3K_u^2$. Its L\'evy exponent has the \emph{signed} kernel
\begin{equation}\label{eq:nu3}
   \nu_3=\nu_A^{[3]}+2\nu_K-3\nu_P ,
\end{equation}
where $\nu_A^{[3]}=3\sum_p\sum_{j\ge1}\tfrac{R_p^j-r_p^j}{j}\delta_{j\log p}$,
$\nu_K$ is \eqref{eq:nuK}, and
$\nu_P(y)dy=\tfrac{e^{-(\eps-u)y}-e^{-\eps y}}{y}dy$.

\begin{theorem}[global Stein--Mecke recurrence]\label{thm:stein}
For polynomial tests,
\begin{equation}\label{eq:mecke}
   \int x\varphi(x)\,d\mu_3(x)
   =\int_0^\infty y\int\varphi(x+y)\,d\mu_3(x)\,\nu_3(dy),
\end{equation}
justified by truncating $\nu_A^{[3]}$ and $\nu_K$ to $y\ge\delta$, differentiating the truncated
transform, and letting $\delta\downarrow0$; the passage is dominated because for each fixed
degree $k$,
$\int_0^\infty y(1+y^k)\,|\nu_3|(dy)<\infty$ --- the singularity $\nu_K(y)\asymp u/(2y)$ has
infinite activity, but the factor $y$ makes it integrable. Setting
$a_k=\int L^{(1)}_k(x)d\mu_3(x)$ and $a_{-1}=0$, the Laguerre recurrence gives
\begin{equation}\label{eq:stein-rec}
   (2k+2)a_k-(k+1)a_{k+1}-(k+1)a_{k-1}
   =\int y\int L^{(1)}_k(x+y)\,d\mu_3(x)\,\nu_3(dy).
\end{equation}
\end{theorem}

\begin{theorem}[the generator is negative before the first prime atom]
\label{thm:stein-witness}
Take $\eps=2$, $u=1$, $c=\tfrac12$ and $y_0=\tfrac12\log2$. Since $0<y_0<\log2$, the measure
$\nu_A^{[3]}$ has no support in a neighbourhood of $y_0$, and the continuous part of
\eqref{eq:nu3} there is
\begin{equation}\label{eq:stein-neg}
   \bigl(2\nu_K-3\nu_P\bigr)(y_0)=\frac{\tfrac52-2\sqrt2}{y_0}\ <\ 0 ,
\end{equation}
exactly, because $2\sqrt2>\tfrac52$ is $8>\tfrac{25}{4}$. By continuity $\nu_3$ is strictly
negative on an open interval contained in $(0,\log2)$. Moreover, at fixed $\eps>0$ and
$u\downarrow0$,
\begin{equation}\label{eq:kappa}
   \frac{\nu_3}{u}\longrightarrow\kappa_\eps(dy)
   =3e^{-(1+\eps)y}\,d\psi(e^y)-3e^{-\eps y}\,dy
   +\frac{2e^{-(3+\eps)y}}{1-e^{-2y}}\,dy ,
\end{equation}
and along the diagonal path $u=c\eps$, $\eps\downarrow0$ the same convergence holds locally, with
continuous density at $y_0$ equal to
\begin{equation}\label{eq:kappa-neg}
   \frac{2e^{-3y_0}}{1-e^{-2y_0}}-3=\sqrt2-3\ <\ 0 .
\end{equation}
\end{theorem}

\begin{nogo}[generator positivity plus jump ordering is not available]\label{ng:stein}
The implication ``positivity of the completed L\'evy measure and ordering of the jumps
$\Rightarrow$ the gate'' has a false premise, by Theorem~\ref{thm:stein-witness}, and the defect
persists exactly in the regime $u=c\eps$, $\eps\downarrow0$ that extracts
$3\lambda_n-A_n$.
\end{nogo}

\noindent The transparent form of \eqref{eq:kappa} is
$\text{primes}-\text{pole}+\text{Gamma}$, and integrating it against the Laguerre response of
\eqref{eq:stein-rec} and then taking the Abel limit recomposes $3\lambda_n-A_n$; integrating the
first term by parts against $d\psi(e^y)$ returns $\psi(e^y)-e^y$, i.e.\ the weighted correlation
of $(\star)$. Stein--Mecke, in other words, is an exact and non-local reorganization which
adds no arithmetic input of its own. What it leaves open --- and this is one of the two live
successors identified in \S\ref{sec:successor} --- is a \emph{specific signed non-local
cancellation} for the Laguerre tests against the real masses $\Lambda(p^j)$.

\begin{remark}[a note on notation]\label{rmk:kappa-scope}
$\kappa_0$ does not denote a measure of finite total variation on the half-line. Every global
recomposition against Laguerre polynomials in this section is to be read as the paired Abel limit
$\eps\downarrow0$, never as an ordinary integral against $\kappa_0$.
\end{remark}

%% ======================================================================
\section{Closed mechanisms IV: analytic positivity}\label{sec:mech-analytic}
%% ======================================================================

The mechanisms of this section abandon arithmetic positivity and look for the sign in function
theory: radial Abel summation, real order on a ray, nonnegative Fej\'er kernels, and the
Carath\'eodory bound $\Re\Phi\le1$. They are the shortest available routes, in the sense that
each of them would deliver the target from a single classical theorem. Two of them die because
the positivity they establish is a positivity of averages that does not transfer to individual
coefficients; the other two die because the analytic hypothesis they require --- holomorphy or
contractivity on a full disc --- is itself equivalent to the Riemann Hypothesis. That second
failure mode is worth naming, because it is the exact analytic shadow of the circularity that the
subject produces repeatedly: a positivity criterion whose hypothesis contains its conclusion.

\subsection{Radial Abel positivity of the incomplete block}\label{sec:abel}

\subsubsection*{Admissible inputs}
The exact contribution of a quartet to the Li coefficients; radial Abel summation; Ces\`aro and
Fej\'er triangular means of fixed order; no information about the real weights.

After the sign erratum of Remark~\ref{rmk:erratum}, the correct coordinate for the target is
\begin{equation}\label{eq:abel-coord}
   (\star_\theta)\iff
   -\lambda_n(\sqrt n)+\widetilde\eps_n+R_n\bigl(T_n(\theta)\bigr)\le(1-\theta)A_n ,
\end{equation}
so what is needed is a \emph{lower} bound for the incomplete Li coefficient
$\lambda_n(\sqrt n)$. An average cannot substitute for it without a signed Tauberian theorem, and
the mechanism of this subsection was the attempt to supply one from radial positivity.

Let $\rho=\beta+i\gam$ with $0<\beta<1$, $\gam\ne0$, $\beta\ne\half$, let
$\mathcal O(\rho)=\{\rho,\bar\rho,1-\rho,1-\bar\rho\}$, and put $w=w_\rho=re^{i\vartheta}$, so
that the symmetries give $w_{\bar\rho}=\bar w$, $w_{1-\rho}=w^{-1}$,
$w_{1-\bar\rho}=\bar w^{-1}$.

\begin{theorem}[radial positivity holds even off the line]\label{thm:abel}
For every quartet $\mathcal O(\rho)$ --- including quartets strictly off the critical line --- the
radial Abel germ of its Li contribution is strictly positive.
\end{theorem}

\begin{nogo}[radial positivity is compatible with exponentially negative
coefficients]\label{ng:abel}
No radial average of fixed order transfers the sign to individual coefficients or to triangular
means, so no Tauberian argument local in the radial variable can supply the lower bound in
\eqref{eq:abel-coord}.
\end{nogo}

\begin{witness}[Abel-positive, coefficient-negative]\label{wit:abel}
For $w=2i$ the exact rational values of the quartet's fourth Li coefficient and its eighth
triangular mean are
\[
   Q_4=-\frac{225}{8},\qquad F_8=-\frac{10885}{128},
\]
both strictly negative, while the radial Abel germ of the same quartet is strictly positive.
\end{witness}

\noindent The structural reason is not subtle and is worth spelling out, because it recurs in
\S\ref{sec:fejer}. Radial summation at $z=\varrho\to1^-$ probes the function near a single
boundary point along a single direction; the coefficients probe all frequencies. An off-line
quartet has $|w|\ne1$, so its coefficient sequence contains $\cosh(na)$, which grows; the
radial germ, being an average with rapidly decaying weights, does not see the growth. This is the
same phenomenon as the Bombieri--Lagarias exponential blowup \cite{BombieriLagarias1999}, and it
means that any averaging device whose kernel decays must be discarded before it is tried.

We record, for completeness, that pointwise Fej\'er positivity on the circle had already been
shown not to survive a reciprocal off-line quartet;
Theorem~\ref{thm:abel} and Witness~\ref{wit:abel} add the exact cumulative and triangular sums and
the radial statement, which is the part that identifies \emph{why} a radial Tauberian argument
cannot be the missing mechanism.

\subsection{Real order and monotonicity of the correlated shift}\label{sec:shift}

\subsubsection*{Admissible inputs}
The absolutely convergent series of $A_u$ in $\Re s>1+u$; the Laguerre generating identity; Abel
summation; the real order of $F(s)=(s-1)\zeta(s)$ on $(1,\infty)$. The mandatory falsifier is the
symmetric product $X_a(s)=\xi(s+a)\xi(s-a)$.

Rather than take the infinitesimal jet, this mechanism keeps the regulator and the Jordan shift
finite and correlated,
\begin{equation}\label{eq:shift}
   s_\eps(z)=1+\eps+\frac{z}{1-z},\qquad u=c\eps,\quad0<c<1 ,
\end{equation}
and retains $H_{c\eps}$ exactly until after the coefficient is extracted. It produces an exact
finite recurrence and an exact renewal equation --- both new and both sign-free --- and then
returns to the starting point.

\begin{nogo}[the real order of the cocycle does not control its Cayley
coefficients]\label{ng:shift}
Removing the regulator returns exactly the cost functional
$C_n=\lambda_n^{\mathrm{prime}}-\lambda_{n+1}^{\mathrm{prime}}$. Real monotonicity of
$(s-1)\zeta(s)$ on the ray is compatible with $C_6>0$; the third difference in $c$ has the sign
opposite to the one a Hausdorff-type argument would need; and uniqueness of Dirichlet series
forbids annihilating the explicit channel identically with a nontrivial finite system of constant
weights.
\end{nogo}

\noindent The lesson of this entry is about the direction of the implication. One has a positive
object ($J_{c\eps}\ge0$) and a contraction property on a ray ($H_{c\eps}<1$ there), and the hoped
conclusion was that $C_n$ has the required sign. Both hypotheses hold; the conclusion does not
follow, because the ray is a one-dimensional probe of a two-dimensional problem, exactly as in
\S\ref{sec:abel}.

\subsection{The branch-free three-channel expansion}\label{sec:three-channel}

\subsubsection*{Admissible inputs}
Nonnegativity of the convolution square $J_u*J_u$; the positive Beta representation of the Gamma
ratio; the functional equation; the exact polar expansion. Falsifier: the shifted product
$X_a$, and Pringsheim's theorem on power series with eventually one-signed coefficients.

Put $F(s)=(s-1)\zeta(s)$, $H_u=F(s-u)/F(s)$,
$K_u(s)=\pi^{u/2}\Gamma(1+(s-u)/2)/\Gamma(1+s/2)$, and
$Y(s)=\xi(s)^2/\bigl(s\pi^{-s/2}\Gamma(s/2)\bigr)$, so that
\begin{equation}\label{eq:Ycoc}
   \mathcal S_u(s)=\frac{Y(s-u)}{Y(s)}=H_u(s)^2K_u(s),
   \qquad
   \log\frac{Y((1-z)^{-1})}{Y(1)}=\sum_{n\ge1}\frac{D_n}{n}z^n ,
\end{equation}
with $D_n=D_n^{[2]}=2\lambda_n-A_n$. The three channels are: the arithmetic square, with
nonnegative coefficients $b_u=J_u*J_u$; the Gamma factor, with the positive Beta measure
\begin{equation}\label{eq:beta-measure}
   d\nu_{\eps,u}(v)=\frac{\pi^{u/2}}{\Gamma(u/2)}
   v^{(1+\eps-u)/2}(1-v)^{u/2-1}\,dv ,
\end{equation}
in which the factor $v^{z/[2(1-z)]}$ shifts the Laguerre argument by $-\tfrac12\log v$; and the
polar square, whose three terms are $1$, $-2u/(s-1)$, $u^2/(s-1)^2$.

\begin{theorem}[exact three-channel expansion]\label{thm:three}
With $d\omega_0=\delta_0(dy)$, $d\omega_1=e^{-\eps y}dy$, $d\omega_2=ye^{-\eps y}dy$ and
\[
   T_{n,j}=\sum_{m\ge1}\frac{b_u(m)}{m^{1+\eps}}
   \int_0^1\!\!\int_0^\infty
   L_n\bigl(\log m+y-\tfrac12\log v\bigr)\,d\omega_j(y)\,d\nu_{\eps,u}(v),
\]
one has, for every $n\ge0$,
\begin{equation}\label{eq:three}
   [z^n]\,\frac{\mathcal S_u(s_\eps(z))}{1-z}=T_{n,0}-2u\,T_{n,1}+u^2\,T_{n,2},
\end{equation}
and, writing
$\bigl(\mathcal S_{c\eps}(s_\eps(z))-1\bigr)\big/\bigl(c\eps(1-z)\bigr)
=\sum_ng_{n,\eps,c}z^n$,
\begin{equation}\label{eq:three-limit}
   g_{0,\eps,c}\to-D_1,\qquad
   g_{n,\eps,c}\to-\Delta D_n\quad(n\ge1),
\end{equation}
the parameter $c$ disappearing in the limit.
\end{theorem}

\begin{proof}
Use $e^{-xz/(1-z)}/(1-z)=\sum_nL_n(x)z^n$ on the arithmetic and Beta channels to get $T_{n,0}$,
and $\bigl(\eps+z/(1-z)\bigr)^{-1}=\int_0^\infty e^{-\eps y}e^{-yz/(1-z)}dy$, with the squared
denominator inserting the factor $y$, to shift $L_n(x)$ to $L_n(x+y)$. Combining the three terms
of the polar square gives \eqref{eq:three}; all series and integrals converge absolutely for
$u<\eps$ and $z$ near the origin. Since $\bigl(\mathcal S_u(s)-1\bigr)/u\to-Y'(s)/Y(s)$, the
limits \eqref{eq:three-limit} follow from \eqref{eq:Ycoc}.
\end{proof}

\begin{nogo}[Jordan positivity plus a Beta measure plus the functional equation do not force the
Cayley sign]\label{ng:three}
Each $T_{n,j}$ is built from positive measures, but $L_n$ has sign and the combination
$1,-2u,u^2$ has sign; by Theorem~\ref{thm:polar}(i) the polar square is not a positive Laplace
mixture for any $c>0$; separating the channels creates divergences of order $\eps^{-1}$ by
\eqref{eq:m1-divergence}; and the shifted falsifier shows that no structural continuation of the
sign is available.
\end{nogo}

\noindent The falsifier deserves its own statement, because it isolates the exact analytic
boundary of the mechanism. Let $X_a(s)=\xi(s+a)\xi(s-a)$ with $0<a<\half$, so
$X_a(1-s)=X_a(s)$ and $X_a$ has zeros at $\half\pm a+i\gam$. Put $Y_a(s)=Y(s+a)Y(s-a)$ and
$\mathcal S_u^{(a)}(s)=Y_a(s-u)/Y_a(s)$. If $0<\eps<a$, a zero of $Y_a$ with real part
$\half+a$ is carried by the Cayley map $z=(s-1-\eps)/(s-\eps)$ into $|z|<1$, and for $u$ outside a
discrete set the numerator does not cancel it, so the Taylor series of
$-\bigl(\mathcal S_u^{(a)}(s_\eps(z))-1\bigr)/\bigl(u(1-z)\bigr)$ has radius $R<1$. Were its
coefficients eventually of one sign, Pringsheim's theorem --- applied after subtracting a
polynomial and dividing by a power of $z$ --- would force a singularity at the positive real point
$z=R$. There is none: for $0<z<1$ one has $s_\eps(z)>1$ and the shifted completed factors are
real, positive and nonvanishing, while the real singularity at $z=1$ is useless because $R<1$.
Applying the argument to the negative of the same series excludes eventual nonpositivity too.
Since the coefficients are real, infinitely many are positive and infinitely many negative.

The scope of this falsifier must be stated exactly, because it is easy to over-read. The positive
Dirichlet series of the shifted factor converges only in a right half-plane, not on the whole
Cayley disc. The argument therefore refutes a continuation of the sign based \emph{only} on
``positive Jordan coefficients plus a positive Beta measure plus the functional equation''. It
does not refute a new arithmetic theorem that uses quantitative information about the real weights
inside the critical strip. And requiring holomorphy on the full disc for a fixed $\eps$ excludes
only zeros with $\Re\rho>\half+\eps$; requiring it for \emph{every} $\eps\downarrow0$, with
generic $u$, forces the Hypothesis --- which brings us to the last analytic mechanism.

\subsection{Fej\'er--Carath\'eodory: the shortest route, and why its hypothesis is the
conclusion}\label{sec:fejer}

\subsubsection*{Admissible inputs}
The Herglotz representation of a function with nonnegative real part; nonnegativity of the Fej\'er
kernel; the Cayley bijection between the disc and the half-plane; Pringsheim's theorem;
Schur/Carath\'eodory bounds.

Fix $u>0$ and put
\begin{equation}\label{eq:Phi}
   \Phi(z)=\mathcal S_u^{[3]}\bigl(s_\eps(z)\bigr)=\sum_{k\ge0}a_kz^k,
   \qquad
   Q(z)=\frac{1-\Phi(z)}{u(1-z)^2}=\sum_{m\ge0}q_mz^m .
\end{equation}
Multiplying series, with no analytic hypothesis whatsoever,
\begin{equation}\label{eq:fejer-coeff}
   q_m=\frac1u\Bigl[(m+1)-\sum_{k=0}^m(m-k+1)a_k\Bigr]
   =\frac{m+1}{u}\bigl(1-\sigma_m(\Phi;1)\bigr),
   \qquad
   \sigma_m(\Phi;1)=\sum_{k=0}^m\Bigl(1-\frac{k}{m+1}\Bigr)a_k ,
\end{equation}
and along $u=c\eps$ one has $q_{n-1}\to3\lambda_n-A_n$.

\begin{theorem}[Fej\'er--Carath\'eodory]\label{thm:fejer}
Suppose $\Phi$ is holomorphic in the unit disc $\Dd$, has real Taylor coefficients, and satisfies
$\Re\Phi(z)\le1$ on $\Dd$. Then $q_m\ge0$ for every $m\ge0$.
\end{theorem}

\begin{proof}
$g=1-\Phi$ has nonnegative real part, so admits a Herglotz representation with a positive measure
$\mu$ on $\Tt$, normalized by $\mu(\Tt)=g(0)=1-a_0$; since all coefficients are real the
imaginary constant may be taken to vanish. Then \eqref{eq:fejer-coeff} becomes
\[
   q_m=\frac{m+1}{u}\int_{-\pi}^{\pi}F_m(t)\,d\mu(t),
   \qquad
   F_m(t)=1+2\sum_{k=1}^m\Bigl(1-\frac{k}{m+1}\Bigr)\cos(kt)
   =\frac{1}{m+1}\Bigl|\sum_{j=0}^me^{ijt}\Bigr|^2\ \ge\ 0 ,
\]
the Fej\'er kernel \cite{Fejer}. The stronger hypothesis $|\Phi|\le1$ suffices as well.
\end{proof}

This is a correct and complete mechanism: it converts one classical inequality into every sign of
the cubic margin at once. The difficulty is entirely in its hypothesis.

\begin{theorem}[the required holomorphy is equivalent to the Hypothesis]
\label{thm:fejer-rh}
For $u_j=c\eps_j$ with $0<c<1$ and $\eps_j\downarrow0$, put
$\Phi_j(z)=Y_3(s_{\eps_j}(z)-u_j)/Y_3(s_{\eps_j}(z))$. Then the following statement is equivalent
to the Riemann Hypothesis:
\begin{quote}
there exists a sequence $\eps_j\downarrow0$ for which every $\Phi_j$ is holomorphic on $\Dd$.
\end{quote}
\end{theorem}

\begin{proof}
The Cayley map is a bijection
$s_\eps(\Dd)=\{s:\Re s>\half+\eps\}$ with inverse $z=(s-1-\eps)/(s-\eps)$. Under the Hypothesis
the denominator has no zeros in that half-plane, so each $\Phi_j$ is holomorphic. Conversely
suppose $\xi(\rho)=0$ with $\Re\rho>\half$. For all large $j$, $\rho$ lies in the half-plane, and
since $G$ is regular and nonvanishing on the critical strip, removability of the singularity in
$\Phi_j$ forces $\xi(\rho-c\eps_j)=0$. A subsequence then yields infinitely many distinct zeros
accumulating at $\rho$, contradicting the isolation of zeros of a nonzero entire function. Hence
there are no zeros to the right of the line, and the functional equation with conjugation
excludes those to the left.
\end{proof}

\begin{nogo}[the analytic hypothesis contains the conclusion]\label{ng:fejer}
Applying Theorem~\ref{thm:fejer} to a sequence that removes the regulator introduces the desired
conclusion in the holomorphy hypothesis, before the inequality $\Re\Phi\le1$ is ever used. And the
coefficientwise weakening does not escape: if, at fixed $\eps>0$, $q_m\ge0$ for all large $m$,
then $Q$ has radius of convergence at least one --- otherwise, after removing a finite prefix,
Pringsheim's theorem forces a singularity at the positive real point $z=R<1$, whereas the
meromorphic formula for $\Phi$ is holomorphic near every point of $(0,1)$ and continues the germ
there. Hence eventual nonnegativity at fixed $\eps$ holomorphically extends $\Phi$ to $\Dd$, and
proving it along a sequence $\eps_j\downarrow0$ again implies the Hypothesis by
Theorem~\ref{thm:fejer-rh}.
\end{nogo}

\noindent Finally, the retreat to a smaller, safe disc fails quantitatively. The Euler--Dirichlet
series is guaranteed only in $\Re s>1+u$; with $\delta=\eps-u$ the largest centred disc that the
Cayley map places inside that half-plane has radius
\begin{equation}\label{eq:safe-radius}
   R_\eps=\frac{\eps-u}{1-\eps+u}\sim(1-c)\eps\longrightarrow0 ,
\end{equation}
and Fej\'er on such a disc controls only means damped by $R_\eps^k$. The obstruction to
undamping is exact.

\begin{witness}[no transfer from a smaller disc]\label{wit:fejer}
Let $0<R<1$, choose $N$ with $R^{-N}>N+1$, and take
\begin{equation}\label{eq:fejer-witness}
   \Phi(z)=a+Mz^N,\qquad (N+1)(1-a)<M\le(1-a)R^{-N},\quad 0\le a<1 .
\end{equation}
Then $|\Phi(z)|\le1$ for $|z|\le R$, while
$q_N=\bigl[(N+1)(1-a)-M\bigr]/u<0$. Concretely, $R=\tfrac12$ and $\Phi(z)=6z^4$ have local
supremum norm $\tfrac38$ on $|z|\le R$ and $q_4=-1$.
\end{witness}

\noindent So the information on any strictly smaller disc does not suffice without a new
inequality special to the arithmetic cocycle. Theorem~\ref{thm:fejer} remains available as a
mechanism; what is missing is a way to establish $\Re\Phi\le1$ on a disc large enough to matter
without assuming what one wants to prove. That statement --- a signed inequality acting directly
on the limit $q_{n-1}\to3\lambda_n-A_n$, requiring neither global holomorphy nor eventual sign at
fixed regulator --- is the open problem this section leaves behind.

%% ======================================================================
\section{Closed mechanisms V: algebraic squares are not metric squares}\label{sec:mech-squares}
%% ======================================================================

The two mechanisms of this section are the ones that come closest to working, and their failure is
the most precisely diagnosable of the sixteen. Both produce an exact factorization of the
target as a square. In both cases the square is an \emph{algebraic} square $C^2$ --- a composition
of commuting operators --- rather than a \emph{metric} square $C^{*}C$, and therefore defines no
norm and carries no sign. And in the quartic case the genuinely quadratic part is destroyed by the
very limit that extracts the coefficient. Together with
Corollary~\ref{cor:shortfall} this is the sharpest form of the paper's thesis: even when a square
exists, it certifies $\ge0$, and $\ge0$ is not what is needed.

\subsection{The local tower square}\label{sec:tower}

\subsubsection*{Admissible inputs}
The exact local coefficient formula on a prime tower; discrete summation by parts; commuting
translation and derivation operators; the positive Beta measure of the Gamma channel; the exact
polar expansion. No use of a norm of $\Cc$, no tower-by-tower sign, no per-background sign.

Let $N=-\partial_x$ and, for $\eps>0$, $\ell>0$,
\begin{equation}\label{eq:ops}
   \Aa_\eps f(x)=\int_0^\infty e^{-r}f(x+r/\eps)\,dr,
   \qquad
   \Uu_\ell f(x)=\int_0^\ell f(x+t)\,dt .
\end{equation}
Since $(\eps-\partial_x)^{-1}=\eps^{-1}\Aa_\eps$, the polar factor of a single cocycle acts on the
Laguerre variable as $\Pp_{\eps,c}=I-c\Aa_\eps$ with $u=c\eps$, and integration by parts in
\eqref{eq:ops} gives $\Aa_\eps N=\eps(I-\Aa_\eps)$, whence the equivalent form
\begin{equation}\label{eq:P-form}
   \Pp_{\eps,c}=(1-c)I+\frac{c}{\eps}\Aa_\eps N .
\end{equation}
Similarly, with $\rho=p^{-1-\eps}$ and $(E_\ell f)(x)=f(x+\ell)$, one has
$\Uu_\ell N=I-E_\ell$ and therefore
\begin{equation}\label{eq:B-form}
   I-\rho E_\ell=(1-\rho)I+\rho\,\Uu_\ell N .
\end{equation}

\begin{theorem}[exact local differences on a prime tower]\label{thm:tower-diff}
Let $\ell=\log p$, $Q=p^u$, and let $\nabla$ be the backward difference with the convention
$a_k=0$ for $k<0$. Then
\begin{equation}\label{eq:tower-2}
   b_u(p^k)=\nabla^2\bigl((k+1)Q^k\bigr),\qquad k\ge0,
\end{equation}
and more generally, for every integer $r\ge1$,
\begin{equation}\label{eq:tower-r}
   \bigl(J_u^{*r}\bigr)(p^k)=\nabla^r\Bigl[\binom{k+r-1}{r-1}Q^k\Bigr],
\end{equation}
both sides having generating function $\bigl(\tfrac{1-x}{1-Qx}\bigr)^r$. In particular the fourth
convolution power is $d_u(p^k)=\nabla^4\bigl[\binom{k+3}{3}Q^k\bigr]$. Consequently, for every
polynomial $f$,
\begin{equation}\label{eq:tower-sum}
   \sum_{k\ge0}b_u(p^k)\rho^kf(x+k\ell)
   =\sum_{k\ge0}(k+1)(Q\rho)^k\,(I-\rho E_\ell)^2f(x+k\ell),
\end{equation}
and analogously with $(I-\rho E_\ell)^r$ and $\binom{k+r-1}{r-1}$; both series converge absolutely
because $Q\rho=p^{-1-(1-c)\eps}<1$.
\end{theorem}

The operators \eqref{eq:P-form} and \eqref{eq:B-form} commute on polynomials, so setting
\begin{equation}\label{eq:C-op}
   \Cc_{p,\eps,c}
   =\bigl((1-\rho)I+\rho\,\Uu_\ell N\bigr)
    \Bigl((1-c)I+\tfrac{c}{\eps}\Aa_\eps N\Bigr)
\end{equation}
one obtains the factorization that names the mechanism.

\begin{theorem}[the coupled local square]\label{thm:tower-square}
A complete tower $p^k$, with all three polar channels retained,
satisfies
\begin{equation}\label{eq:tower-square}
   \sum_{k\ge0}b_u(p^k)\rho^k\,\Pp_{\eps,c}^2f(x+k\ell)
   =\sum_{k\ge0}(k+1)(Q\rho)^k\,\Cc_{p,\eps,c}^2f(x+k\ell),
\end{equation}
and, expanding $\Cc=a_0+a_1N+a_2N^2$ with
\[
   a_0=(1-\rho)(1-c)I,\qquad
   a_1=(1-\rho)\tfrac{c}{\eps}\Aa_\eps+\rho(1-c)\Uu_\ell,\qquad
   a_2=\tfrac{\rho c}{\eps}\Uu_\ell\Aa_\eps ,
\]
all of which are positive combinations of translations and mutually commute, the relation
$N^jL_n=L^{(j)}_{n-j}$ gives
\begin{equation}\label{eq:C2L}
   \Cc^2L_n=a_0^2L_n+2a_0a_1L^{(1)}_{n-1}+\bigl(a_1^2+2a_0a_2\bigr)L^{(2)}_{n-2}
   +2a_1a_2L^{(3)}_{n-3}+a_2^2L^{(4)}_{n-4}.
\end{equation}
\end{theorem}

The signed measure of the three-channel expansion has vanished inside a square, and all
operatorial weights in \eqref{eq:C2L} are positive. This is genuine progress in organization, and
it is where the mechanism stops.

\begin{nogo}[$\Cc^2$ is not $\Cc^{*}\Cc$, and the real blocks change sign]\label{ng:tower}
Equation \eqref{eq:tower-square} is an algebraic composition, not a Gram form: it defines no norm,
and every raised Laguerre polynomial in \eqref{eq:C2L} still oscillates. Worse, the sign failure is
not formal --- the blocks assembled from the \emph{real} arithmetic weights take both signs.
\end{nogo}

\noindent Consider a tower $p^k$ against a background $m$ coprime to $p$. After integrating the
Beta measure of the Gamma channel, the degree-one block equals a positive prefactor
$b_u(m)m^{-1-\eps}$ times
\begin{equation}\label{eq:block-1}
   K_0M_p\Bigl[(1-c)^2\bigl(1-\log m-\bar\eta-\ell\mu_p\bigr)+\frac{2c(1-c)}{\eps}\Bigr],
\end{equation}
where $K_0=K_u(1+\eps)>0$, $M_p=\bigl(\tfrac{1-\rho}{1-Q\rho}\bigr)^2=\sum_kb_u(p^k)\rho^k>0$,
$\mu_p=\tfrac{2\rho(Q-1)}{(1-\rho)(1-Q\rho)}>0$ is the exact mean of $k$ under those normalized
weights, and
$\bar\eta=\tfrac12\bigl[\psi\bigl(\tfrac{3+\eps}{2}\bigr)-\psi\bigl(\tfrac{3+\eps-u}{2}\bigr)\bigr]>0$
is the mean of $-\tfrac12\log v$ under the normalized Beta measure. The block changes sign exactly
when
\begin{equation}\label{eq:threshold-m}
   \log m=1-\bar\eta-\ell\mu_p+\frac{2c}{(1-c)\eps} ,
\end{equation}
so for fixed $p$ and $c\in(0,1)$ there is $\eps_0(p,c)>0$ such that for
$0<\eps<\eps_0(p,c)$ the block at $m=1$ is on the positive side while every sufficiently large
coprime background is on the negative side. Both signs occur with the real weights $b_u$. More
generally, the leading coefficient in $\log m$ of the degree-$n$ block is
\begin{equation}\label{eq:leading-n}
   \frac{K_0M_p(1-c)^2}{n!}(-1)^n ,
\end{equation}
so for every \emph{even} $n$, including every even $n\ge150$, there are infinitely many coprime
backgrounds whose canonical block is positive. Hence the global sign cannot be closed by imposing
a sign on each canonical block. The quartic analogue is identical in form: with
$M_{p,4}=\bigl(\tfrac{1-\rho}{1-Q\rho}\bigr)^4$ and
$\mu_{p,4}=\tfrac{4\rho(Q-1)}{(1-\rho)(1-Q\rho)}$, the degree-one block is a positive prefactor
times
\begin{equation}\label{eq:block-4}
   (1-c)^4\bigl(1-\log m-3\bar\eta-\ell\mu_{p,4}\bigr)+\frac{4c(1-c)^3}{\eps},
\end{equation}
changing sign at $\log m=1-3\bar\eta-\ell\mu_{p,4}+\tfrac{4c}{(1-c)\eps}$.

\begin{remark}[what survives: the global multiplicative sum]\label{rmk:tower-alive}
This no-go kills the strategy ``sign per tower, then sum''. It does not refute the sign of the
complete sum. With the exact global identity
$g_{n,\eps,c}\to-\Delta D_n$ of Theorem~\ref{thm:three}, the statement
\begin{equation}\label{eq:global-gate}
   \exists\,c\in(0,1)\ \forall n\ge149:\qquad
   \limsup_{\eps\downarrow0}g_{n,\eps,c}\le0
\end{equation}
is pointwise equivalent to $\Delta D_n\ge0$ (the limit exists), and together with the finite
certificates would give the strong margin, hence $(\star)$. A proof of \eqref{eq:global-gate}
may not use a sign per tower, a sign per background, or a norm of $\Cc$: it must retain the
multiplicative interactions between distinct primes. That is the narrowest live point of this
mechanism, and it is one of the two successors we return to in \S\ref{sec:successor}.
\end{remark}

\begin{remark}[numerical diagnostic, not evidence in any proof]\label{rmk:diag-global}
A Cauchy/Borwein extraction in double precision found $g_{n,\eps,c}<0$ up to $n=1200$ for
$(\eps,c)\in\{(10^{-3},\tfrac12),(10^{-2},\tfrac12),(10^{-1},\tfrac12),(1,\tfrac12),
(10^{-1},\tfrac9{10})\}$. This is not a certified interval and enters no proof. Its only use is
directional: no counterexample to the global sum appeared, even though \eqref{eq:block-1} and
\eqref{eq:leading-n} prove rigorously that its local summands take both signs.
\end{remark}

\subsection{The fourth power as a sum of squares}\label{sec:sos}

\subsubsection*{Admissible inputs}
The exact fourth difference \eqref{eq:tower-r} with $r=4$; the formal factorization of the quartic
cocycle; the infinite divisibility of the Euler and Gamma channels for real exponents; symmetry
and the functional equation. Falsifier: an off-line quartet inside the target range.

The near-quartic margin $r_*=2002/501$ of \eqref{eq:rstar} sits only $\delta=2/501$ below the
integer four, and at the integer the cocycle admits a formal square root:
\begin{equation}\label{eq:sos-root}
   B_u(s)=H_u(s)^2K_u(s)^{3/2},
   \qquad
   \mathcal S_u^{[4]}(s)=H_u(s)^4K_u(s)^3=B_u(s)^2 .
\end{equation}
This is the last, and most tempting, source of positivity in the whole family. It fails in three
independent ways.

\begin{nogo}[the local fourth power is algebraic, not a norm]\label{ng:sos-local}
By \eqref{eq:tower-r} with $r=4$, the tower sum is
\begin{equation}\label{eq:sos-tower}
   \sum_{k\ge0}d_u(p^k)\rho^k\Pp^4f(x+k\ell)
   =\sum_{k\ge0}\binom{k+3}{3}(Q\rho)^k\,\Cc^4f(x+k\ell),
\end{equation}
and although $\Cc^4=(\Cc^2)^2$, this is a composition, not $\Cc^{*}\Cc$. It defines no norm; and by
Witness~\ref{wit:pf2-quartic} the fourth convolution power is not even $\mathrm{PF}_2$, while
\eqref{eq:block-4} shows that its complete real blocks take both signs.
\end{nogo}

\begin{nogo}[the square disappears in the extraction limit]\label{ng:sos-limit}
The limit that extracts $4\lambda_n-A_n$ is linear, and the genuinely quadratic part is lost before
the coefficient is reached.
\end{nogo}

\noindent Indeed, from \eqref{eq:sos-root},
\begin{equation}\label{eq:sos-lin}
   \frac{1-\mathcal S_u^{[4]}}{u}=\frac{1-B_u}{u}\bigl(1+B_u\bigr),
   \qquad\text{whereas}\qquad
   \frac{(1-B_u)^2}{u}\longrightarrow0
\end{equation}
as $u\downarrow0$, since $B_u\to1$. So the extraction sees the first factor only; the square is
formally present and analytically absent. The writing \eqref{eq:sos-root} is not a sum-of-squares
certificate.

\begin{nogo}[a symmetry-only sum of squares would falsify an off-line divisor]\label{ng:sos-offline}
Any sum-of-squares argument using only conjugation, the functional equation and the fourth exponent
would apply verbatim to an artificial divisor with off-line zeros, whose quartic contribution is
strictly negative inside the target range.
\end{nogo}

\begin{witness}[off-line quartet at the first target index]\label{wit:sos-offline}
Take
\begin{equation}\label{eq:sos-w}
   w=\tfrac12e^{74\pi i/75},\qquad \rho=\frac{1}{1-w},
\end{equation}
so that $\half<\Re\rho<1$ and $\Im\rho\ne0$. For the associated symmetric quartet, with
$a=\log2$ and $\vartheta=74\pi/75$, the contribution to $4\lambda_n$ is
\begin{equation}\label{eq:sos-16}
   16\bigl(1-\cosh(na)\cos(n\vartheta)\bigr).
\end{equation}
At $n=150$ one has $n\vartheta=148\pi$, so $\cos(n\vartheta)=1$ and \eqref{eq:sos-16} is strictly
negative; repeating the quartet with multiplicity defeats any fixed archimedean term. The
phenomenon appears already at $n=2$: for $w=\tfrac12e^{5\pi i/6}$,
\begin{equation}\label{eq:sos-n2}
   \lambda_2^{\mathrm{quartet}}
   =4-4\cosh(2\log2)\cos\tfrac{5\pi}{3}
   =4-4\cdot\tfrac{17}{8}\cdot\tfrac12=-\tfrac14 .
\end{equation}
\end{witness}

\noindent Finally, the alternative of a positive generator for the quartic cocycle fails at the
same point as everything else in \S\ref{sec:mech-measure}: before the first atom $\log2$, the
first variation of the completed inverse is
$f_{4,\eps}(x)=-4e^{-\eps x}+3\tfrac{e^{-(3+\eps)x}}{1-e^{-2x}}$, and at $x_0=\tfrac12\log2$,
\begin{equation}\label{eq:sos-halflog2}
   f_{4,\eps}(x_0)=e^{-\eps x_0}\Bigl(-4+\frac{3}{\sqrt2}\Bigr)<0 .
\end{equation}

We can now state, in the coordinates of the cocycle, what the \emph{floor-$1000$ quartic
transfer} would have to supply. By Theorem~\ref{thm:threshold} its convenient slack is
$A_n/1001$, and in cocycle language the correction that subtracts it is fractional:
\begin{equation}\label{eq:frac-corrector}
   \mathcal S_u^{[r_*]}
   =\mathcal S_u^{[4]}\Bigl(\frac{\xi(s-u)}{\xi(s)}\Bigr)^{\!-2/501} .
\end{equation}
The fractional corrector in \eqref{eq:frac-corrector} removes precisely the slack of
\eqref{eq:threshold} and reintroduces the zero and branch gate. This explains the failure of that
particular fixed-floor transfer. It is not a correction required by Li's criterion or by an
adapted cofinal cutoff.

%% ======================================================================
\section{Closed mechanisms VI: the sign of the tail from an absolute envelope}
\label{sec:mech-envelope}
%% ======================================================================

There remains one route that would not need any positivity at all: at a prescribed cutoff, remove
the residual tail slack by improving Corollary~\ref{cor:a0plus} from a two-sided bound
$|R_n(T_n)|\le A_n/[4(1+T_n)]$ to a favourable \emph{sign}, $R_n(T_n)\le0$. By \eqref{eq:CnT}
this would remove the tail from the requirement entirely. We show that this is unavailable from the
inputs currently in hand.

\subsubsection*{Admissible inputs}
The tail estimate of Theorem~\ref{thm:a0}; the Vinogradov--Korobov envelope $|E(u)|\le W(u)=C_{\mathrm{VK}}e^{u-\eta(u)}$ with
$\eta$ increasing on the final ray; eventual monotonicity of $e^u+E(u)$; the location of the
largest zero of $L^{(2)}_{n-1}$. Nothing about the support of $d\psi$ on prime powers.

Let $x_{*,n}$ be the largest zero of $L^{(2)}_{n-1}$ and $s_n=(-1)^{n-1}$, so that beyond
$x_{*,n}$ the kernel has the constant sign $K_n(u)=s_ne^{-u}|L^{(2)}_{n-1}(u)|$.

\begin{theorem}[both tail signs saturate the same envelope]\label{thm:tail}
For $\sigma\in\{-1,+1\}$ define, on the ray $u\ge T>x_{*,n}$, the discrepancy
$E_\sigma(u)=\sigma s_nW(u)$. Both choices satisfy the envelope $|E_\sigma|\le W$ exactly, and
\begin{equation}\label{eq:tail-signs}
   R_{n,\sigma}(T)=-\sigma\,C_{\mathrm{VK}}
   \int_T^\infty e^{-\eta(u)}\bigl|L^{(2)}_{n-1}(u)\bigr|\,du ,
\end{equation}
so $\sigma=-1$ gives $R_{n,\sigma}(T)>0$ and $\sigma=+1$ gives $R_{n,\sigma}(T)<0$, for every
finite cutoff. The witness is not degenerate: factoring through the positive real Laguerre zeros,
\begin{equation}\label{eq:tail-quant}
   \bigl|L^{(2)}_{n-1}(u)\bigr|=\frac{1}{(n-1)!}\prod_{j=1}^{n-1}(u-x_{j,n})
   \ \ge\ \frac{(T-x_{*,n})^{n-1}}{(n-1)!}\qquad(T\le u\le T+1),
\end{equation}
whence, if $\eta$ is increasing on that ray,
\begin{equation}\label{eq:tail-lower}
   \bigl|R_{n,\sigma}(T)\bigr|
   \ \ge\ \frac{C_{\mathrm{VK}}\,e^{-\eta(T+1)}(T-x_{*,n})^{n-1}}{(n-1)!} .
\end{equation}
\end{theorem}

\begin{theorem}[monotonicity does not decide the sign either]\label{thm:tail-mono}
Put $\phi_T(u)=1-e^{-(u-T)}$ and $\widetilde E_\sigma(u)=\sigma s_nW(u)\phi_T(u)$, so that the
initial datum $E(T)=0$ is shared. Then, restricted to $[T+1,T+2]$,
\begin{equation}\label{eq:tail-smooth}
   \bigl|\widetilde R_{n,\sigma}(T)\bigr|
   \ \ge\ \frac{C_{\mathrm{VK}}(1-e^{-1})e^{-\eta(T+2)}(T+1-x_{*,n})^{n-1}}{(n-1)!},
\end{equation}
with opposite signs for the two values of $\sigma$. Moreover, if $T$ is enlarged until
\begin{equation}\label{eq:tail-mono-cond}
   C_{\mathrm{VK}}e^{-\eta(u)}
   \Bigl(\phi_T(u)+\bigl|\phi_T'(u)-\eta'(u)\phi_T(u)\bigr|\Bigr)<1
   \qquad(u\ge T),
\end{equation}
then $\tfrac{d}{du}\bigl(e^u+\widetilde E_\sigma(u)\bigr)>0$ for \emph{both} signs. The usual
Vinogradov--Korobov functions satisfy \eqref{eq:tail-mono-cond} on every sufficiently distant ray,
and enlarging an admissible cutoff preserves the tail estimate.
\end{theorem}

\begin{nogo}[the envelope does not determine the sign of the tail]\label{ng:tail}
\begin{equation}\label{eq:tail-nogo}
   \text{absolute envelope}\ +\ \text{monotonicity on the final ray}
   \ \not\Longrightarrow\ R_n(T_n)\le0 .
\end{equation}
\end{nogo}

\noindent Nor is ``take the first $T$ with $R_n(T)\le0$'' an unconditional construction: within the
class \eqref{eq:tail-signs} that set may be empty, and verifying that it is nonempty for the
\emph{actual} arithmetic discrepancy already requires new arithmetic information. We stress the
scope: the witnesses of Theorems~\ref{thm:tail} and \ref{thm:tail-mono} are not supported on prime
powers and do not model the von Mangoldt weights. What they prove is that the class of inputs
currently available is insufficient --- and, by elimination, that any theorem about the sign of the
real tail must use the support on prime powers and the exact values $\Lambda(m)$ in an essential
way. That is the same conclusion as every other section above.

%% ======================================================================
\section{The three universal falsifiers}\label{sec:falsifiers}
%% ======================================================================

Sixteen mechanisms were closed, but only three falsifiers were needed. Understanding the division
of labour among them is what turns the catalogue from a list into a classification, because each
falsifier refutes a different \emph{kind} of hypothesis, and together they exhaust the ways a
semidefiniteness certificate can be assembled.

We first state the methodological rule that governs their use, because it was arrived at painfully
and it is the reason the conclusions reached here are stable.

\begin{convention}[the pre-registration rule]\label{conv:prereg}
Every mechanism opens with (i) its statement, (ii) an explicit enumeration of the axioms and inputs
it actually uses, and (iii) a falsifier \emph{constructed so as to satisfy exactly those axioms}.
A falsifier that does not satisfy the axioms in use validates nothing: if a mechanism survives a
model it was never meant to apply to, that is not evidence of anything, and the correct response is
to rebuild the model, not to declare victory.
\end{convention}

\begin{remark}[the logical limit of the rule]\label{rmk:prereg-limit}
The rule has a boundary which must be stated honestly. The deepest input level available fixes
$q=p^k$ and $\ell=\log p$ literally; that is, it fixes the arithmetic of $\zeta$. An ``off-line
divisor'' satisfying \emph{that} axiom would be $\zeta$ itself with an off-line zero, whose
existence is exactly the question at issue. So an off-line falsifier cannot honestly be demanded of
a step that genuinely uses the real weights without presupposing the negation of the Hypothesis.
The valid form of the rule is therefore: whenever an estimate \emph{relaxes} the arithmetic level,
it must pass the falsifier appropriate to the axioms it actually retains. An argument that uses a
special quantitative property of the real weights is not refuted by any falsifier in this
section --- but it must \emph{declare and prove} that property. Not one of the sixteen mechanisms
did.
\end{remark}

\subsection{Falsifier I: the off-line quartet}\label{sec:fals-quartet}

The first falsifier refutes every hypothesis built from symmetry alone: conjugation, the functional
equation, the multiset of zero heights, lobe occupancy, radial averages. In Cayley coordinates
$w=e^{-a+i\vartheta}$, a symmetric quartet $\{\rho,\bar\rho,1-\rho,1-\bar\rho\}$ contributes to
$m\lambda_n$ the quantity
\begin{equation}\label{eq:quartet-general}
   4m\bigl(1-\cosh(na)\cos(n\vartheta)\bigr),
\end{equation}
which for $a\ne0$ is unbounded below along the subsequence of $n$ with $\cos(n\vartheta)$ near
$1$, and grows exponentially in $n$ --- this is Bombieri--Lagarias
\cite{BombieriLagarias1999}, Theorem~1(c). Three instantiations were used above:
\begin{itemize}[leftmargin=*]
\item \S\ref{sec:m3}: $\vartheta=\pi/50$, $r=\tfrac{1+\cos\vartheta}{2}$, with $100\mid n$ and
      $n\ge300$, placing the stationary point uniformly in the bulk;
\item \S\ref{sec:sos}: $w=\tfrac12e^{74\pi i/75}$, giving strict negativity already at the first
      target index $n=150$;
\item \S\ref{sec:three-channel}: the entire family $X_a(s)=\xi(s+a)\xi(s-a)$, which retains the
      whole positive Jordan hierarchy and the functional equation while having a full line of
      off-line zeros.
\end{itemize}
There is also a purely archimedean instance which is useful because it involves no off-line zeros
at all: $X(s)=(s-\half)^2+\tfrac14$ has all its zeros on the critical line, yet
$\Delta\lambda_2[X]=-2$. This kills any derivation of a first-difference margin from the Hypothesis
itself or from per-zero squares, and it shows that even the on-line world does not supply the
monotonicity that several mechanisms wanted.

\begin{remark}\label{rmk:quartet-scope}
What the quartet family cannot do is refute an inequality that uses the Euler product. Every
statement in \S\S\ref{sec:m3}, \ref{sec:sos}, \ref{sec:three-channel} that invokes it is therefore
phrased as: \emph{no argument using only these symmetry inputs can succeed}. That is exactly the
scope claimed and no more.
\end{remark}

\subsection{Falsifier II: sign changes of the genuinely arithmetic blocks}\label{sec:fals-arith}

The second falsifier is the one that cannot be dismissed as an artefact of an abstract model,
because it is computed from the real primes in exact rational arithmetic. It refutes hypotheses of
the form ``the arithmetic block has a fixed sign'', which is how every local or blockwise
positivity argument must begin. Three instances were used:
\begin{itemize}[leftmargin=*]
\item the cumulative arithmetic correlator changes sign:
      $H(2969)<-21$ and $H(3167)>110$ for
      $H(x)=\sum_{q\le x}\Lambda(q)(q-1)-\psi(x)^2/2$ \eqref{eq:H-witness};
\item the $\mathrm{PF}_2$ minor on the fibre $3\cdot2^k$ equals $-\tfrac{1}{2916}$
      \eqref{eq:pf2-witness}, and the quartic minor equals $\tfrac{11-8\sqrt2}{8}<0$
      \eqref{eq:pf2-quartic};
\item the canonical block of the tower square changes sign as the coprime background varies,
      at the exact threshold \eqref{eq:threshold-m}, with leading coefficient
      \eqref{eq:leading-n} of sign $(-1)^n$ --- so for every even degree there are infinitely many
      positive blocks and infinitely many negative ones.
\end{itemize}
The first two use the primes $2$ and $3$ and the actual von Mangoldt values; the third uses the
actual multiplicative structure. These are not relaxations. They prove that positivity is not a
property of the arithmetic pieces at any of the natural granularities.

\subsection{Falsifier III: the negative density below the first prime atom}
\label{sec:fals-halflog2}

The third falsifier is the sharpest and the most structural, and it is the subject of the next
section. It refutes every hypothesis of the form ``the completed cocycle is the transform of a
positive measure'', including Bernstein, complete monotonicity, stochastic order,
$\mathrm{PF}_\infty$, decreasing variation, and generator positivity in the Stein--Mecke sense. It
does so by evaluating an exact density at the single rational point
\begin{equation}\label{eq:x0}
   x_0=\tfrac12\log2 .
\end{equation}
Its instances above are \eqref{eq:cm-witness}, \eqref{eq:stein-neg}, \eqref{eq:kappa-neg} and
\eqref{eq:sos-halflog2}. That one point should carry so much weight is not an accident, and the
reason is worth a section of its own.

%% ======================================================================
\section{The wall at $\tfrac12\log2$: arithmetic propagation in the Laplace variable}
\label{sec:halflog2}
%% ======================================================================

There is a structural obstruction which recurs across essentially every arithmetic approach to the
Hypothesis, and which we shall call the \emph{arithmetic-propagation} wall: information available
where the Euler product converges does not propagate into the region $\half<\Re s<1$ where the
Hypothesis lives. We now show that the
recurring witness at $x_0=\tfrac12\log2$ is precisely that wall, transported into the Laplace
variable of the completed cocycle, and that it is not a near miss but a robust structural
obstruction with a rational certificate.

\begin{theorem}[the completed first variation is negative below the first atom]
\label{thm:halflog2}
Fix $r>1$ and $\eps>0$, normalize the polar--Gamma product at $s_0=1+\eps$,
\[
   B_{r,u}(t)=\Bigl(\frac{P_u(s_0+t)}{P_u(s_0)}\Bigr)^{\!r}
              \Bigl(\frac{K_u(s_0+t)}{K_u(s_0)}\Bigr)^{\!r-1},
   \qquad P_u(s)=1-\frac{u}{s-1},
\]
and let $f_{r,\eps}$ be the continuous part of the inverse Laplace transform of
$\partial_uB_{r,u}\big|_{u=0}$. Then
\begin{equation}\label{eq:frq}
   f_{r,\eps}(x)=-r\,e^{-\eps x}+(r-1)\frac{e^{-(3+\eps)x}}{1-e^{-2x}},
\end{equation}
and at the point $x_0=\tfrac12\log2$, where $e^{-2x_0}=\tfrac12$ and hence
$e^{-3x_0}/(1-e^{-2x_0})=1/\sqrt2$,
\begin{equation}\label{eq:frq-neg}
   \boxed{\ f_{r,\eps}(x_0)=e^{-\eps x_0}\Bigl(-r+\frac{r-1}{\sqrt2}\Bigr)\ <\ 0
   \qquad\text{for every }r>1.\ }
\end{equation}
The normalized arithmetic factor $A_u^r$ has positive mass at zero and all its remaining support in
$[\log2,\infty)$; since $0<x_0<\log2$, it cannot alter the density in a neighbourhood of $x_0$.
Hence, testing against a nonnegative bump supported in that neighbourhood, for every fixed $r>1$
and $\eps>0$ there is $u_0(r,\eps)>0$ such that for $0<u<u_0(r,\eps)$ the normalized completed
cocycle is not the Laplace transform of a positive measure.
\end{theorem}

The quantifiers also cover the correlated path $u=c\eps$, but this does not follow from a Taylor
expansion uniform in $u$, because the polar normalization contains $(1-c)^{-r}$. It is verified
directly. Let $\varphi\ge0$ be smooth, not identically zero, supported in a compact interval around
$x_0$ contained in $(0,\log2)$, and let $\mu_{r,c\eps,\eps}$ be the inverse Laplace transform of
the full normalized cocycle. Since there are no arithmetic atoms in that support, and since the
mass at zero of the Euler channel tends to $(1-c)^r$ and cancels the polar factor $(1-c)^{-r}$, the
diagonal expansion gives
\begin{equation}\label{eq:diag-expansion}
   \langle\varphi,\mu_{r,c\eps,\eps}\rangle
   =c\eps\int_0^\infty\varphi(x)
   \Bigl[-r+(r-1)\frac{e^{-3x}}{1-e^{-2x}}\Bigr]dx+o(\eps),
\end{equation}
with the $o(\eps)$ understood for $c$, $r$ and the compact test $\varphi$ fixed. The bracket in
\eqref{eq:diag-expansion} equals $-r+(r-1)/\sqrt2<0$ at $x_0$ and remains negative on a small
enough support, so the left-hand side is negative for all sufficiently small $\eps>0$ along
$u=c\eps$. This is the regime that extracts the coefficient, so the defect is present exactly where
it matters.

\begin{obstruction}[arithmetic propagation, localized in the Laplace variable]\label{ob:halflog2}
Let $x_0=\tfrac12\log2$, strictly below the first prime atom $\log2$. In a neighbourhood of $x_0$
the arithmetic channel of the completed cocycle contributes nothing beyond its mass at the origin:
its remaining support is contained in $[\log2,\infty)$. Consequently the density there is produced
by the Gamma and polar channels alone, and for every real exponent $r>1$ that combination is
strictly negative, by \eqref{eq:frq-neg}. Any route that asks the completed cocycle to be a
positive measure, a Bernstein function, or a positive generator therefore asks the archimedean
factor to perform, in a region containing no primes, work that only the primes could do.
\end{obstruction}

\begin{proposition}[the wall is not a near miss]\label{prop:not-near-miss}
The defect at $x_0$ is
\begin{equation}\label{eq:defect}
   -r+\frac{r-1}{\sqrt2}=-\Bigl(1-\frac1{\sqrt2}\Bigr)r-\frac1{\sqrt2}
   \approx-0.2929\,r-0.7071 ,
\end{equation}
so it is bounded away from zero by a constant proportion of $r$, uniformly. In particular there is
no exponent $r>0$ for which the combination is nonnegative: that would require
$r(1-\sqrt2)\ge1$, impossible for $r>0$.
\end{proposition}

\begin{proof}
Immediate from \eqref{eq:frq-neg}; the last claim is $-r+(r-1)/\sqrt2\ge0\iff
r(1/\sqrt2-1)\ge1/\sqrt2\iff r(1-\sqrt2)\ge1$, and $1-\sqrt2<0$.
\end{proof}

Proposition~\ref{prop:not-near-miss} is what elevates the witness from a computation to a wall.
Had the defect been a small negative number, one could reasonably hope that a sharper treatment of
the Gamma factor, a better normalization, or a cleverer choice of correlation $c$ would push it
across zero. It is not small, it does not depend on $c$, and it does not improve with $r$: it gets
\emph{worse} linearly in $r$, so the ``more copies of the Euler factor'' strategy that motivated the
whole ladder $2\to3\to4$ makes this particular obstruction monotonically stronger.

\begin{remark}[the honest reading]\label{rmk:halflog2-reading}
Wall~\ref{ob:halflog2} does not say that the global sign of the gate is wrong. It says that the
region $(0,\log2)$ --- which is fully explicit, containing no arithmetic --- has a definite negative
contribution to any pointwise positivity claim. Whether the arithmetic mass on $[\log2,\infty)$
dominates that explicit negative mass in the \emph{integrated} sense is a completely different
question, and it is not answered by this wall. Indeed the natural response to
Wall~\ref{ob:halflog2} is not to look for another positive representation but to split the Laplace
variable at $\log2$ --- the point where arithmetic begins, and not an adjustable parameter --- and
compare the exactly computable negative mass on $(0,\log2)$ against the arithmetic mass on
$[\log2,\infty)$. That comparison is a signed, non-local statement, not a positivity, and it is
therefore of the right type. We return to it in \S\ref{sec:successor}.
\end{remark}

%% ======================================================================
\section{Synthesis: why sixteen failures are one failure}\label{sec:synthesis}
%% ======================================================================

We now assemble the catalogue and prove the classification that the paper's title asserts.

\subsection{The catalogue}

Table~\ref{tab:census} records the sixteen mechanisms. The column ``output'' states the strongest
conclusion the mechanism can reach about the target, assuming everything it wanted to be true were
true. The column ``class'' assigns it to Definition~\ref{def:semidef}. The column ``witness''
names the falsifier of \S\ref{sec:falsifiers} that closed it. We note that the count depends on
how one bundles the families --- \S\ref{sec:cm} alone closes four classical routes, and the earlier
running summary of this work counted fourteen rather than sixteen --- and that the precise number is
immaterial. What matters is the last two columns, which are constant down the table.

{\small
\begin{longtable}{@{}p{0.235\textwidth}p{0.215\textwidth}p{0.115\textwidth}p{0.185\textwidth}p{0.12\textwidth}@{}}
\caption{The sixteen closed mechanisms.}\label{tab:census}\\
\toprule
\textbf{mechanism} & \textbf{object produced} & \textbf{class} &
\textbf{witness} & \textbf{output} \\
\midrule
\endfirsthead
\multicolumn{5}{@{}l}{\small\itshape Table~\ref{tab:census}, continued.}\\
\toprule
\textbf{mechanism} & \textbf{object produced} & \textbf{class} &
\textbf{witness} & \textbf{output} \\
\midrule
\endhead
\bottomrule
\endfoot
Local PSD of the order kernel $w(\max(i,j))$ (\S\ref{sec:m1-local})
& regulated identity \eqref{eq:m1-finite}; block decomposition \eqref{eq:m1-blocks}; inertia
\eqref{eq:inertia}
& (S1)
& $H(2969)<-21<110<H(3167)$
& $\ \ge0$ \\[2pt]
Global Gram/energy of the same kernel (\S\ref{sec:m1-global})
& max-energy identity \eqref{eq:max-energy}; affine identity \eqref{eq:affine-exact}
& (S1)
& Hessian $\equiv0$; load \eqref{eq:load}
& $\ \ge0$ \\[2pt]
Non-local zero--lobe Gram (\S\ref{sec:m3})
& Weil form \eqref{eq:weil-a1}; stationary map \eqref{eq:stationary}
& (S1)
& bulk quartet, $\vartheta=\pi/50$
& $\ \ge0$ \\[2pt]
M\"obius inversion as positive adjoint (\S\ref{sec:m5-gram})
& degree convolution \eqref{eq:degree-conv}, \eqref{eq:m5-paired}
& (S1)
& quotient comb on $\{1,p\}$
& $\ \ge0$ \\[2pt]
Positive Euler--Gamma Koszul complex (\S\ref{sec:koszul})
& $d_E^2=0$; mixed-channel cancellation; \eqref{eq:MdZ}
& (S7)
& $N_p$ null-homotopic
& $\ \ge0$ \\[2pt]
Jordan hierarchy $\mu*\log^k\ge0$ (\S\ref{sec:jordan})
& cocycle \eqref{eq:jordan} and the full positive hierarchy
& (S4)
& $C_1<0<C_6$; $X_a$; Schur $\Leftrightarrow\RH$
& $\ \ge0$ \\[2pt]
Selberg $\mu*\log^2\ge0$ + $O(x)$ + Riccati (\S\ref{sec:selberg})
& exact degree recurrence
& (S4)
& dual load $\ge n(n-1)^n-1$
& $\ \ge0$ \\[2pt]
$\mathrm{PF}_2/\mathrm{TP}_2$, log-concavity (\S\ref{sec:pf2})
& minors \eqref{eq:pf2}, \eqref{eq:pf2-general}
& (S3)
& $-\tfrac1{2916}$; $\tfrac{11-8\sqrt2}{8}$
& $\ \ge0$ \\[2pt]
Compound-Poisson + convolution order (\S\ref{sec:cp})
& L\'evy measures \eqref{eq:nuA}, \eqref{eq:nuK}; coupling in $u$
& (S4),(S5)
& oscillatory tests are not monotone
& $\ \ge0$ \\[2pt]
Complete monotonicity / Bernstein (\S\ref{sec:cm})
& exact density \eqref{eq:h-exact}
& (S4)
& $-\sqrt2+\artanh\tfrac1{\sqrt2}<0$
& $\ \ge0$ \\[2pt]
Stein--Mecke generator positivity (\S\ref{sec:stein})
& Mecke identity \eqref{eq:mecke}; recurrence \eqref{eq:stein-rec}
& (S4)
& $(\tfrac52-2\sqrt2)/y_0<0$; $\sqrt2-3<0$
& $\ \ge0$ \\[2pt]
Radial Abel / Ces\`aro--Fej\'er means (\S\ref{sec:abel})
& exact quartet sums; radial germ positivity
& (S6)
& $w=2i$: $Q_4=-\tfrac{225}{8}$
& $\ \ge0$ \\[2pt]
Real order of the correlated shift (\S\ref{sec:shift})
& finite recurrence and renewal equation
& (S5)
& $C_6>0$; third difference in $c$
& $\ \ge0$ \\[2pt]
Three-channel Jordan--Beta expansion (\S\ref{sec:three-channel})
& expansion \eqref{eq:three}; limit \eqref{eq:three-limit}
& (S4)
& $\delta_0+c(cr-2)e^{-r}dr$; Pringsheim
& $\ \ge0$ \\[2pt]
Local tower algebraic square $\Cc_p^2$ (\S\ref{sec:tower})
& differences \eqref{eq:tower-2}, \eqref{eq:tower-r}; square \eqref{eq:tower-square}
& (S2)
& threshold \eqref{eq:threshold-m}; sign $(-1)^n$
& $\ \ge0$ \\[2pt]
Quartic sum of squares (\S\ref{sec:sos})
& factorization \eqref{eq:sos-root}; tower sum \eqref{eq:sos-tower}
& (S2)
& $(1-B_u)^2/u\to0$; $n=150$ quartet
& $\ \ge0$ \\[2pt]
Favourable tail sign from the VK envelope (\S\ref{sec:mech-envelope})
& two-sign saturation \eqref{eq:tail-signs}
& ---
& \eqref{eq:tail-lower}, \eqref{eq:tail-smooth}
& $\ \ge0$ \\
\end{longtable}
}

\subsection{The classification theorem}

\begin{theorem}[the category mismatch]\label{thm:classification}
Let $\mathcal M$ be any of the mechanisms of Table~\ref{tab:census}, and suppose --- contrary to the
witnesses --- that the positivity it seeks were established. Then the conclusion $\mathcal M$
delivers about the target is of the form $D_n^{[m]}\ge0$ for an even integer $m\in\{2,4\}$, or of the
form $\Delta D_n^{[m]}\ge0$, or $\lambda_n^{\mathrm{prime}}\ge-A_n/m$; in every case a
semidefiniteness statement in the sense of Definition~\ref{def:semidef}. If $m=2$ this is
sufficient for $(\star)$; if $m=4$, its bare nonnegativity does not force the compact certificate
at the prescribed floor through \eqref{eq:a0plus}, and the convenient floor-$1000$ sufficient
condition exceeds it by
\[
   \frac{A_n}{1001}=\frac{n\log n}{2002}\bigl(1+o(1)\bigr)\longrightarrow\infty .
\]
This last statement concerns only the formal transfer from the listed auxiliary
nonnegativities: those facts alone yield no positive lower bound on a pairing with a
sign-changing test. It does not preclude an additional arithmetic argument, and it does not make
a proportional bound necessary for RH.
\end{theorem}

\begin{remark}[the quartic case fails by non-existence, not by weakness]\label{rmk:quartic-exists}
It must be said plainly that the $m=4$ clause above is \emph{not} the reason the quartic
mechanism fails. If a certificate delivered $D_n^{[4]}\ge0$ for every $n\ge150$, then
$\lambda_n\ge\tfrac14A_n>0$ there, and Li's criterion together with the verified finite range
would give the Riemann Hypothesis outright, with no passage through $(\star)$ and no slack to
recover. The quartic route fails for a different and more basic reason: \emph{no such
certificate exists}. Section~\ref{sec:sos} proves that the fourth power is an algebraic square
and not a metric one, that the local fourth convolution fails $\mathrm{PF}_2$
(Witness~\ref{wit:pf2-quartic}), that the genuinely quadratic part vanishes in the extraction
limit \eqref{eq:sos-lin}, and that a certificate using only symmetry would falsify an off-line
divisor already at $n=150$ (Witness~\ref{wit:sos-offline}). The same holds at $m=2$ and $m=3$.
The catalogue is a statement about existence, and only secondarily about strength.
\end{remark}

\begin{proof}
The first assertion is a reading of the table: \S\S\ref{sec:m1-local}--\ref{sec:m1-global} and
\ref{sec:three-channel}, \ref{sec:tower} produce $D_n^{[2]}$ or $\Delta D_n$;
\S\S\ref{sec:selberg}, \ref{sec:shift}, \ref{sec:jordan} produce the first difference $C_n$;
\S\S\ref{sec:cm}, \ref{sec:stein}, \ref{sec:fejer} produce $D_n^{[3]}$; \S\ref{sec:sos} produces
$D_n^{[4]}$. The quantitative statement is Theorem~\ref{thm:threshold}. For the last assertion, note
that each mechanism's final step has the form
\[
   \text{target}=\int T\,d\pi\quad\text{or}\quad\text{target}=\langle T,\mathcal K\,T\rangle,
\]
with $\pi\ge0$ or $\mathcal K\succeq0$ the auxiliary object and $T$ the Laguerre test. Nonnegativity
of $\pi$ or $\mathcal K$ bounds the expression below by $0$ when $T\ge0$ or when the form is a
square; it gives no bound at all when $T$ changes sign, and it certainly gives no bound proportional
to a divergent quantity $A_n$ built from a different (archimedean) channel. To produce
$\kappa A_n$ one would have to bound below the pairing of a positive object with an oscillatory test
by a positive multiple of an independent scale, which is a coercivity statement in the sense of
Definition~\ref{def:coerc}. On the other hand, if a mechanism really established
$D_n^{[4]}\ge0$ for all large $n$, then $\lambda_n\ge A_n/4>0$ and it would prove RH directly.
The classification therefore concerns the attempted transfer to a prefixed compact certificate,
not the logical sufficiency of quartic positivity.
\end{proof}

\begin{corollary}[scope of the common obstruction]\label{cor:one-failure}
The sixteen entries of Table~\ref{tab:census} are instances of one obstruction to the proposed
\emph{fixed-cutoff transfers}: auxiliary nonnegativity neither controls the signed tail nor gives
a positive margin against a sign-changing test. No purely formal recombination preserving only
those auxiliary facts repairs that transfer. A new signed arithmetic input, an adapted-cutoff
argument proving strict quartic positivity, or a density-relaxed criterion is outside this
corollary.
\end{corollary}

\subsection{Two invariants of the catalogue}

Beyond the classification, two quantities recur across the entries and are worth naming, because
they measure the two ways the mechanisms fail.

\begin{definition}[the absolute-value load]\label{def:load}
Whenever a mechanism exits by taking absolute values on the kernel --- which the discipline of the
discipline adopted here forbids, but which is the only available exit once a Hessian cancels --- the cost is the
absolute load $\mathcal B_n$ of \eqref{eq:load}, satisfying
$\mathcal B_n\ge\exp\bigl(\tfrac23n\log n+\tfrac13n\log\log n-O(n)\bigr)$, against a
first-difference budget of size $O(\log n)$. In the regulated coordinate the same loss appears as
$\eps^{-n-1}$, i.e.\ $n^{n+1}$ at the natural scale $\eps=1/n$ (\S\ref{sec:selberg}).
\end{definition}

\begin{definition}[the propagation defect]\label{def:defect}
Whenever a mechanism requires a positive representation of the completed cocycle, the cost is the
defect \eqref{eq:defect} at $x_0=\tfrac12\log2$, namely $-(1-1/\sqrt2)r-1/\sqrt2$, which is
bounded away from zero and increases linearly in the exponent $r$.
\end{definition}

Every entry of Table~\ref{tab:census} pays one of these two. Definition~\ref{def:load} is the price
of abandoning the coupling; Definition~\ref{def:defect} is the price of demanding a positive
measure in the particular cocycle coordinate. These exhaust the transfers catalogued here, not
all possible positivity arguments for the Li coefficients.

\begin{remark}[four recurring obstructions, seen from here]\label{rmk:walls}
It is worth naming the four general obstructions that the sixteen entries instantiate, since each is
familiar from other approaches to the Hypothesis. \emph{Arithmetic propagation} is
Wall~\ref{ob:halflog2}, localized in the Laplace variable. \emph{Sign versus magnitude} appears as
Definition~\ref{def:load}: absolute envelopes measure magnitude, and the target needs sign.
\emph{Uniform norm}, or rank-one escape, appears in \S\ref{sec:successor}, where the relevant
operator bound turns out to be saturated exactly. And \emph{Weil positivity} appears as
Theorem~\ref{thm:weil-a1}: the target \emph{is} a Weil positivity statement, which is why no
reformulation of it can be easier than itself.
\end{remark}

%% ======================================================================
\section{What remains after the cutoff and quantifier corrections}\label{sec:successor}
%% ======================================================================

A negative result earns its keep only if it narrows the search. This section first audits two
stronger routes retained from the fixed-cutoff programme: operator coercivity is unavailable,
whereas fixed-vector non-alignment is not ruled out. It then uses the dominant-mode structure of
Bombieri--Lagarias to weaken the quantifier itself. The resulting block, density and nonlinear
partition criteria are logically sufficient for RH without coefficientwise positivity or a
proportional margin.

\subsection{Operator-norm coercivity is unavailable, not merely hard}\label{sec:no-operator-gap}

The most natural coercivity certificate in the sense of Definition~\ref{def:coerc}(C1) would be a
spectral gap for the whitened arithmetic jet: writing $\mathcal P_\Lambda$ for the prime block and
$S$ for the archimedean whitening, one would want
\begin{equation}\label{eq:gap-false}
   \lambda_{\max}\bigl(S\mathcal P_\Lambda S\bigr)\ \le\ 1-c
\end{equation}
for a fixed $c>0$. Structural and numerical work in the broader investigation from which this paper
is drawn (\S\ref{sec:availability}) establishes three things about that quantity: the sharp form of
the prime-side reduction is exactly $\lambda_{\max}\le1$; $\zeta$ sits \emph{at} the marginal value;
and the measured margin decays to zero exponentially in the window size, with twelve digits of
saturation already at window $N=8$. Consequently:

\begin{proposition}[no proportional gap in the operator coordinate]\label{prop:no-gap}
Under the Riemann Hypothesis the extremal quantity is pinned at the positivity floor,
$\lambda_{\max}(S\mathcal P_\Lambda S)=1$. Hence a statement of the form
$\lambda_{\max}\le1-c$ with fixed $c>0$ is \emph{false} if the Hypothesis holds, and cannot serve
as a lemma toward it. The same applies to the quaternionic Hodge--Riemann polarization, whose
spectrum on the relevant eigenspace decays exponentially with no uniform lower bound.
\end{proposition}

Proposition~\ref{prop:no-gap} eliminates a whole quadrant of the coercivity family in one stroke,
and it does so for a reason worth internalizing: the extremal problem \emph{is} Weil positivity,
and Weil positivity is marginal by construction. Any faithful reformulation of it will be marginal
too. Coercivity, if it exists at all, cannot live at the level of the operator.

\subsection{The extremal problem is not the problem $(\star)$ poses}\label{sec:fixed-family}

Here is a distinction that is easy to miss and that materially changes the outlook. Write, at a
finite cutoff,
\begin{equation}\label{eq:mu-n}
   \mu_n:=\frac{\bigl|\lambda_n^{\mathrm{prime}}\bigr|}{A_n}
   =\frac{\text{prime-side energy of }h_n}{\text{archimedean energy of }h_n},
\end{equation}
so that by \eqref{eq:split} and Remark~\ref{rmk:75} the target \eqref{eq:frontier} is
\begin{equation}\label{eq:mu-target}
   \mu_n\ \le\ 1-c\ \approx\ \tfrac34 .
\end{equation}
The ``gapless'' and ``marginal'' verdicts one meets in this subject concern
$\sup_h\mu(h)$, the supremum over \emph{all} admissible test functions. That supremum equals $1$;
that is the content of Weil positivity, and it is marginal by definition. But $(\star)$ does
not take a supremum. It concerns one explicit sequence $h_n=f_n*f_n^{\#}$, and for that sequence
the true value is, under the Hypothesis,
\begin{equation}\label{eq:mu-true}
   \mu_n=\frac{O(\sqrt n\log n)}{\tfrac12n\log n\,(1+o(1))}=O\bigl(n^{-1/2}\bigr)\longrightarrow0 .
\end{equation}
So the fixed family is not near the extremum, and a gap for it is not excluded by the marginality
of the extremum. We state this explicitly because reading the gapless verdicts quickly would
suggest that the route of this paper is self-defeating, and it is not: marginality is a property of
the supremum, not of every member.

\begin{remark}[what the distinction costs]\label{rmk:distinction-cost}
The price of the distinction is that one loses every tool that works uniformly over the test class.
A statement about a fixed sequence cannot be proved by bounding an operator; it requires knowing
something about that sequence's position relative to the operator's near-extremal directions. This
is precisely Definition~\ref{def:coerc}(C3), and it is the only cell of the coercivity table that
Proposition~\ref{prop:no-gap} does not empty. The density-relaxed route below does not belong to
that table at all.
\end{remark}

\subsection{A remaining stronger route: fixed-vector coercivity and non-alignment}\label{sec:nonalign}

\begin{problem}[uniform non-alignment]\label{prob:nonalign}
Prove that there exists $c>\tfrac14$ such that, for all $n\ge150$,
\begin{equation}\label{eq:nonalign}
   \bigl\langle f_n,\bigl(I-S\mathcal P_\Lambda S\bigr)f_n\bigr\rangle
   \ \ge\ c\,\|f_n\|^2 ,
\end{equation}
i.e.\ that the Laguerre test vector is uniformly separated from the near-extremal subspace of the
whitened arithmetic jet.
\end{problem}

Problem~\ref{prob:nonalign} is a coercivity statement for a fixed vector against a \emph{gapless}
operator, and that is a recognized and distinct kind of problem: it is not solved by exhibiting a
spectral gap, because there is none, but by a resolvent or quantitative-unique-continuation
estimate showing that $f_n$ cannot concentrate in the subspace where the operator approaches $1$ ---
a subspace whose dimension grows with the window. It satisfies the five requirements that
\S\ref{sec:synthesis} imposes on any successor: it is proportional rather than semidefinite; it is
not an absolute bound, so it does not pay the load of Definition~\ref{def:load}; it does not require
a positive representation, so it does not pay the defect of Definition~\ref{def:defect}; it uses
$\Lambda(m)\ge0$ non-termwise, through the operator, so it is not refuted by Falsifier~II; and it
does not appear in Table~\ref{tab:census}. It is, of course, of the strength of the Riemann
Hypothesis. The claim is only that it is a coherent stronger route, not that every proof must take
this form.

\begin{remark}[a cheap diagnostic for the optional non-alignment route]\label{rmk:diag-overlap}
Before developing Problem~\ref{prob:nonalign} it is worth computing the overlap
$\langle f_n,v_{\max}\rangle/\|f_n\|$ of the Laguerre test with the top eigenvector of the whitened
jet, at moderate window sizes and several base points. If the overlap decays, non-alignment is
plausible and the investment is justified; if it is $\Theta(1)$, Problem~\ref{prob:nonalign} is dead
before it starts and weeks are saved. This is a diagnostic in the sense of
Convention~\ref{conv:standard}, and would enter no proof.
\end{remark}

\subsection{A weaker route: syndetic excursions and density-relaxed Li criteria}
\label{sec:density-criterion}

The coefficientwise quantifier can be weakened much further.  The required input is already
present in the dominant-mode proof of Bombieri--Lagarias
\cite[Theorem~1(c)]{BombieriLagarias1999}, but the syndetic conclusion is useful enough to state
separately.

\begin{theorem}[syndetic negative excursions if RH fails]\label{thm:syndetic-li}
If RH is false, there exist $R>1$, $c>0$, $L\in\mathbb N$, $n_0$, and a set
$\mathcal S\subset\mathbb N$ with gaps at most $L$ such that
\begin{equation}\label{eq:syndetic-excursion}
   \lambda_n\le-cR^n\qquad(n\in\mathcal S,\ n\ge n_0).
\end{equation}
In particular $\mathcal S$ has positive lower natural and logarithmic density.  The constants
depend on the exterior zeros; no universal positive density follows from this argument.
\end{theorem}

\begin{proof}
For $u_\rho=\rho/(\rho-1)$, failure of RH gives $|u_\rho|>1$ for some zero on the right of the
critical line. Functional symmetry gives $\lambda_{-n}=\lambda_n$, so the dominant-mode
calculation with $u_\rho^n$ has the orientation used below. The hypotheses used in the proof of
Bombieri--Lagarias imply that the maximal
modulus $R>1$ is attained by finitely many modes
$u_j=Re^{i\phi_j}$ and that, for some $R_1<R$,
\[
   \lambda_n=-R^n\sum_{j=1}^{K}\cos(n\phi_j)
              +O(n^2R_1^n+n^2).
\]
Put $\alpha=(\phi_1,\ldots,\phi_K)$ and
$H=\overline{\{n\alpha:n\in\mathbb Z\}}\subset\mathbb T^K$.  On the compact monothetic group
$H$, rotation by $\alpha$ is minimal.  A sufficiently small neighbourhood $U$ of the identity
satisfies $\sum_j\cos x_j>K/2$ on $U$.  Finitely many translates of $U$ cover $H$, so the return
set $\{n:n\alpha\in U\}$ has bounded gaps.  On its tail the displayed error is
$o(R^n)$, proving \eqref{eq:syndetic-excursion}.  A set with gaps at most $L$ has lower density at
least $1/L$, and partial summation gives the logarithmic-density assertion.
\end{proof}

For $E\subset\mathbb N$ write
\[
 \overline\delta_{\log}(E)=\limsup_{X\to\infty}\frac{1}{\log X}
       \sum_{\substack{n\le X\\n\in E}}\frac{1}{n}.
\]

\begin{theorem}[subexponential density and block criterion]\label{thm:density-block}
Let $b_n\ge0$ satisfy $\log(1+b_n)=o(n)$.  The following are equivalent:
\begin{enumerate}[label=\textup{(\roman*)},leftmargin=*]
\item RH;
\item there is $E$ with $\overline\delta_{\log}(E)=0$ such that
      $\lambda_n\ge-b_n$ for every sufficiently large $n\notin E$;
\item there are consecutive integer intervals $I_j$, with $|I_j|\to\infty$, on each of which
      $\lambda_n\ge-b_n$ at every index;
\end{enumerate}
The theorem remains true with $b_n=1$.  It also implies the quartic variant: RH is equivalent to
$4\lambda_n>A_n$ outside an exceptional set of logarithmic density zero.
\end{theorem}

\begin{proof}
Under RH, Li's criterion gives $\lambda_n\ge0$ for every $n$, so (ii)--(iii) follow (and the
quartic statement follows eventually from
$\lambda_n=A_n+O(\sqrt n\log n)$).  If RH is false, choose the syndetic set of
Theorem~\ref{thm:syndetic-li}.  Since $b_n=e^{o(n)}$, its tail satisfies
$b_n<cR^n$, so every sufficiently large element of $\mathcal S$ violates the lower bound.
Positive logarithmic density rules out (ii), and bounded gaps rule out (iii).
\end{proof}

Thus one need not prove $(\star)$, quartic positivity, or a proportional estimate at every degree.  It
would be enough to prove the constant lower bound $\lambda_n\ge-1$ on consecutive blocks of
unbounded length.  This is a genuine weakening of the target for arbitrary sequences, although
it remains RH-equivalent for the Li coefficients.

There is also a bounded nonlinear detector which retains the density information while discarding
the exponentially large amplitude.  Put
\[
 D_n=4\lambda_n-A_n,\qquad
 H_{n,b}=\frac{(-D_n)_+}{b_n+(-D_n)_+},
\]
where $b_n>0$ and $\log b_n=o(n)$.  Theorem~\ref{thm:syndetic-li} gives the following immediate
criterion:
\begin{equation}\label{eq:saturated-barrier}
 \boxed{\ \RH\quad\Longleftrightarrow\quad
 \exists I_j:\ \min I_j,|I_j|\to\infty,
 \quad \frac{1}{|I_j|}\sum_{n\in I_j}H_{n,b}\longrightarrow0.\ }
\end{equation}
Indeed, under RH one has $D_n>0$ eventually.  If RH is false, then on the syndetic set of
Theorem~\ref{thm:syndetic-li} one has $-D_n\gg R^n$, so $H_{n,b}\to1$ there; every sufficiently
long interval therefore has average bounded below by a positive constant.  The saturation is
essential: it detects a density of bad degrees rather than paying for their exponential height.
It is neither a linear Abel mean nor a polynomial additive identity.

For a smooth detector with an exact arithmetic representation, assume additionally that
$b_n\to\infty$, fix $t>0$, put $H_X=\sum_{n\le X}n^{-1}$, and define the bounded Fermi--Dirac
average
\begin{equation}\label{eq:fermi}
 \mathfrak F_{t,b}(X)=\frac1{H_X}\sum_{n\le X}\frac1n
       \frac1{1+\exp\bigl(t(\lambda_n+b_n)\bigr)}.
\end{equation}

\begin{theorem}[bounded Fermi criterion]\label{thm:fermi}
For every nonnegative $b_n\to\infty$ with $\log(1+b_n)=o(n)$, and every fixed $t>0$,
\begin{equation}\label{eq:fermi-bound}
   \boxed{\ \RH\quad\Longleftrightarrow\quad
           \mathfrak F_{t,b}(X)\longrightarrow0.\ }
\end{equation}
\end{theorem}

\begin{proof}
Under RH, $\lambda_n\ge0$, so the summand in \eqref{eq:fermi} is at most
$(1+e^{tb_n})^{-1}\to0$; its logarithmic mean tends to zero.  If RH is false, then on the
positive-logarithmic-density set $\mathcal S$ of Theorem~\ref{thm:syndetic-li} one has
$\lambda_n+b_n\to-\infty$.  The corresponding Fermi summands tend to one, so
$\liminf_X\mathfrak F_{t,b}(X)>0$.
\end{proof}

The Fermi criterion has an exact arithmetic form which keeps the pole, the archimedean factor,
all primes and all prime powers coupled until after the nonlinear observable is formed.  For
$\varepsilon>0$ set
\[
 p_n(\varepsilon)=n\sum_{k=1}^{n}\binom{n-1}{k-1}
       \frac{(-1)^{k-1}}{k\varepsilon^k},
\qquad
 \mathcal E_{n,\varepsilon}=p_n(\varepsilon)
 -\sum_{m\ge2}\frac{\Lambda(m)}{m^{1+\varepsilon}}
       L_{n-1}^{(1)}(\log m).
\]
The prime--Laguerre identity is
$\lambda_n=A_n+\lim_{\varepsilon\downarrow0}\mathcal E_{n,\varepsilon}$; hence, for each finite
$X$,
\begin{equation}\label{eq:fermi-arith}
\boxed{
 \mathfrak F_{t,b}(X)=\lim_{\varepsilon\downarrow0}\frac1{H_X}
 \sum_{n\le X}\frac1n
 \left[1+e^{t(A_n+b_n+p_n(\varepsilon))}
 \prod_{m\ge2}\exp\!\left(
 -\frac{t\Lambda(m)}{m^{1+\varepsilon}}L_{n-1}^{(1)}(\log m)
 \right)\right]^{-1}. }
\end{equation}
For fixed $\varepsilon>0$ the logarithm of the product converges absolutely, and the finite outer
sum permits passage to the limit by continuity.  Thus proving \eqref{eq:fermi-bound} directly
from \eqref{eq:fermi-arith}, for example with $b_n=\log(n+1)$, is the principal density-relaxed
arithmetic front.  The identity itself proves no estimate.

There is a related but strictly stronger-shaped exponential auxiliary:
\begin{equation}\label{eq:partition-bound}
 \mathcal Z_t(X):=\sum_{n\le X}\frac{e^{-t\lambda_n}}n=O_t(\log X).
\end{equation}
It too is equivalent to RH, and its exact regularized representation is
\begin{equation}\label{eq:partition-arith}
 \mathcal Z_t(X)=\lim_{\varepsilon\downarrow0}
 \sum_{n\le X}\frac{e^{-t(A_n+p_n(\varepsilon))}}{n}
 \prod_{m\ge2}
 \exp\!\left(\frac{t\Lambda(m)}{m^{1+\varepsilon}}
                 L_{n-1}^{(1)}(\log m)\right).
\end{equation}
Indeed, RH gives $e^{-t\lambda_n}\le1$ and hence $\mathcal Z_t(X)\le H_X$; if RH is false,
one term along the syndetic set of Theorem~\ref{thm:syndetic-li} grows like
$n^{-1}\exp(tcR^n)$ and violates every $O(\log X)$ bound.
But \eqref{eq:partition-bound} does not exploit the density relaxation.  Taking $X=n$ and keeping
the $n$-th nonnegative summand gives
\begin{equation}\label{eq:partition-pointwise}
 \lambda_n\ge-\frac1t\bigl(\log n+\log\log n+O_t(1)\bigr)
 \qquad(n\ge2).
\end{equation}
Thus $\mathcal Z_t$ is retained only for its simple algebraic factorization.  The Fermi average is
the weaker front: it is bounded by one before averaging, so an isolated excursion cannot replace
a positive density of excursions.  Jensen has the wrong direction, taking absolute values
destroys the cancellation as $\varepsilon\downarrow0$, and the linear Abel means catalogued above
control neither nonlinear observable.

\begin{remark}[post-audit of the first two Fermi attacks]\label{rmk:fermi-postaudit}
Two natural continuations of \eqref{eq:fermi-arith} have been closed quantitatively.  First, the
Fourier identity
\[
 (1+e^y)^{-1}=\frac12-\int_0^\infty\frac{\sin(sy)}{\sinh(\pi s)}\,ds
\]
lifts the complete prime--pole block to an exact product of unit-modulus phases, so it removes the
artificial explosion caused by separating positive exponentials.  It does not remove the decisive
scale: an excursion of depth $Y$ is detected at $s=O(Y^{-1})$.  For the rational off-line quartet
$w=2i$, one has $Y_n\ge2^n$ on $n\equiv0\pmod4$, so every fixed-frequency or polynomial-resolution
estimate misses the Fermi signal.  The diagonal change $s=e^{-nv}$ moves an exponential rate to a
fixed $v$, but the Laguerre generating function then controls a \emph{geometric product} of phases,
whereas Fermi requires their \emph{arithmetic mean}; at the critical scale all Hadamard powers
survive.

Second, the cross-degree Parseval and Christoffel--Darboux identities are exact, but their
archimedean--pole square, prime square and mixed term are each of order
$\varepsilon^{-4N}$ on degrees $N<n\le2N$.  Only their complete signed cancellation has a finite
limit.  Positivity of the Christoffel--Darboux kernel controls one piece and supplies no upper bound
for the coupled energy.  Thus the unitary microfrequency branch and the bare Parseval/CD branch are
both no-gos.  They add useful exact identities, but \eqref{eq:fermi-bound} remains an equivalence
without an unconditional estimate.
\end{remark}

\subsection{The other live successor: splitting the Laplace variable at $\log2$}
\label{sec:split}

Wall~\ref{ob:halflog2} refutes \emph{pointwise} positivity in a region containing no primes. It
says nothing about the signed total, and Remark~\ref{rmk:halflog2-reading} already indicated the
natural response: do not ask the completed cocycle to be a positive measure; split at the point
where arithmetic begins,
\begin{equation}\label{eq:split-log2}
   \int_0^\infty\varphi\,d\mu
   =\underbrace{\int_0^{\log2}\varphi\,d\mu}_{\text{no primes: exact, closed form}}
   +\underbrace{\int_{\log2}^\infty\varphi\,d\mu}_{\text{arithmetic: }\Lambda\ge0,\ \text{atoms at }k\log p} ,
\end{equation}
and compare the exactly computable negative mass of the first block against the arithmetic mass of
the second. The first block involves only the Gamma and polar channels, both in closed form, with
first variation \eqref{eq:frq}; the second is where $\Lambda(m)\ge0$ lives. The cut at $\log2$ is
not an adjustable parameter --- earlier attempts of ours were forced to fix an arbitrary $\delta$
because shrinking windows scaled out of the estimate --- it is the location of the first atom.

This is exactly the ``specific signed non-local cancellation'' that \S\ref{sec:stein} left open, and
the ``global multiplicative sum retaining inter-prime interactions'' that
Remark~\ref{rmk:tower-alive} left open, with the correct partition supplied. We record the obvious
risk: the negative mass on $(0,\log2)$ may be of order $A_n$ and so may the positive mass, leaving
another marginality. If so, the matter closes quickly and with a witness, which is the standard
adopted throughout.

\subsection{The cushion in the stronger fixed-ratio route}\label{sec:cushion}

For the optional uniform fixed-ratio route, Remark~\ref{rmk:75} says that the prime side may
consume almost three quarters of the archimedean block, while under the Hypothesis it consumes
asymptotically none of it. So:

\begin{proposition}[a factor-four cushion against the truth]\label{prop:cushion}
Under the Riemann Hypothesis,
\[
   \frac{\lambda_n}{A_n}=1+O\bigl(n^{-1/2}\bigr),
\]
whereas the sufficient condition \eqref{eq:frontier} requires only $\lambda_n/A_n>\tfrac14$, and
the fixed-ratio variant of Theorem~\ref{thm:scalefree} only $\lambda_n/A_n\ge c$ for some fixed
$c>\tfrac14$ --- for instance $c=501/2002=0.250250\ldots$ at the declared cutoff floor. Either
way the route carries a cushion of a factor four against the truth and needs no sharp constant.
\end{proposition}

Proposition~\ref{prop:cushion} locates the difficulty of the proportional route. It does not say
that a proportional lower bound is necessary: Theorem~\ref{thm:density-block} permits a merely
subexponential lower bound on sparse blocks. Absolute envelopes still fail for the fixed-ratio route:
an estimate that loses $\exp\bigl(\tfrac23n\log n\bigr)$, as in Definition~\ref{def:load}, is not
``a bit too lossy'' for a target with a factor-four cushion --- it is off by a superpolynomial
factor, and no amount of cushion helps.

%% ======================================================================
\section{Scope, retractions and the open problems}\label{sec:scope}
%% ======================================================================

\subsection{What each no-go does and does not exclude}

Every no-go in this paper is relative to a declared list of admissible inputs, and we restate the
uniform disclaimer once, in full, because it is the difference between a catalogue and an
over-claim. Let $\mathcal M$ be a mechanism of Table~\ref{tab:census} with declared input list
$\mathcal I$. The no-go asserts:
\begin{quote}
no derivation of the target using only the inputs in $\mathcal I$ can succeed, because there exists
an object satisfying every element of $\mathcal I$ for which the target fails.
\end{quote}
It does \emph{not} assert that the target is false, nor that no derivation exists, nor that some
other mechanism using $\mathcal I$ plus additional arithmetic information must fail. In particular,
not one of the sixteen excludes:
\begin{enumerate}[label=\textup{(\arabic*)},leftmargin=*]
\item a signed, non-local inequality valid specifically for the sequence $\Lambda(m)$ ---
      required and unproved by \S\S\ref{sec:m1-global}, \ref{sec:m3}, \ref{sec:selberg},
      \ref{sec:stein}, \ref{sec:tower};
\item a non-local differential that kills only quadratic combinations of quotients without making
      the generators $N_p$ exact, together with an independent construction of its exact Gamma
      diagonal (\S\ref{sec:koszul});
\item a global renormalization or regrouping across coprime backgrounds, including the global term
      of the three-channel identity (\S\ref{sec:tower});
\item a signed inequality acting directly on the limit $q_{n-1}\to3\lambda_n-A_n$, requiring
      neither global holomorphy nor eventual sign at fixed regulator (\S\ref{sec:fejer});
\item a sign statement for the real tail that uses the support on prime powers and the exact values
      $\Lambda(m)$ essentially (\S\ref{sec:mech-envelope}).
\end{enumerate}

\subsection{Retractions and corrections carried by this paper}

We list the corrections applied in the course of this work, because an audit trail that hides its
own repairs is not an audit trail.
\begin{itemize}[leftmargin=*]
\item \emph{The sign of the incomplete Li block.} See Remark~\ref{rmk:erratum}. The corrected
      unconditional form is $\lambda_n=A_n+\lambda_n(\sqrt n)+O(\sqrt n\log n)$. Every corollary
      resting on an unconditional bound for $A_n-\lambda_n-\lambda_n(\sqrt n)$ is withdrawn,
      including a finite-range corollary derived from an effective constant. Making that constant
      effective is not an outstanding technical task: the asserted estimate itself had the wrong
      sign. The contour proof additionally requires repairing the horizontal segment, whose printed
      geometric ratio is $e/2>1$; taking $T=2\sqrt n+\eps_n$ gives ratio $e^{1/4}/2<1$ via
      $G_n(s+1)/G_n(s)=(1-1/(s+1)^2)^n$, after which one returns to $\sqrt n$ at cost
      $O(\sqrt n\log n)$.
\item \emph{A claimed free relaxation from lowering $\theta$.} Withdrawn: by
      Theorem~\ref{thm:theta} the transport is signed and the net bracket
      $-2\theta A_n\le\mathcal N_n(\theta)\le\tfrac12A_n$ has indeterminate sign. There is no
      discount.
\item \emph{A claimed equivalence between $(\star)$ and an upper bound for
      $\lambda_n(\sqrt n)$ ``with the same implied constant both ways''.} Withdrawn: the
      $O(\cdot)$ hides exactly the required sign. The correct statements are one sufficient and one
      necessary condition with an explicit gap of $2\theta A_n+2C_*\sqrt n\log n$.
\item \emph{A claimed derivation of $(\star)$ from $\lambda_n>0$.} Withdrawn: the exact
      condition is $\lambda_n\ge\tfrac14A_n+R_n(T_n)$. Li positivity alone is not enough; see the
      second registered trap below.
\item \emph{The constants $2002/501$ and $A_n/1001$ as the intrinsic frontier.} Withdrawn.
      They are the same coarse datum, obtained by replacing the actual cutoff with its validity
      floor $T\ge1000$.  Admissibility is upward closed.  Thus $D_n^{[4]}>0$ permits an adapted
      finite cutoff proving $(\star)$, while $D_n^{[4]}\ge0$ for all large $n$ proves RH directly.
      The fixed-floor threshold remains a valid sufficient condition and a valid historical audit
      of the transfers tested against it.
\item \emph{Two editions of the range formula.} The published version of \cite{Lagarias2007}
      (Ann.\ Inst.\ Fourier \textbf{57} (2007), 1689, p.~1694) prints
      $n\le T^2/2(\log T)^2$; the arXiv source prints $T^2/4(\log T)^2$. At $T=3\times10^{12}$
      these are $\approx5.5\times10^{21}$ and $\approx2.7\times10^{21}$. We register the editorial
      discrepancy and do not dismiss the larger range; determining why the smaller is stated
      remains open.
\item \emph{A numerical table without a stability guard.} Withdrawn and replaced: agreement between
      two Cauchy radii is necessary but not sufficient, since both radii share the same zeta
      evaluation, transform length and floating-point format, so a common systematic error is
      invisible. Rows are reported as ``stable between two radii'', never as ``a reliable range''.
\end{itemize}

\subsection{Two registered traps}

\begin{enumerate}[label=\textup{(T\arabic*)},leftmargin=*]
\item There is no unconditional bound
      $\bigl|A_n-\lambda_n-\lambda_n(\sqrt n)\bigr|\ll\sqrt n\log n$. The small remainder has the
      opposite sign, $A_n-\lambda_n+\lambda_n(\sqrt n)$, and using the old remainder as small
      injects exactly $\RH$-strength content.
\item For $(\star)$ it does not suffice that $\lambda_n>0$. The exact condition is
      $\lambda_n\ge\tfrac14A_n+R_n(T_n)$. Under the crude form of the tail estimate the strong margin
      $2\lambda_n-A_n\ge0$ was sufficient but never necessary; retaining the factor of
      Corollary~\ref{cor:a0plus}, $r_*\lambda_n-A_n\ge0$ with $r_*=2002/501$ is one uniform
      sufficient condition, the actual $T_n$-dependent margin is much smaller, and strict
      $D_n^{[4]}>0$ suffices after adapting the cutoff.  Separately, non-strict
      $D_n^{[4]}\ge0$ already proves RH without passing through $(\star)$.
\end{enumerate}

\subsection{The open problems, stated exactly}

\begin{problem}\label{prob:main}
Prove directly from the arithmetic side any one of the equivalent conditions of
Theorem~\ref{thm:density-block}; for example, prove $\lambda_n\ge-1$ on consecutive blocks of
unbounded length.  The principal smooth version is to prove
$\mathfrak F_{t,b}(X)\to0$ from \eqref{eq:fermi-arith}, for one $t>0$ and for example
$b_n=\log(n+1)$.  Unlike the exponential auxiliary $\mathcal Z_t$, this does not imply a
coefficientwise lower bound and retains the density relaxation.
\end{problem}

\begin{problem}\label{prob:a1-cofinal}
For the literal compact route, prove $D_n^{[4]}>0$ for every $n\ge150$.  By
Theorem~\ref{thm:four} this gives an admissible adapted cutoff at every index and proves
$(\star)$.  The stronger condition $D_n^{[4]}\ge A_n/1001$ is sufficient at the floor
$T=1000$, but is not equivalent to this problem and is not intrinsic.
\end{problem}

\begin{problem}\label{prob:global-sum}
Determine the sign of the global multiplicative sum \eqref{eq:global-gate}, retaining the
interactions between distinct primes, without a sign per tower, a sign per background, or a norm of
$\Cc$.
\end{problem}

\begin{problem}\label{prob:mass}
In the splitting \eqref{eq:split-log2}, compare the exact negative mass of the arithmetic-free block
$(0,\log2)$ against the arithmetic mass of $[\log2,\infty)$, at $r=4$ (or at a separately
declared fixed $r<4$) along $u=c\eps$.
\end{problem}

\begin{problem}\label{prob:effective}
Make the remainder constant $C_*$ in $\widetilde\eps_n=O(\sqrt n\log n)$ effective for $\zeta$;
the implied constants in \cite{Lagarias2007} depend on the automorphic datum and are not effective
as stated.
\end{problem}

Problem~\ref{prob:nonalign} is an optional stronger route. Problems~\ref{prob:main},
\ref{prob:a1-cofinal} and \ref{prob:nonalign} are of the strength of the Riemann Hypothesis;
Problems~\ref{prob:global-sum} and \ref{prob:mass} are concrete possible arithmetic inputs;
Problem~\ref{prob:effective} is bookkeeping that was deliberately deferred.

%% ======================================================================
\section{Reproducibility}\label{sec:repro}
%% ======================================================================

Every sign decision quoted as a certificate in this paper is exact. We record here which
statements were verified mechanically and in what arithmetic, so that a reader can reconstruct the
audit rather than trust it.

\begin{itemize}[leftmargin=*]
\item \emph{Rational arithmetic, no floating point.} The cumulative sign changes
      \eqref{eq:H-witness}, computed with outward-rounded rational intervals for the prime
      logarithms; the finite identities \eqref{eq:max-energy}, \eqref{eq:collapse} and
      \eqref{eq:hess-boundary}, verified with two perturbations --- one of total mass zero and one
      with $F(X)\ne0$ --- so that a wrong sign or an omitted boundary term makes the check fail; the
      $\mathrm{PF}_2$ minor \eqref{eq:pf2-witness} and the local identity \eqref{eq:pf2}; the
      integral certificate for the $\artanh$ inequality of \eqref{eq:cm-witness}; the operator
      identities \eqref{eq:P-form}, \eqref{eq:tower-2}, \eqref{eq:B-form}, the commutation
      $\Bb^2\Pp^2=(\Bb\Pp)^2$ on test polynomials, and the recomposition of the degree-one formula
      on both sides of its threshold \eqref{eq:threshold-m}, together with the leading coefficient
      \eqref{eq:leading-n} up to degree eight; the value $r_*=2002/501$, the limiting equality of
      Corollary~\ref{cor:a0plus}, the binomial signs of the polar channel and twelve coefficients of the
      hypergeometric sum \eqref{eq:1F1}; the Fej\'er coefficient identity \eqref{eq:fejer-coeff}
      and the witness $R=\tfrac12$, $\Phi(z)=6z^4$ of Witness~\ref{wit:fejer}; the quartic
      degree-two witness, the exact relation \eqref{eq:threshold-id} and the off-line phase at
      index $150$ of Witness~\ref{wit:sos-offline}; the exact quartet values of
      Witness~\ref{wit:abel}; the three-channel convolution on prime powers and the polar
      identities of Theorem~\ref{thm:polar}.
\item \emph{Double precision, diagnostic only.} The two Cauchy-extraction probes of
      Remarks~\ref{rmk:diag-cubic} and \ref{rmk:diag-global}. Both compare two radii; both share the
      same zeta evaluation, transform length and floating-point format, so their agreement is
      necessary but not sufficient, and a common systematic error would be invisible. Neither is
      used in the proof of anything, and rows are described as stable between two radii, never as a
      reliable range. For non-integer exponents the probe uses pointwise principal logarithms rather
      than an unwound analytic branch, so its non-integer rows are exploratory and are not cited as
      evidence anywhere in this paper.
\end{itemize}

\begin{remark}[the numerical trap that the guard was installed to catch]\label{rmk:guard}
At one index the two-radius disagreement was $1.1\times10^{2}$ while the quantity being certified
was $17.3$. Without an automatic discard rule --- reject a row unless the two-radius disagreement is
below one per cent of the quantity --- the table would have appeared to confirm a theorem on a row
with no content. This is the reason for the wording of Convention~\ref{conv:standard}.
\end{remark}

\begin{remark}[the diagnostic for the cubic margin]\label{rmk:diag-cubic}
For $\eps=0.1$, $c=\tfrac12$, a double-precision extraction of the normalized cubic cocycle gave
$h^{[3]}_{n-1,\eps,c}<0$ for $1\le n\le1201$, the largest observed value in the target range
$150\le n\le1201$ being approximately $-390.1$ at $n=150$. This is evidence, not a certificate.
\end{remark}

%% ======================================================================
\section{Epilogue}\label{sec:epilogue}
%% ======================================================================

The result of this work, stated without decoration, has two layers.  The first is a catalogue of
sixteen failed transfers, each closed by an explicit witness.  The second is a correction to the
interpretation originally placed on that catalogue.  The descent from $2$ through $3$ to
$2002/501$ concerned a uniform sufficient transfer at the coarse floor $T=1000$; it was not the
intrinsic frontier.  Because cutoffs are upward closed, strict quartic positivity realizes the
compact certificate at adapted cutoffs, and non-strict quartic positivity already proves RH
directly.  Thus the paper's valid negative result concerns the listed transfer mechanisms, not
semidefiniteness as a universal proof category.  What the title names survives that correction on
different ground, stated as Proposition~\ref{prop:thesis}: the object whose nonnegativity is at
stake is provably not a positive object, the defect is a fixed proportion rather than a small
perturbation, and what remains is a signed cancellation between two divergent blocks with a
proportional margin.

We do not think this reflects badly on the mechanisms, and we would build several of them again.
The order kernel identity of Theorem~\ref{thm:max-energy}, the tower differences of
Theorem~\ref{thm:tower-diff}, the hypergeometric polar density of
Theorem~\ref{thm:polar-1F1}, the Stein--Mecke recurrence of Theorem~\ref{thm:stein} and the
Fej\'er coefficient identity of \eqref{eq:fejer-coeff} are exact, and they are the kind of object
that survives its own no-go. What we do think is that the discipline which produced them --- state
the axioms first, build the falsifier to satisfy exactly those axioms, and accept the witness when
it arrives --- is the reason the sixteen deaths can now be read as one theorem rather than as
sixteen disappointments.

The quantifier audit then weakens the live task in a different direction.  Failure of RH forces
exponentially negative Li coefficients on a syndetic set.  It is therefore enough to prove a
subexponential lower bound on arbitrarily long consecutive blocks, or the bounded Fermi limit
\eqref{eq:fermi-bound}.  Equation \eqref{eq:fermi-arith} keeps pole, Gamma, primes and prime powers
coupled and states the density-relaxed arithmetic front exactly.  We do not estimate it here.
The exponential partition $\mathcal Z_t$ remains an algebraically simple but stronger auxiliary:
its $O(\log X)$ bound already implies the pointwise estimate \eqref{eq:partition-pointwise}.
Between the Fermi identity and its vanishing limit still lies the Riemann Hypothesis; this paper
supplies a narrower target and an audit trail, not its proof.

%% ======================================================================
\section{Source material and supporting computations}\label{sec:availability}
%% ======================================================================

The investigation from which this paper is drawn --- the full derivations, the verification scripts
in exact rational arithmetic, the numerical diagnostics with their stability guards, and the
complete record of the mechanisms catalogued above, including the ones whose deaths are only
summarized here --- is maintained in the public repository
\begin{center}
\url{https://github.com/dtrejopizzo/riemann}
\end{center}
A reader who wishes to reconstruct rather than trust the audit will find the underlying material in
three places within it.

\emph{The reduction of the criterion.} The identification of Li--Keiper positivity with the
localized Weil form, the completed autocorrelation of Theorem~\ref{thm:weil-a1}, the whitened
arithmetic jet and the numerical facts about $\lambda_{\max}(S\mathcal P_\Lambda S)$ quoted in
\S\ref{sec:no-operator-gap}, and the gapless quaternionic polarization referred to in
\S\ref{sec:fixed-family}, are developed in \texttt{03-research/phase-102-omega7-closure-campaign}.

\emph{The direct route and its finite certificates.} The compact--tail identity
\eqref{eq:CnT}, the unconditional tail estimate of Theorem~\ref{thm:a0} in its original form, the
rational certificates for the finitely many indices below the target range, the collective
order-kernel identity behind \S\ref{sec:m1-local}, the earlier non-local and Sturm--Liouville
closures referred to in \S\S\ref{sec:m3} and \ref{sec:selberg}, and the absolute-load estimate
\eqref{eq:load}, are in \texttt{03-research/phase-103-direct-a1-closure}.

\emph{The sixteen mechanisms of this paper, and their successors.} Every entry of
Table~\ref{tab:census}, with its declared inputs, its exact identities, its witness, and its
mechanical verifier, together with the margin ladder of \S\ref{sec:ladder}, the cutoff audit of
Remark~\ref{rmk:1000}, the sign correction of Remark~\ref{rmk:erratum} and the bibliographic
audit of \S\ref{sec:biblio}, is in
\texttt{03-research/phase-104-unconditional-a1-closure}. That directory also records the further
mechanisms closed after the sixteen catalogued here --- among them the connected-cumulant
classification, which shows that every additive polynomial identity in the joint moments of
several prime towers collapses to single-tower cumulants, and the Gibbs/Chernoff gate, which
fails because all nontrivial exponential moments of the unit-Palm defect diverge. Neither is
needed for the results of this paper, and both sharpen the conclusion of \S\ref{sec:successor}
that the surviving front is non-additive and non-exponential. The verification scripts live in the
\texttt{tools} subdirectory of that directory; those that decide a sign use Python's
\texttt{Fraction} arithmetic exclusively, and the two that do not are named as diagnostics in their
own headers.

No document in that repository may be cited as a proof of the Riemann Hypothesis, and none of them
claims to be.

\begin{thebibliography}{99}

%% --- Foundational and classical ---
\bibitem{Riemann} B.~Riemann, \emph{\"Uber die Anzahl der Primzahlen unter einer gegebenen
Gr\"osse}, Monatsber.\ Berliner Akad.\ (1859), 671--680.
\bibitem{Titchmarsh} E.~C.~Titchmarsh, \emph{The Theory of the Riemann Zeta-Function}, 2nd ed.\
(rev.\ D.~R.~Heath-Brown), Oxford Univ.\ Press, 1986.
\bibitem{Edwards} H.~M.~Edwards, \emph{Riemann's Zeta Function}, Academic Press, 1974; Dover, 2001.
\bibitem{Ivic} A.~Ivi\'c, \emph{The Riemann Zeta-Function: Theory and Applications}, Dover, 2003.
\bibitem{Davenport} H.~Davenport, \emph{Multiplicative Number Theory}, 3rd ed., rev.\ by
H.~L.~Montgomery, Grad.\ Texts in Math.\ \textbf{74}, Springer, 2000.
\bibitem{IwaniecKowalski} H.~Iwaniec and E.~Kowalski, \emph{Analytic Number Theory}, Amer.\ Math.\
Soc.\ Colloq.\ Publ.\ \textbf{53}, 2004.
\bibitem{MontgomeryVaughan} H.~L.~Montgomery and R.~C.~Vaughan, \emph{Multiplicative Number
Theory I: Classical Theory}, Cambridge Stud.\ Adv.\ Math.\ \textbf{97}, Cambridge Univ.\ Press,
2007.
\bibitem{ConreyNotices} J.~B.~Conrey, \emph{The Riemann Hypothesis}, Notices Amer.\ Math.\ Soc.\
\textbf{50} (2003), no.~3, 341--353.
\bibitem{Bombieri} E.~Bombieri, \emph{Problems of the Millennium: The Riemann Hypothesis}, Clay
Mathematics Institute, 2000.
\bibitem{Sarnak} P.~Sarnak, \emph{Problems of the Millennium: The Riemann Hypothesis} (2004), Clay
Mathematics Institute report.

%% --- Li--Keiper coefficients ---
\bibitem{Li1997} X.-J.~Li, \emph{The positivity of a sequence of numbers and the Riemann
hypothesis}, J.\ Number Theory \textbf{65} (1997), no.~2, 325--333.
\bibitem{BombieriLagarias1999} E.~Bombieri and J.~C.~Lagarias, \emph{Complements to Li's criterion
for the Riemann hypothesis}, J.\ Number Theory \textbf{77} (1999), no.~2, 274--287.
\bibitem{Keiper1992} J.~B.~Keiper, \emph{Power series expansions of Riemann's $\xi$ function},
Math.\ Comp.\ \textbf{58} (1992), no.~198, 765--773.
\bibitem{Lagarias2007} J.~C.~Lagarias, \emph{Li coefficients for automorphic $L$-functions}, Ann.\
Inst.\ Fourier (Grenoble) \textbf{57} (2007), no.~5, 1689--1740. arXiv:math/0404394.
\bibitem{ConreyLi} J.~B.~Conrey and X.-J.~Li, \emph{A note on some positivity conditions related to
zeta and $L$-functions}, Int.\ Math.\ Res.\ Not.\ \textbf{2000}, no.~18, 929--940.
\bibitem{Coffey2005} M.~W.~Coffey, \emph{Toward verification of the Riemann hypothesis: application
of the Li criterion}, Math.\ Phys.\ Anal.\ Geom.\ \textbf{8} (2005), no.~3, 211--255.
arXiv:math/0402168.
\bibitem{Coffey2008} M.~W.~Coffey, \emph{New results concerning power series expansions of the
Riemann $\xi$ function and the Li/Keiper constants}, Proc.\ R.\ Soc.\ A \textbf{464} (2008),
no.~2091, 711--731.
\bibitem{Voros2006} A.~Voros, \emph{Sharpenings of Li's criterion for the Riemann hypothesis},
Math.\ Phys.\ Anal.\ Geom.\ \textbf{9} (2006), no.~1, 53--63.
\bibitem{VorosBook} A.~Voros, \emph{Zeta Functions over Zeros of Zeta Functions}, Lecture Notes of
the Unione Matematica Italiana \textbf{8}, Springer, 2010.
\bibitem{OmarMazhouda} S.~Omar and K.~Mazhouda, \emph{Le crit\`ere de Li et l'hypoth\`ese de
Riemann pour la classe de Selberg}, J.\ Number Theory \textbf{125} (2007), no.~1, 50--58.
\bibitem{Maslanka} K.~Ma\'slanka, \emph{Li's criterion for the Riemann hypothesis --- numerical
approach}, Opuscula Math.\ \textbf{24} (2004), no.~1, 103--114.

%% --- Weil positivity and the explicit formula ---
\bibitem{Weil1952} A.~Weil, \emph{Sur les ``formules explicites'' de la th\'eorie des nombres
premiers}, Comm.\ S\'em.\ Math.\ Univ.\ Lund (1952), 252--265.
\bibitem{Weil1948} A.~Weil, \emph{Sur les courbes alg\'ebriques et les vari\'et\'es qui s'en
d\'eduisent}, Actualit\'es Sci.\ Ind.\ \textbf{1041}, Hermann, Paris, 1948.
\bibitem{Deligne} P.~Deligne, \emph{La conjecture de Weil.~I}, Publ.\ Math.\ IH\'ES \textbf{43}
(1974), 273--307.
\bibitem{connes1999} A.~Connes, \emph{Trace formula in noncommutative geometry and the zeros of the
Riemann zeta function}, Selecta Math.\ (N.S.) \textbf{5} (1999), no.~1, 29--106.
arXiv:math/9811068.
\bibitem{ConnesConsani2021} A.~Connes and C.~Consani, \emph{Weil positivity and trace formula: the
archimedean place}, Selecta Math.\ (N.S.) \textbf{27} (2021), no.~4, art.~77. arXiv:2006.13771.
\bibitem{ConnesConsaniRR} A.~Connes and C.~Consani, \emph{Riemann--Roch for $\Spec\Z$}, C.~R.\
Math.\ Acad.\ Sci.\ Paris \textbf{362} (2024), 1727--1732. arXiv:2306.00456.
\bibitem{Burnol} J.-F.~Burnol, \emph{Sur les formules explicites.~I: analyse invariante}, C.~R.\
Acad.\ Sci.\ Paris S\'er.\ I Math.\ \textbf{331} (2000), no.~6, 423--428.
\bibitem{Suzuki2025} M.~Suzuki, \emph{On the Hilbert space derived from the Weil distribution},
Canad.\ J.\ Math.\ (2025), published online. arXiv:2301.00421.
\bibitem{Yoshida} H.~Yoshida, \emph{On Hermitian forms attached to zeta functions}, in: Zeta
Functions in Geometry, Adv.\ Stud.\ Pure Math.\ \textbf{21}, Kinokuniya, 1992, 281--325.
\bibitem{Deninger} C.~Deninger, \emph{Some analogies between number theory and dynamical systems on
foliated spaces}, in: Proc.\ Int.\ Congress of Mathematicians, Vol.~I (Berlin, 1998), Doc.\ Math.\
Extra Vol.~I (1998), 163--186.

%% --- Total positivity, PF sequences, log-concavity ---
\bibitem{Karlin} S.~Karlin, \emph{Total Positivity}, Vol.~I, Stanford Univ.\ Press, 1968.
\bibitem{Schoenberg} I.~J.~Schoenberg, \emph{On smoothing operations and their generating
functions}, Bull.\ Amer.\ Math.\ Soc.\ \textbf{59} (1953), 199--230.
\bibitem{ASW} M.~Aissen, I.~J.~Schoenberg, and A.~M.~Whitney, \emph{On the generating functions of
totally positive sequences.~I}, J.\ Analyse Math.\ \textbf{2} (1952), 93--103.
\bibitem{Edrei} A.~Edrei, \emph{On the generating functions of totally positive sequences.~II},
J.\ Analyse Math.\ \textbf{2} (1952), 104--109.
\bibitem{FallatJohnson} S.~M.~Fallat and C.~R.~Johnson, \emph{Totally Nonnegative Matrices},
Princeton Univ.\ Press, 2011.
\bibitem{BrandenSurvey} P.~Br\"and\'en, \emph{Unimodality, log-concavity, real-rootedness and
beyond}, in: Handbook of Enumerative Combinatorics, CRC Press, 2015, 437--483.
\bibitem{BorceaBranden} J.~Borcea and P.~Br\"and\'en, \emph{The Lee--Yang and P\'olya--Schur
programs.~I}, Comm.\ Pure Appl.\ Math.\ \textbf{62} (2009), no.~12, 1595--1631.
\bibitem{Stanley} R.~P.~Stanley, \emph{Log-concave and unimodal sequences in algebra,
combinatorics, and geometry}, Ann.\ New York Acad.\ Sci.\ \textbf{576} (1989), 500--535.

%% --- Sums of squares and real algebraic geometry ---
\bibitem{Hilbert1888} D.~Hilbert, \emph{\"Uber die Darstellung definiter Formen als Summe von
Formenquadraten}, Math.\ Ann.\ \textbf{32} (1888), 342--350.
\bibitem{Motzkin} T.~S.~Motzkin, \emph{The arithmetic--geometric inequality}, in: Inequalities
(Proc.\ Sympos.\ Wright--Patterson AFB, 1965), Academic Press, 1967, 205--224.
\bibitem{ChoiLam} M.~D.~Choi and T.~Y.~Lam, \emph{Extremal positive semidefinite forms}, Math.\
Ann.\ \textbf{231} (1977), no.~1, 1--18.
\bibitem{Reznick} B.~Reznick, \emph{Some concrete aspects of Hilbert's 17th problem}, in: Real
Algebraic Geometry and Ordered Structures, Contemp.\ Math.\ \textbf{253}, Amer.\ Math.\ Soc.,
2000, 251--272.
\bibitem{Scheiderer} C.~Scheiderer, \emph{Sums of squares of regular functions on real algebraic
varieties}, Trans.\ Amer.\ Math.\ Soc.\ \textbf{352} (2000), no.~3, 1039--1069.
\bibitem{Marshall} M.~Marshall, \emph{Positive Polynomials and Sums of Squares}, Math.\ Surveys
Monogr.\ \textbf{146}, Amer.\ Math.\ Soc., 2008.
\bibitem{PrestelDelzell} A.~Prestel and C.~N.~Delzell, \emph{Positive Polynomials: From Hilbert's
17th Problem to Real Algebra}, Springer, 2001.
\bibitem{Schmudgen} K.~Schm\"udgen, \emph{The $K$-moment problem for compact semi-algebraic sets},
Math.\ Ann.\ \textbf{289} (1991), no.~2, 203--206.
\bibitem{Putinar} M.~Putinar, \emph{Positive polynomials on compact semi-algebraic sets}, Indiana
Univ.\ Math.\ J.\ \textbf{42} (1993), no.~3, 969--984.
\bibitem{Lasserre} J.~B.~Lasserre, \emph{Global optimization with polynomials and the problem of
moments}, SIAM J.\ Optim.\ \textbf{11} (2001), no.~3, 796--817.
\bibitem{Parrilo} P.~A.~Parrilo, \emph{Semidefinite programming relaxations for semialgebraic
problems}, Math.\ Program.\ \textbf{96} (2003), no.~2, 293--320.
\bibitem{Blekherman} G.~Blekherman, P.~A.~Parrilo, and R.~R.~Thomas (eds.), \emph{Semidefinite
Optimization and Convex Algebraic Geometry}, MOS-SIAM Ser.\ Optim.\ \textbf{13}, SIAM, 2013.

%% --- Moment problems, Laplace transforms, complete monotonicity ---
\bibitem{Widder} D.~V.~Widder, \emph{The Laplace Transform}, Princeton Univ.\ Press, 1941.
\bibitem{Bernstein} S.~Bernstein, \emph{Sur les fonctions absolument monotones}, Acta Math.\
\textbf{52} (1929), 1--66.
\bibitem{Hausdorff} F.~Hausdorff, \emph{Summationsmethoden und Momentfolgen.~I}, Math.\ Z.\
\textbf{9} (1921), 74--109.
\bibitem{SSV} R.~L.~Schilling, R.~Song, and Z.~Vondra\v{c}ek, \emph{Bernstein Functions: Theory and
Applications}, 2nd ed., de Gruyter Stud.\ Math.\ \textbf{37}, 2012.
\bibitem{Akhiezer} N.~I.~Akhiezer, \emph{The Classical Moment Problem and Some Related Questions in
Analysis}, Oliver \& Boyd, 1965.
\bibitem{ShohatTamarkin} J.~A.~Shohat and J.~D.~Tamarkin, \emph{The Problem of Moments}, Math.\
Surveys \textbf{1}, Amer.\ Math.\ Soc., 1943.
\bibitem{Aronszajn} N.~Aronszajn, \emph{Theory of reproducing kernels}, Trans.\ Amer.\ Math.\ Soc.\
\textbf{68} (1950), 337--404.
\bibitem{HornJohnson} R.~A.~Horn and C.~R.~Johnson, \emph{Matrix Analysis}, 2nd ed., Cambridge
Univ.\ Press, 2013.

%% --- Infinite divisibility, Levy processes, Stein's method ---
\bibitem{Sato} K.~Sato, \emph{L\'evy Processes and Infinitely Divisible Distributions}, Cambridge
Stud.\ Adv.\ Math.\ \textbf{68}, Cambridge Univ.\ Press, 1999.
\bibitem{SteutelVanHarn} F.~W.~Steutel and K.~van Harn, \emph{Infinite Divisibility of Probability
Distributions on the Real Line}, Marcel Dekker, 2004.
\bibitem{Bondesson} L.~Bondesson, \emph{Generalized Gamma Convolutions and Related Classes of
Distributions and Densities}, Lecture Notes in Statist.\ \textbf{76}, Springer, 1992.
\bibitem{Feller} W.~Feller, \emph{An Introduction to Probability Theory and Its Applications},
Vol.~II, 2nd ed., Wiley, 1971.
\bibitem{LastPenrose} G.~Last and M.~Penrose, \emph{Lectures on the Poisson Process}, IMS Textb.\
\textbf{7}, Cambridge Univ.\ Press, 2018.
\bibitem{Mecke} J.~Mecke, \emph{Station\"are zuf\"allige Ma\ss e auf lokalkompakten Abelschen
Gruppen}, Z.\ Wahrsch.\ Verw.\ Gebiete \textbf{9} (1967), 36--58.
\bibitem{ChenGoldsteinShao} L.~H.~Y.~Chen, L.~Goldstein, and Q.-M.~Shao, \emph{Normal Approximation
by Stein's Method}, Springer, 2011.
\bibitem{BarbourChen} A.~D.~Barbour and L.~H.~Y.~Chen (eds.), \emph{An Introduction to Stein's
Method}, Lect.\ Notes Ser.\ Inst.\ Math.\ Sci.\ \textbf{4}, Singapore Univ.\ Press, 2005.
\bibitem{NourdinPeccati} I.~Nourdin and G.~Peccati, \emph{Normal Approximations with Malliavin
Calculus: From Stein's Method to Universality}, Cambridge Univ.\ Press, 2012.
\bibitem{ShakedShanthikumar} M.~Shaked and J.~G.~Shanthikumar, \emph{Stochastic Orders}, Springer,
2007.
\bibitem{MullerStoyan} A.~M\"uller and D.~Stoyan, \emph{Comparison Methods for Stochastic Models
and Risks}, Wiley, 2002.

%% --- Fourier positivity, Caratheodory--Fejer, Tauberian theory ---
\bibitem{Fejer} L.~Fej\'er, \emph{\"Uber trigonometrische Polynome}, J.\ Reine Angew.\ Math.\
\textbf{146} (1916), 53--82.
\bibitem{Caratheodory} C.~Carath\'eodory, \emph{\"Uber den Variabilit\"atsbereich der Koeffizienten
von Potenzreihen, die gegebene Werte nicht annehmen}, Math.\ Ann.\ \textbf{64} (1907), 95--115.
\bibitem{Toeplitz} O.~Toeplitz, \emph{\"Uber die Fouriersche Entwicklung positiver Funktionen},
Rend.\ Circ.\ Mat.\ Palermo \textbf{32} (1911), 191--192.
\bibitem{Herglotz} G.~Herglotz, \emph{\"Uber Potenzreihen mit positivem, reellen Teil im
Einheitskreis}, Ber.\ Verh.\ S\"achs.\ Akad.\ Wiss.\ Leipzig \textbf{63} (1911), 501--511.
\bibitem{Schur} I.~Schur, \emph{\"Uber Potenzreihen, die im Innern des Einheitskreises beschr\"ankt
sind}, J.\ Reine Angew.\ Math.\ \textbf{147} (1917), 205--232.
\bibitem{Zygmund} A.~Zygmund, \emph{Trigonometric Series}, 3rd ed., Cambridge Univ.\ Press, 2002.
\bibitem{Katznelson} Y.~Katznelson, \emph{An Introduction to Harmonic Analysis}, 3rd ed., Cambridge
Univ.\ Press, 2004.
\bibitem{FoiasFrazho} C.~Foias and A.~E.~Frazho, \emph{The Commutant Lifting Approach to
Interpolation Problems}, Birkh\"auser, 1990.
\bibitem{Korevaar} J.~Korevaar, \emph{Tauberian Theory: A Century of Developments}, Grundlehren
math.\ Wiss.\ \textbf{329}, Springer, 2004.
\bibitem{HardyDivergent} G.~H.~Hardy, \emph{Divergent Series}, Oxford Univ.\ Press, 1949.
\bibitem{Landau} E.~Landau, \emph{Darstellung und Begr\"undung einiger neuerer Ergebnisse der
Funktionentheorie}, 2nd ed., Springer, 1929. \emph{(Pringsheim's theorem on power series with
one-signed coefficients.)}
\bibitem{Remmert} R.~Remmert, \emph{Classical Topics in Complex Function Theory}, Grad.\ Texts in
Math.\ \textbf{172}, Springer, 1998.

%% --- Coercivity, spectral gaps, Poincare inequalities ---
\bibitem{BGL} D.~Bakry, I.~Gentil, and M.~Ledoux, \emph{Analysis and Geometry of Markov Diffusion
Operators}, Grundlehren math.\ Wiss.\ \textbf{348}, Springer, 2014.
\bibitem{Gross} L.~Gross, \emph{Logarithmic Sobolev inequalities}, Amer.\ J.\ Math.\ \textbf{97}
(1975), no.~4, 1061--1083.
\bibitem{ReedSimon} M.~Reed and B.~Simon, \emph{Methods of Modern Mathematical Physics}, vols.~I,
IV, Academic Press, 1972--1978.
\bibitem{Simon} B.~Simon, \emph{Trace Ideals and Their Applications}, 2nd ed., Amer.\ Math.\ Soc.,
2005.
\bibitem{Davies} E.~B.~Davies, \emph{Spectral Theory and Differential Operators}, Cambridge Stud.\
Adv.\ Math.\ \textbf{42}, Cambridge Univ.\ Press, 1995.

%% --- Laguerre asymptotics ---
\bibitem{Szego} G.~Szeg\H{o}, \emph{Orthogonal Polynomials}, 4th ed., Amer.\ Math.\ Soc.\ Colloq.\
Publ.\ \textbf{23}, 1975.
\bibitem{Erdelyi} A.~Erd\'elyi, \emph{Asymptotic Expansions}, Dover, 1956.
\bibitem{Olver} F.~W.~J.~Olver, \emph{Asymptotics and Special Functions}, Academic Press, 1974;
reprint A.~K.~Peters, 1997.
\bibitem{DKMVZ} P.~Deift, T.~Kriecherbauer, K.~T.-R.~McLaughlin, S.~Venakides, and X.~Zhou,
\emph{Strong asymptotics of orthogonal polynomials with respect to exponential weights}, Comm.\
Pure Appl.\ Math.\ \textbf{52} (1999), no.~12, 1491--1552.

%% --- Effective prime number theory and zero-free regions ---
\bibitem{Vinogradov} I.~M.~Vinogradov, \emph{A new estimate of the function $\zeta(1+it)$},
Izv.\ Akad.\ Nauk SSSR Ser.\ Mat.\ \textbf{22} (1958), 161--164.
\bibitem{Korobov} N.~M.~Korobov, \emph{Estimates of trigonometric sums and their applications},
Uspekhi Mat.\ Nauk \textbf{13} (1958), no.~4, 185--192.
\bibitem{Ford} K.~Ford, \emph{Vinogradov's integral and bounds for the Riemann zeta function},
Proc.\ London Math.\ Soc.\ (3) \textbf{85} (2002), no.~3, 565--633.
\bibitem{MTY} M.~J.~Mossinghoff, T.~S.~Trudgian, and A.~Yang, \emph{Explicit zero-free regions for
the Riemann zeta-function}, Res.\ Number Theory \textbf{10} (2024), Paper No.~11.
\bibitem{JohnstonYang} D.~R.~Johnston and A.~Yang, \emph{Some explicit estimates for the error term
in the prime number theorem}, J.\ Math.\ Anal.\ Appl.\ \textbf{523} (2023), Paper No.~127029.
arXiv:2204.01980.
\bibitem{PlattTrudgian} D.~J.~Platt and T.~S.~Trudgian, \emph{The Riemann hypothesis is true up to
$3\times10^{12}$}, Bull.\ Lond.\ Math.\ Soc.\ \textbf{53} (2021), no.~3, 792--797.
arXiv:2004.09765.
\bibitem{Buthe} J.~B\"uthe, \emph{Estimating $\pi(x)$ and related functions under partial RH
assumptions}, Math.\ Comp.\ \textbf{85} (2016), no.~301, 2483--2498.
\bibitem{GuthMaynard} L.~Guth and J.~Maynard, \emph{New large value estimates for Dirichlet
polynomials}, preprint (2024). arXiv:2405.20552.

%% --- Arithmetic identities: Selberg, Mobius, Jordan totient ---
\bibitem{Selberg1949} A.~Selberg, \emph{An elementary proof of the prime-number theorem}, Ann.\ of
Math.\ (2) \textbf{50} (1949), 305--313.
\bibitem{Erdos1949} P.~Erd\H{o}s, \emph{On a new method in elementary number theory which leads to
an elementary proof of the prime number theorem}, Proc.\ Nat.\ Acad.\ Sci.\ U.S.A.\ \textbf{35}
(1949), 374--384.
\bibitem{Apostol} T.~M.~Apostol, \emph{Introduction to Analytic Number Theory}, Springer, 1976.
\bibitem{MSSS} P.~Moree, S.~Saad Eddin, A.~Sedunova, and Y.~Suzuki, \emph{Jordan totient
quotients}, J.\ Number Theory \textbf{209} (2020), 147--166. arXiv:1810.04742.
\bibitem{Eisenbud} D.~Eisenbud, \emph{Commutative Algebra with a View Toward Algebraic Geometry},
Grad.\ Texts in Math.\ \textbf{150}, Springer, 1995. \emph{(Koszul complexes.)}
\bibitem{Voisin} C.~Voisin, \emph{Hodge Theory and Complex Algebraic Geometry}, Vols.~I--II,
Cambridge Univ.\ Press, 2002--2003.

%% --- Related criteria and falsifiers ---
\bibitem{Beurling} A.~Beurling, \emph{A closure problem related to the Riemann zeta-function},
Proc.\ Nat.\ Acad.\ Sci.\ U.S.A.\ \textbf{41} (1955), no.~5, 312--314.
\bibitem{BaezDuarte} L.~B\'aez-Duarte, \emph{A strengthening of the Nyman--Beurling criterion for
the Riemann hypothesis}, Atti Accad.\ Naz.\ Lincei \textbf{14} (2003), no.~1, 5--11.
arXiv:math/0202141.
\bibitem{RodgersTao} B.~Rodgers and T.~Tao, \emph{The de Bruijn--Newman constant is non-negative},
Forum Math.\ Pi \textbf{8} (2020), e6.
\bibitem{PolymathDBN} D.~H.~J.~Polymath, \emph{Effective approximation of heat flow evolution of the
Riemann $\xi$ function, and a new upper bound for the de Bruijn--Newman constant}, Res.\ Math.\
Sci.\ \textbf{6} (2019), no.~3, Paper No.~31.
\bibitem{GORZ} M.~Griffin, K.~Ono, L.~Rolen, and D.~Zagier, \emph{Jensen polynomials for the
Riemann zeta function and other sequences}, Proc.\ Natl.\ Acad.\ Sci.\ USA \textbf{116} (2019),
no.~22, 11103--11110.
\bibitem{DavenportHeilbronn} H.~Davenport and H.~Heilbronn, \emph{On the zeros of certain Dirichlet
series} I, II, J.\ London Math.\ Soc.\ \textbf{11} (1936), 181--185 and 307--312.
\bibitem{Montgomery} H.~L.~Montgomery, \emph{The pair correlation of zeros of the zeta function},
in: Analytic Number Theory, Proc.\ Sympos.\ Pure Math.\ \textbf{24}, Amer.\ Math.\ Soc., 1973,
181--193.
\bibitem{deBranges} L.~de~Branges, \emph{Hilbert Spaces of Entire Functions}, Prentice--Hall,
Englewood Cliffs NJ, 1968.
\bibitem{Bognar} J.~Bogn\'ar, \emph{Indefinite Inner Product Spaces}, Springer, 1974.

\end{thebibliography}

\end{document}
